\documentclass[11pt,a4paper]{article}
\usepackage[T1]{fontenc}
\usepackage[utf8]{inputenc}
\usepackage{lmodern,microtype,amsmath,amssymb,amsthm}
\usepackage[margin=23mm,headheight=14pt]{geometry}
\usepackage{enumitem,xcolor,tikz,fancyhdr,needspace,booktabs,longtable,array}
\usepackage[colorlinks=true,urlcolor=blue,linkcolor=black]{hyperref}
\usepackage{bookmark}
\setlength{\parindent}{0pt}
\setlength{\parskip}{6pt}
\setlength{\emergencystretch}{2em}
\pagestyle{fancy}
\fancyhf{}
\fancyhead[L]{\small Quantum speedups for discrete computational problems}
\fancyfoot[C]{\small\thepage}
\renewcommand{\headrulewidth}{0.3pt}
\hypersetup{pdftitle={Quantum Speedups for Discrete Computational Problems: 102 Entries}}

% Each cell is two body-font ems wide and one body-font em high.
% RGB primaries are blended with white; the three dimensions are independent.
\definecolor{strengthPolynomial}{RGB}{255,224,224}
\definecolor{strengthPolynomialTime}{RGB}{255,0,0}
\definecolor{evidenceProved}{RGB}{0,0,255}
\definecolor{evidenceStrong}{RGB}{85,85,255}
\definecolor{hardwareSmall}{RGB}{64,255,64}
\definecolor{hardwarePolynomial}{RGB}{128,255,128}
\definecolor{hardwareExponential}{RGB}{191,255,191}
\newcommand{\entrybar}[3]{%
  \begin{tikzpicture}[baseline=-2pt,x=22pt,y=11pt]
    \fill[#1] (0,0) rectangle (1,1);
    \fill[#2] (1,0) rectangle (2,1);
    \fill[#3] (2,0) rectangle (3,1);
    \draw[black,line width=0.35pt] (0,0) rectangle (3,1);
    \draw[black,line width=0.35pt] (1,0)--(1,1) (2,0)--(2,1);
  \end{tikzpicture}%
}
\newcounter{catalogueentry}
\newcommand{\entrytitle}[5]{%
  \Needspace{10\baselineskip}%
  \setcounter{catalogueentry}{#1}\addtocounter{catalogueentry}{-1}%
  \refstepcounter{catalogueentry}\label{entry:#1}%
  \addcontentsline{toc}{subsection}{\protect\numberline{#1}#2}%
  \par\vspace{12pt}{\Large\bfseries\raggedright #1. #2\nobreak\hspace{0.7em}\entrybar{#3}{#4}{#5}\par}\vspace{5pt}
}
\newcommand{\colorcell}[1]{%
  \begin{tikzpicture}[baseline=-2pt,x=22pt,y=11pt]
    \fill[#1] (0,0) rectangle (1,1);
    \draw[black,line width=0.35pt] (0,0) rectangle (1,1);
  \end{tikzpicture}%
}
\newcommand{\dimensionlabel}[2]{\par\noindent \textbf{#1.}\nobreak\hspace{0.4em}\colorcell{#2}\hspace{0.5em}}
\newtheoremstyle{problemstyle}{5pt}{9pt}{\normalfont}{}{\bfseries}{.}{\newline}{}
\theoremstyle{problemstyle}
\newtheorem{problem}{Problem}
\newcommand{\subhead}[1]{\Needspace{4\baselineskip}\par\vspace{5pt}{\bfseries #1\par}\vspace{1pt}}
\newcommand{\speedupheading}{\Needspace{5\baselineskip}\par\vspace{7pt}{\fontsize{13}{16}\selectfont\bfseries Quantum speedup\par}\vspace{3pt}}
\newcommand{\complexity}[2]{\par\noindent\textbf{\boldmath #1 complexity: #2}\ }
\newenvironment{algorithmindex}[1]{%
  \begin{enumerate}[label={[#1.\arabic*]},leftmargin=3.4em,labelsep=0.6em,itemsep=4pt,topsep=3pt,parsep=0pt]
  \small
}{\end{enumerate}}


\definecolor{strengthIntermediate}{RGB}{255,208,208}
\definecolor{strengthExpToSubexp}{RGB}{255,112,112}
\definecolor{evidencePartial}{RGB}{170,170,255}
\definecolor{hardwareDemonstrated}{RGB}{0,255,0}
\setcounter{tocdepth}{2}
\begin{document}
\begin{titlepage}
\vspace*{24mm}
{\small ALGORITHMS\quad /\quad COMPLEXITY\quad /\quad CRYPTOGRAPHY\par}
\vspace{14mm}
{\Huge\bfseries Quantum Speedups for\par Discrete Computational Problems\par}
\vspace{12mm}
{\Large 102 problem and complexity entries\par}
\vspace{10mm}
{\large A research catalogue for a knowledge graph\par}
\vfill
Each entry connects a formal problem with quantum algorithms, classical algorithms, complexity comparisons, evidence, and implementation resources. Complexity-theoretic entries describe the bounds and implications linking these results.\par
\vspace{20mm}
{\large 9 November 2026\par}
\end{titlepage}

\section*{Introduction}
\addcontentsline{toc}{section}{Introduction}
The first three chapters study decision, search, optimization, counting and sampling on combinatorial structures, integers, finite fields and groups, and lattices. Numerical weights, approximate outputs and continuous or simulation subroutines are included for discrete target problems. Continuous optimization, numerical linear algebra, numerical differential equations and Hamiltonian simulation are outside their scope as independent problems. The fourth chapter, Complexity bounds and implications, also considers continuous and quantum-state problems to study quantum complexity classes.

The organizing unit is a mathematical problem with its input representation, promise or approximation conditions, and success criterion. Quantum and classical bounds are compared in the same regime. Algorithms for the same problem appear together; their formal definitions and complexity statements are located in the literature under Algorithm.

Each entry carries three title cells: strength (red), evidence (blue), and hardware (green). Greater saturation indicates a higher category. The corresponding colors also appear beside Strength, Evidence and Hardware, whose text specifies the comparison and implementation resources. White denotes the lowest category or a dimension that does not apply to a complexity-theoretic constraint.

\subhead{Strength}
Red indicates the strongest supported speedup in an entry, across its parameter regimes and input models. The four categories use nearby pale shades for the first two and progressively more saturated red for the last two. Complexity bounds retain their general parameter dependence; the accompanying text states the assumptions, evidence and hardware resources.

\begin{tabular}{@{}ll@{}}
\colorcell{strengthPolynomial} & Polynomial speedup\\[3pt]
\colorcell{strengthIntermediate} & Improved exponential/subexponential complexity\\[3pt]
\colorcell{strengthExpToSubexp} & Exponential to subexponential\\[3pt]
\colorcell{strengthPolynomialTime} & Exponential or subexponential to polynomial
\end{tabular}

\emph{Polynomial speedup} includes fixed-power relations between the classical and quantum running times. In particular, $2^{an}\to2^{bn}$ for constants $0<b<a$ satisfies $T_Q=T_C^{b/a}$ and belongs in this category. Quadratic and higher-degree improvements are identified in the accompanying text, together with their origin in amplitude amplification, quantum walks or other methods.

\emph{Improved exponential/subexponential complexity} covers superpolynomial advantages for which both runtimes remain superpolynomial within the same exponential or subexponential regime. The quantum time is a vanishing power of the classical time, $\log T_Q=o(\log T_C)$. The other two categories describe the changes from exponential to subexponential time and from exponential or subexponential to polynomial time.

\emph{Power saving} remains a textual descriptor: it means a decrease in the power of the input size inside the exponent. Thus $2^{n^{2/3}}\to2^{n^{1/3}}$ has power saving, whereas $2^{\sqrt n\log n}\to2^{\sqrt n}$ improves a logarithmic factor. Both comparisons belong to the second color category. The change $2^{\Theta(n)}\to2^{O(n/\sqrt{\log n})}$ belongs to the third category and has no power saving. A change of parameter is stated explicitly; the oracle address length and the full oracle-table size are distinct quantities.

\Needspace{8\baselineskip}
\subhead{Evidence}
\begin{tabular}{@{}ll@{}}
\colorcell{evidenceProved} & Fully rigorous: matching bounds in the same computational model\\[3pt]
\colorcell{evidenceStrong} & Strong conditional evidence from an explicit hardness assumption\\[3pt]
\colorcell{evidencePartial} & Partial evidence from the best known classical algorithms\\[3pt]
\colorcell{white} & No established evidence for a speedup
\end{tabular}

A proved quantum upper bound can coexist with partial evidence for a speedup. Strong conditional evidence identifies the assumption or implication at issue; for factoring and discrete logarithms, this includes the corresponding cryptographic hardness assumptions. Complexity-theoretic entries use blue for the status of the stated theorem or conditional implication.

\Needspace{12\baselineskip}
\subhead{Hardware}
\begin{tabular}{@{}ll@{}}
\colorcell{hardwareDemonstrated} & Demonstrated advantage at a substantial problem scale\\[3pt]
\colorcell{hardwareSmall} & Small-instance physical implementation or benchmark\\[3pt]
\colorcell{hardwarePolynomial} & Prospective fault-tolerant implementation with polynomial resources\\[3pt]
\colorcell{hardwareExponential} & Prospective implementation with superpolynomial depth or storage\\[3pt]
\colorcell{white} & Implementation resources unspecified or not applicable
\end{tabular}

The hardware text identifies the implemented task, logical-qubit requirements, coherent memory, depth, and precision. Logical qubits encode protected quantum information; physical qubits implement the chosen error-correcting code. Analog atoms or spins and noisy circuit qubits are identified separately. A query bound determines the number of oracle calls; a circuit realization additionally specifies the oracle and its workspace.

\subhead{Complexity conventions}
The last sentence of each Problem states the resource being measured. Oracle or query complexity counts calls to the specified input oracle. Gate complexity counts elementary quantum gates. Classical running time may count bit operations, word-RAM operations, arithmetic operations, or real-RAM operations, with the model stated in the entry. Sampling complexity counts observations; continuous-time algorithms specify evolution time and the Hamiltonian energy normalization. Mixed comparisons explicitly identify each resource.

Bounded error means success probability at least $2/3$ unless another value is stated. $O^*(\cdot)$ suppresses factors polynomial in the input length, and $\widetilde O(\cdot)$ suppresses polylogarithmic factors. Variables are defined locally. In cryptography, bit length and group or field cardinality are distinct parameters. The standard notation
\[
L_X[\alpha,c]=\exp\!\big((c+o(1))(\log X)^\alpha(\log\log X)^{1-\alpha}\big)
\]
describes subexponential number-theoretic algorithms. Quantum random-access memory (QRAM) provides coherent address access; its storage size is separate from the address-register size.

\subhead{Reading paths}
The catalogue has four chapters. Entries~\ref{entry:1}--\ref{entry:30} and~\ref{entry:102} cover exact optimization and search, including recursive Fourier sampling. Entries~\ref{entry:31}--\ref{entry:44}, \ref{entry:46}, \ref{entry:49}--\ref{entry:56}, \ref{entry:59}--\ref{entry:60} and~\ref{entry:103}--\ref{entry:107} cover graph algorithms and discrete optimization, including variational approaches to partitioning, scheduling, portfolios and community detection. Entries~\ref{entry:61}--\ref{entry:90} and~\ref{entry:101} cover cryptanalysis and post-quantum problems. Entries~\ref{entry:91}--\ref{entry:100} connect the algorithms to query lower bounds, fine-grained hypotheses, Hamiltonian complexity and variational optimization. The final table lists relations for the knowledge graph.
\clearpage
\tableofcontents
\clearpage

\clearpage
\section{Exact combinatorial optimization and search}

\entrytitle{1}{Unstructured search}{strengthPolynomial}{evidenceProved}{hardwareSmall}
\setcounter{problem}{0}
\begin{problem}[Unstructured search]
Given coherent access to a Boolean oracle
\[
 f:\{0,\ldots,N-1\}\longrightarrow\{0,1\},
\]
output an $x$ with $f(x)=1$, or report that no such $x$ exists, with success probability at least $2/3$.
Let $M=|f^{-1}(1)|$. In a combinatorial search problem, $x$ represents a candidate solution and $f$ its feasibility predicate.
Complexity is measured in oracle queries for both models, with an additional quantum gate bound expressed in terms of the oracle implementation cost.
\end{problem}

\speedupheading
\complexity{Quantum}{$\Theta(\sqrt N)$ oracle queries, Grover search and amplitude amplification.}
For $M\ge1$, the expected query bound is $O(\sqrt{N/M})$, including the case of unknown $M$.
If a reversible evaluation of $f$ costs $C_f$ gates, the direct implementation uses
$O((C_f+\log N)\sqrt{N/M})$ gates.
For a general reversible trial with success probability $p$, amplitude amplification reduces the required number of trials to $O(p^{-1/2})$.

\complexity{Classical}{$\Theta(N)$ oracle queries, exhaustive or randomized search.}
Under a promise of $M\ge1$ marked inputs, randomized sampling requires $\Theta(N/M)$ queries in the worst case over their locations. A general trial with success probability $p$ needs $\Theta(p^{-1})$ independent repetitions. The upper and lower bounds concern the same Boolean-oracle access model.

\dimensionlabel{Strength}{strengthPolynomial} Tight quadratic polynomial speedup in the search-space size, arising from Grover search and amplitude amplification. With $N=2^n$, the comparison is $2^n$ versus $2^{n/2}$, an improvement in the exponential constant without power saving.

\dimensionlabel{Evidence}{evidenceProved} Fully rigorous: the classical and quantum query lower bounds match the stated algorithms; see \href{https://arxiv.org/pdf/quant-ph/9701001}{Bennett--Bernstein--Brassard--Vazirani (1997)}.

\dimensionlabel{Hardware}{hardwareSmall} Small-instance search circuits have been implemented. A useful large search requires $\lceil\log_2N\rceil$ logical candidate qubits, reversible-oracle workspace, and $\Theta(\sqrt N)$ coherent iterations in the worst case.

\subhead{Algorithm}
\begin{algorithmindex}{G}
\item \href{https://arxiv.org/pdf/quant-ph/9605043}{Grover, \emph{A Fast Quantum Mechanical Algorithm for Database Search} (STOC 1996)}. Section 3, ``Algorithm'', preprint p.~3, steps (i)--(iii).
\item \href{https://arxiv.org/pdf/quant-ph/0005055}{Brassard, H{\o}yer, Mosca and Tapp, \emph{Quantum Amplitude Amplification and Estimation} (2002)}. Section 2, Theorem 2 for known success probability; Theorem 3 and \textsc{QSearch} for unknown success probability, preprint pp.~7--9.
\end{algorithmindex}

\entrytitle{2}{Black-box discrete minimization}{strengthPolynomial}{evidenceProved}{hardwareExponential}
\setcounter{problem}{1}
\begin{problem}[Black-box minimum finding]
Given an oracle evaluating $F:[N]\to\{0,\ldots,2^b-1\}$, output $x_*\in\operatorname*{arg\,min}_{x\in[N]}F(x)$. This is the optimization primitive used inside quantum dynamic programming and threshold-based exact optimization.
Complexity is measured in value-oracle queries for both models, with a quantum gate bound parameterized by the evaluation circuit.
\end{problem}
\speedupheading
\complexity{Quantum}{$\Theta(\sqrt N)$ queries, D\"urr--H{\o}yer minimum finding.}
The algorithm returns a minimizing index with bounded error. Failure probability at most $\delta$ costs $O(\sqrt N\log(1/\delta))$ evaluations. If reversible evaluation costs $C_F$ gates, constant-error gate complexity is $O((C_F+b+\log N)\sqrt N)$. Infeasible candidates can be represented by values above the feasible objective range.
\complexity{Classical}{$\Theta(N)$ queries, linear scanning.}
A scan retains the smallest value and its index, using constant additional words. The matching lower bound concerns arbitrary values accessed through the same oracle. This primitive is used inside the structured optimization algorithms below, whose classical comparisons exploit their explicit input representations.
\dimensionlabel{Strength}{strengthPolynomial} Tight quadratic polynomial speedup in $N$, obtained through repeated amplitude amplification.
\dimensionlabel{Evidence}{evidenceProved} Fully proved in the value-oracle model; reduction from unstructured search supplies the quantum lower bound, and an unseen table entry supplies the classical lower bound.
\dimensionlabel{Hardware}{hardwareExponential} The search and value registers need $\lceil\log_2N\rceil+b$ logical qubits plus reversible evaluation workspace. A compact objective can be evaluated arithmetically; a table supplied through coherent memory additionally requires storage for its entries. Useful large searches require fault tolerance and $\Theta(\sqrt N)$ coherent evaluations.
\subhead{Algorithm}
\begin{algorithmindex}{2}
\item \href{https://arxiv.org/pdf/quant-ph/9607014}{D\"urr and H{\o}yer, \emph{A Quantum Algorithm for Finding the Minimum} (1996/1999)}. Theorem 1 and the numbered minimum-finding procedure, preprint pp.~1--2; the linear scan is the classical comparison in the introduction.
\end{algorithmindex}

\entrytitle{3}{Constraint satisfaction by backtracking}{strengthPolynomial}{evidencePartial}{hardwareExponential}
\setcounter{problem}{2}
\begin{problem}[Constraint satisfaction]
For variables $x_1,\ldots,x_n\in[d]$ and explicitly given constraints $C_1,\ldots,C_m$, find $x\in[d]^n$ with $\bigwedge_jC_j(x)=1$, or decide unsatisfiability. Assume a classical branching rule and a predicate that rejects inconsistent partial assignments. Let $T$ be the size of their backtracking tree and $D\leq n$ its depth.
The tree comparison counts branching and predicate-oracle calls, while the explicit bounded-arity comparison uses quantum gates and classical bit operations.
\end{problem}
\speedupheading
\complexity{Quantum}{$O(\sqrt{TD})$ detection queries and $O(\sqrt T D^{3/2}\log D)$ witness queries, Montanaro's quantum backtracking.}
These are constant-error bounds for the supplied branching rule and pruning predicate. Reversible implementation costs multiply the query count. The quantum walk explores the same partial-assignment tree as the specified classical backtracking procedure; its depth dependence is explicit in the bounds.
\complexity{Classical}{$O(T)$ tree queries, depth-first backtracking; $2^{o(n)}\operatorname{poly}(n,m)$ bit operations, Boolean tensor contraction for fixed domain and arity.}
For fixed domain $d$ and arity $r$, the time bound is $\operatorname{poly}(n,m)2^{O_{d,r}(n\log\log n/\sqrt{\log n})}$. Encoding each variable in $O(\log d)$ bits and forbidding disallowed local assignments produces a fixed-width CNF with $O(n)$ variables. Exact counting and self-reduction decide feasibility and recover a witness in polynomial space. When $m=O(n)$, the exponent improves to $O_{d,r}(n/\sqrt{\log n})$. The tree-query comparison remains useful for an instance whose pruning rules yield a small $T$; the bounded-arity expression is the unrestricted explicit-input comparator in that regime.
\dimensionlabel{Strength}{strengthPolynomial} Nearly quadratic in the size of a specified search tree, through a quantum walk. Fixed-domain, fixed-arity instances also admit the displayed subexponential classical bound, whose exponent has no power saving.
\dimensionlabel{Evidence}{evidencePartial} Rigorous acceleration of the supplied backtracking procedure; evidence against unrestricted classical solvers depends on the constraint language and instances.
\dimensionlabel{Hardware}{hardwareExponential} A path requires $O(D\log d)$ qubits, together with reversible constraint-evaluation and phase-estimation workspace. The register count is polynomial, while a tree of exponential size entails exponentially many coherent walk operations. Useful large instances require fault tolerance.
\subhead{Algorithm}
\begin{algorithmindex}{3}
\item \href{https://arxiv.org/pdf/1509.02374}{Montanaro, \emph{Quantum Walk Speedup of Backtracking Algorithms} (Theory of Computing 2018)}. Theorems 1--2; Algorithm 1 defines classical backtracking, Algorithm 2 defines quantum detection, and the section on finding a marked vertex gives witness recovery.
\item \href{https://github.com/openai/math/blob/main/preprints/The-Exponential-Time-Hypothesis-is-False-October-25-2026/paper.pdf}{OpenAI, \emph{The Exponential Time Hypothesis Is False} (2026)}. Theorem 1.1 and Appendices D.1--D.4; applied to the fixed-width encoding described above.
\end{algorithmindex}

\entrytitle{4}{Discrete optimization by branch and bound}{strengthPolynomial}{evidencePartial}{hardwareExponential}
\setcounter{problem}{3}
\begin{problem}[Branch-and-bound optimization]
Minimize $F(x)$ over a finite feasible set represented by a rooted branching tree. Each node represents a subset of candidates and carries a computable lower bound on their objective values. Return an optimal feasible leaf and its value. Integer programming and scheduling commonly use this representation.
Complexity is measured in branching, bounding and heuristic-oracle queries; quantum gate complexity additionally includes their reversible implementation costs.
\end{problem}
\speedupheading
\complexity{Quantum}{$\widetilde O(\sqrt Q\,D)$ queries, Incremental Quantum Branch and Bound.}
Here $Q$ is the number of queries used by the specified classical branch-and-bound algorithm, $D$ is the tree depth, and the heuristic satisfies the branch-local conditions of the theorem. Logarithmic factors include the objective range, heuristic range and failure probability. Montanaro's construction gives an earlier near-quadratic bound in the truncated optimal-cost tree size. The incremental construction follows the portion of the tree explored by the chosen heuristic and returns an optimal leaf, or a leaf within a prescribed additive tolerance when that stopping criterion is supplied.
\complexity{Classical}{$O(Q)$ queries, the specified branch-and-bound procedure.}
The comparator uses the same branching rule, objective bounds and search heuristic. Its work can be far smaller than the size of the full search tree. For a concrete integer-programming or scheduling formulation, specialized algorithms and branch-and-cut methods supply additional classical comparisons; their costs depend on the input representation and bounding rules.
\dimensionlabel{Strength}{strengthPolynomial} Nearly quadratic polynomial speedup in the work of the selected classical tree-search procedure, based on quantum tree search and subtree estimation.
\dimensionlabel{Evidence}{evidencePartial} The query acceleration is rigorous under the stated heuristic conditions. Its significance relative to all classical algorithms is problem- and instance-dependent.
\dimensionlabel{Hardware}{hardwareExponential} A path, objective registers and reversible bounding workspace use polynomially many logical qubits. Long coherent tree searches require fault tolerance; an exponential $Q$ leaves an exponential coherent running time. The implementation cost of the bound and heuristic is part of the resource estimate.
\subhead{Algorithm}
\begin{algorithmindex}{4}
\item \href{https://arxiv.org/pdf/1906.10375}{Montanaro, \emph{Quantum Speedup of Branch-and-Bound Algorithms} (Physical Review Research 2020)}. Algorithm 2 and Theorem 1.
\item \href{https://arxiv.org/pdf/2210.03210}{Chakrabarti, Minssen, Yalovetzky and Pistoia, \emph{Universal Quantum Speedup for Branch-and-Bound, Branch-and-Cut, and Tree-Search Algorithms} (2022)}. Theorem 1.3 and Section 2.4, Algorithm 2, ``Incremental-Quantum-Branch-and-Bound''; Section 1 defines the classical heuristic conditions.
\end{algorithmindex}

\entrytitle{5}{Fixed-width Boolean satisfiability}{white}{white}{hardwareExponential}
\setcounter{problem}{4}
\begin{problem}[$k$-SAT]
Fix an integer $k\ge3$. Given
\[
 F(x_1,\ldots,x_n)=\bigwedge_{j=1}^{m}C_j,
 \qquad |C_j|\le k,
\]
where each $C_j$ is a disjunction of literals, output $x\in\{0,1\}^n$ satisfying $F(x)=1$, or decide that none exists. The clause list is supplied explicitly. SAT represents Boolean feasibility in constraint programming and exact optimization.
The comparison uses quantum gate complexity and classical bit complexity on explicit formula inputs.
\end{problem}

\speedupheading
\complexity{Quantum}{$O^*((2-2/k)^{n/2})$, quantum Sch\"oning; $O^*(1.1433^n)$ for 3-SAT, amplified PPSZ.}
These are bounded-error trial-amplification bounds, with polynomial factors for formula evaluation suppressed. They assume coherent evaluation and reversal of the classical randomized trial. Running the classical counting algorithm below on a quantum computer also supplies a subexponential upper bound for explicit $k$-SAT.

\complexity{Classical}{$(n+m)^{O_k(1)}2^{O_k(n\log\log n/\sqrt{\log n})}$, exact counting by Boolean tensor contraction.}
The algorithm is deterministic and uses polynomial space. For $m\le Dn$ with fixed $D$, the time is
$(n+m)^{O_{k,D}(1)}2^{O_{k,D}(n/\sqrt{\log n})}$.
Positivity of the count decides satisfiability; conditional counts recover an assignment with polynomial overhead.

\dimensionlabel{Strength}{white} The subexponential classical bound is asymptotically smaller than the displayed amplified-trial bounds. Its exponent is $n^{1-o(1)}$, so the reduction from exponential to subexponential time has no power saving.

\dimensionlabel{Evidence}{white} The stated algorithms have rigorous upper bounds. The comparison supplies no asymptotic quantum-speedup claim for amplified Sch\"oning or PPSZ.

\dimensionlabel{Hardware}{hardwareExponential} The amplified-trial constructions require fault tolerance, $n$ assignment qubits, polynomial reversible workspace and exponentially many coherent trial evaluations.

{\small [Before ``ETH is false'', the known general classical 3-SAT bound was $O^*(1.30704^n)$, and amplified PPSZ claimed a quadratic improvement to $O^*(1.1433^n)$.]\par}

\subhead{Algorithm}
\begin{algorithmindex}{S}
\item \href{https://www.cs.utep.edu/vladik/2005/tr05-26.pdf}{Dantsin, Kreinovich and Wolpert, \emph{On Quantum Versions of Record-Breaking Algorithms for SAT}, SIGACT News 36(4), 103--108 (2005)}. Section 1: ``Ambainis' observation'' and ``The fastest algorithm for $k$-SAT'', author-preprint pp.~2--3.
\item \href{https://www.eecs.harvard.edu/~michaelm/CS223/schoning99.pdf}{Sch\"oning, \emph{A Probabilistic Algorithm for $k$-SAT and Constraint Satisfaction Problems} (FOCS 1999)}. Section 2, ``The Algorithm'', p.~410. \href{https://theoretics.episciences.org/13222/pdf}{Scheder, \emph{PPSZ Is Better Than You Think} (TheoretiCS 2024)}: Section 1.1, Algorithm 1, p.~3.
\item \href{https://github.com/openai/math/blob/main/preprints/The-Exponential-Time-Hypothesis-is-False-October-25-2026/paper.pdf}{OpenAI, \emph{The Exponential Time Hypothesis Is False} (25 October 2026)}. Theorem 1.1, p.~2; Appendices D.1--D.3, pp.~53--54; recursive contraction in D.4.1--D.4.7, pp.~57--64. Witness self-reduction follows Theorem 1.1.
\end{algorithmindex}

\entrytitle{6}{Approximate counting of feasible assignments}{strengthPolynomial}{evidenceProved}{hardwareSmall}
\setcounter{problem}{5}
\begin{problem}[Approximate counting]
For an oracle $f:\{0,1\}^n\to\{0,1\}$, let $a=2^{-n}\sum_x f(x)$. Output $\widehat a$ such that $|\widehat a-a|\leq\epsilon$ with probability at least $2/3$. The count is $2^n\widehat a$; this estimates solution density, feasible volume, or a probability under a reversibly samplable candidate distribution.
The black-box bounds count oracle queries, with gate costs obtained from the reversible oracle; the explicit-CNF comparator counts classical bit operations.
\end{problem}
\speedupheading
\complexity{Quantum}{$O(1/\epsilon)$ queries, quantum counting and amplitude estimation.}
For additive error $\epsilon$ and constant failure probability, coherent sampling and its inverse give the stated bound. Relative error $\epsilon a$ for $a>0$ costs $O(1/(\epsilon\sqrt a))$ calls. Gate complexity multiplies these counts by the state-preparation, oracle and reflection costs. The finite domain permits exact enumeration when that is cheaper than the displayed precision-dependent bound.
\complexity{Classical}{$\Theta(1/\epsilon^2)$ worst-case sampling queries, Monte Carlo estimation.}
Independent samples give the additive-error bound, capped by $2^n$ evaluations of the entire domain. Relative error costs $O(1/(a\epsilon^2))$ samples. When $f$ is an explicit fixed-width CNF, deterministic tensor contraction computes the count exactly in $\operatorname{poly}(n,m)2^{O(n\log\log n/\sqrt{\log n})}$ bit operations and polynomial space; for $m=O(n)$ the exponent is $O(n/\sqrt{\log n})$. This exact method is an additional comparator at high precision or very small feasible density.
\dimensionlabel{Strength}{strengthPolynomial} Quadratic polynomial improvement in inverse additive precision in the sampling-oracle model, from amplitude estimation. Relative-error improvement also depends on the feasible density $a$.
\dimensionlabel{Evidence}{evidenceProved} Tight worst-case sampling-oracle separation in the precision range before finite-domain saturation. Explicit formula structure gives the additional classical comparison stated above.
\dimensionlabel{Hardware}{hardwareSmall} Small amplitude-estimation circuits have been demonstrated. The input register needs $n$ qubits plus oracle workspace; phase-estimation implementations add $O(\log(1/\epsilon))$ ancillas, while iterative variants reduce this register. Useful high precision requires long coherent oracle sequences and fault tolerance.
\subhead{Algorithm}
\begin{algorithmindex}{6}
\item \href{https://arxiv.org/pdf/quant-ph/0005055}{Brassard, H{\o}yer, Mosca and Tapp, \emph{Quantum Amplitude Amplification and Estimation} (2002)}. Section 4, \textsc{Est Amp} and Theorem 12; Theorem 13 defines \textsc{Count}, and Theorem 15 gives relative approximation.
\item \href{https://github.com/openai/math/blob/main/preprints/The-Exponential-Time-Hypothesis-is-False-October-25-2026/paper.pdf}{OpenAI, \emph{The Exponential Time Hypothesis Is False} (2026)}. Theorem 1.1 and Appendix D: exact fixed-width counting.
\end{algorithmindex}

\entrytitle{7}{Travelling salesman}{strengthPolynomial}{evidencePartial}{hardwareExponential}
\setcounter{problem}{6}
\begin{problem}[Travelling salesman problem]
Given a weighted graph $G=(V,E,w)$, $|V|=n$, find a Hamiltonian cycle $C$ minimizing $\sum_{e\in C}w(e)$, or report infeasibility. Integer weights have maximum magnitude $L$ and are binary encoded. The problem models routing and cyclic sequencing.
The general quantum bound uses QRAM time and the bounded-degree bounds use quantum gate time, while classical bounds use RAM time with polynomial weight-bit costs.
\end{problem}
\speedupheading
\complexity{Quantum}{$O^*(1.728^n)$, quantum subset dynamic programming; $O^*(1.110^n)$ and $O^*(1.301^n)$ on degree-three and degree-four graphs, quantum backtracking.}
The general algorithm assumes coherent access to its exponentially large precomputed table. The bounded-degree algorithms use polynomial workspace and have additional factors $\log L\log\log L$ for integer weights. Their graph restrictions are part of the stated comparisons.
\complexity{Classical}{$O(n^2 2^n)$, Held--Karp dynamic programming; $O^*(1.1757^n)$ and $O^*(1.4214^n)$ on degree-three and degree-four graphs, sparse-graph path-decomposition algorithms.}
Held--Karp uses $O(n2^n)$ words and optimizes paths indexed by a visited set and endpoint. The sparse-graph bounds use exponential memory and the weighted connectivity algorithm of Gaspers. A polynomial-space cubic-graph comparator runs in $O^*(1.232^n)$ time. Thus each quantum bound is compared with the applicable degree and memory regime, with the bit cost of objective arithmetic retained in polynomial factors.
\dimensionlabel{Strength}{strengthPolynomial} Polynomial speedups through smaller exponential constants. General subset dynamic programming uses minimum finding and amplitude amplification; the bounded-degree constructions use quantum backtracking.
\dimensionlabel{Evidence}{evidencePartial} Rigorous upper bounds and partial evidence relative to the strongest displayed classical algorithms in their respective graph classes.
\dimensionlabel{Hardware}{hardwareExponential} The general algorithm needs polynomial active registers and exponential QRAM. The bounded-degree versions need polynomial logical workspace but exponentially long coherent tree searches. Both useful regimes require fault tolerance; degree restrictions reduce the exponential time base.
\subhead{Algorithm}
\begin{algorithmindex}{7}
\item \href{https://arxiv.org/pdf/1807.05209}{Ambainis et al., \emph{Quantum Speedups for Exponential-Time Dynamic Programming Algorithms} (SODA 2019)}. Theorem 8 and Algorithm 2; the TSP section states the classical Held--Karp recurrence.
\item \href{https://arxiv.org/pdf/1612.06203}{Moylett, Linden and Montanaro, \emph{Quantum Speedup of the Traveling-Salesman Problem for Bounded-Degree Graphs} (2017)}. Theorem 1; Theorems 4--5 give the degree-three and degree-four algorithms and their classical search subroutines.
\item \href{https://serggasp.github.io/pathwidth-site/papers/sparse-pathwidth/14-apps-connectivity.html}{Gaspers, \emph{Path Decompositions of Sparse Graphs via Gaussian Sweeps and Faster Exponential Algorithms} (2026)}. Theorem 14.1(d): weighted TSP applications.
\end{algorithmindex}

\entrytitle{8}{Minimum set cover}{strengthPolynomial}{evidencePartial}{hardwareExponential}
\setcounter{problem}{7}
\begin{problem}[Minimum set cover]
Given a universe $U$, $|U|=n$, and an explicit family $\mathcal F\subseteq2^U$, $|\mathcal F|=m$, minimize $|\mathcal C|$ subject to $\mathcal C\subseteq\mathcal F$ and $\bigcup_{S\in\mathcal C}S=U$. Typical interpretations select facilities, tests, or resources that jointly cover all requirements.
The quantum bound counts QRAM time, the subset algorithms count classical RAM operations, and the CNF-based bounds count bit operations on explicit inputs.
\end{problem}
\speedupheading
\complexity{Quantum}{$O^*(1.7274^n)$, quantum set-cover dynamic programming.}
Polynomial factors include the explicit family size $m$. The method combines subset preprocessing with recursive quantum minimum finding and assumes coherent table access. Its exponential parameter is universe size $n$, allowing the family to be larger than a linear number of sets.
\complexity{Classical}{$O^*(\min\{2^n,2^m\})$, subset dynamic programming or family enumeration.}
For $m=\Theta(n)$, a stronger bound is $O_\epsilon(2^{\epsilon n}\operatorname{poly}(n,m))$ for every fixed $\epsilon>0$, from long-clause shortening and fixed-width counting. Selection variables, element-cover clauses and a linear-size cardinality circuit give $O(n)$ variables and clauses in this regime. If element frequency is fixed and $m=O(n)$, the encoding has fixed width and linear size, giving $\operatorname{poly}(n,m)2^{O(n/\sqrt{\log n})}$ time. The unrestricted-family comparison and these linear-family regimes use different relations between $n$ and $m$.
\dimensionlabel{Strength}{strengthPolynomial} Polynomial speedup through a smaller exponential constant in the unrestricted-family, universe-size comparison, from amplitude amplification. For a linear-size explicit family, the subexponential classical bound is asymptotically smaller.
\dimensionlabel{Evidence}{evidencePartial} The quantum upper bound is rigorous in QRAM. The broader Set Cover Conjecture concerns specified universe-size and family-size regimes; the linear-family subexponential bound must be included when selecting that comparison.
\dimensionlabel{Hardware}{hardwareExponential} Polynomially many active subset and objective qubits accompany exponentially large coherent preprocessing memory. The storage and long nested searches make useful instances a fault-tolerant target.
{\small [Before ``ETH is false'', the explicit-family comparison used $O^*(\min\{2^n,2^m\})$ classically and $O^*(1.7274^n)$ quantumly, including the linear-family regime now covered by the subexponential bound.]\par}
\subhead{Algorithm}
\begin{algorithmindex}{8}
\item \href{https://arxiv.org/pdf/1807.05209}{Ambainis et al., \emph{Quantum Speedups for Exponential-Time Dynamic Programming Algorithms} (SODA 2019)}. Theorem 10 and Algorithm 3: Set Cover.
\item \href{https://arxiv.org/pdf/2311.08042}{Gaspers and Li, \emph{Quantum Algorithms for Graph Coloring and Other Partitioning, Covering, and Packing Problems} (ICALP 2024)}. Theorems 4--6 and Corollary 1: general covering framework and refined bounds.
\item \href{https://github.com/anthropics/formal-math/tree/main/eth}{Claude (Anthropic), \emph{A $2^{o(n)}$ Time Algorithm for 3-SAT} (2026)}. Long-clause shortening and Corollary 7; combined with \href{https://github.com/openai/math/blob/main/preprints/The-Exponential-Time-Hypothesis-is-False-October-25-2026/paper.pdf}{OpenAI, \emph{The Exponential Time Hypothesis Is False}, Theorem 1.1 and Appendix D}, for the explicit-family encodings above.
\end{algorithmindex}

\entrytitle{9}{Set multicover with repeated sets}{strengthPolynomial}{evidencePartial}{hardwareExponential}
\setcounter{problem}{8}
\begin{problem}[Set multicover with repeated sets]
Given $\mathcal F=\{S_1,\ldots,S_m\}\subseteq2^{[n]}$ and a fixed demand $D\geq1$, minimize $\sum_jx_j$ over integers $x_j\geq0$ satisfying $\sum_{j:i\in S_j}x_j\geq D$ for every $i$. Here a set may be selected repeatedly. This models redundant coverage or repeated resource provision.
The quantum algorithm is measured in QRAM time, classical lattice dynamic programming in RAM time, and the bounded-incidence encoding in bit operations.
\end{problem}
\speedupheading
\complexity{Quantum}{$\operatorname{poly}(m,n)^{\log n}T_D^n$, quantum lattice dynamic programming.}
For fixed $D$, the base satisfies $T_D<D+1$, with $T_2<2.660$, $T_3<3.529$ and $T_4<4.421$. The quasipolynomial prefactor is part of the bound. Coherent access to precomputed lattice states is assumed. Truncating coverage coordinates at $D$ retains all relevant states, and no set needs multiplicity above $D$.
\complexity{Classical}{$O(m(D+1)^n)$, demand-vector dynamic programming.}
The state space consists of capped coverage vectors in $\{0,\ldots,D\}^n$. Let $I=\sum_j|S_j|$ be the incidence count. When $m+I=O(n)$ and $D$ is fixed, a linear-size bounded-width Boolean encoding gives the stronger deterministic bound $\operatorname{poly}(m,n)2^{O_D(n/\sqrt{\log n})}$ with polynomial space. This includes the multiplicity variables, capped local coverage sums and objective threshold.
\dimensionlabel{Strength}{strengthPolynomial} Polynomial speedup through a smaller exponential constant for fixed demand in the general lattice comparison, using quantum search. Linear-incidence instances have a subexponential classical comparator.
\dimensionlabel{Evidence}{evidencePartial} Rigorous QRAM upper bound and partial evidence against the best general classical lattice algorithm.
\dimensionlabel{Hardware}{hardwareExponential} The active coverage register has $O(n\log(D+1))$ qubits, with polynomial workspace and exponential coherent preprocessing storage. The dependence on $D$ affects both the lattice size and reversible arithmetic.
{\small [Before ``ETH is false'', the lattice comparison used $O(m(D+1)^n)$ classically and $\operatorname{poly}(m,n)^{\log n}T_D^n$ quantumly, including the linear-incidence regime now covered by the subexponential bound.]\par}
\subhead{Algorithm}
\begin{algorithmindex}{9}
\item \href{https://arxiv.org/pdf/2104.14384}{Glos, Kokainis, Mori and Vihrovs, \emph{Quantum Speedups for Dynamic Programming on $n$-Dimensional Lattice Graphs} (MFCS 2021)}. Algorithm 1 and Theorem 5: lattice search; Sections 6--7 and Theorem 9: QRAM implementation and repeated-set multicover.
\item \href{https://github.com/openai/math/blob/main/preprints/The-Exponential-Time-Hypothesis-is-False-October-25-2026/paper.pdf}{OpenAI, \emph{The Exponential Time Hypothesis Is False} (2026)}. Theorem 1.1 and Appendix D, applied to the bounded-demand, linear-incidence encoding.
\end{algorithmindex}

\entrytitle{10}{Chromatic number}{strengthPolynomial}{evidencePartial}{hardwareExponential}
\setcounter{problem}{9}
\begin{problem}[Graph coloring and chromatic number]
Given $G=(V,E)$, find the smallest $k$ admitting $c:V\to[k]$ such that $c(u)\neq c(v)$ for every edge $uv$. The decision variant supplies $k$ and asks for such a coloring. Colors represent mutually incompatible time slots, frequencies, or resources.
The unrestricted quantum comparison uses QRAM time and classical RAM time, while fixed-color CNF algorithms are measured in bit operations.
\end{problem}
\speedupheading
\complexity{Quantum}{$O^*(1.7956^n)$, Gaspers--Li graph-coloring algorithm.}
Here $n=|V|$ and the number of colors is optimized without a fixed bound. The algorithm combines restricted-family preprocessing, maximal-independent-set enumeration and quantum minimum finding. Its table is accessed coherently in the QRAM model.
\complexity{Classical}{$O^*(2^n)$, inclusion--exclusion over independent sets; $\operatorname{poly}(n+|E|)2^{O_q(n\log\log n/\sqrt{\log n})}$ for fixed $q$-colorability, Boolean tensor contraction.}
The first bound computes the unrestricted chromatic number. For a fixed number $q$ of colors, one-hot vertex-color variables give $qn$ Boolean variables and fixed-width clauses for vertex and edge conditions, even on dense graphs. Exact counting and self-reduction recover a coloring in polynomial space. On sparse graphs the exponent becomes $O_q(n/\sqrt{\log n})$.
\dimensionlabel{Strength}{strengthPolynomial} Polynomial speedup through a smaller exponential constant for unrestricted chromatic number, based on amplitude amplification. Fixed-$q$ decision has the displayed subexponential classical comparator.
\dimensionlabel{Evidence}{evidencePartial} Rigorous quantum upper bound with partial evidence relative to the general classical chromatic-number algorithm. The fixed-$q$ comparison follows the explicit bounded-width encoding.
\dimensionlabel{Hardware}{hardwareExponential} Polynomial active subset and color registers coexist with exponentially many coherently accessible preprocessing values. Useful instances require fault tolerance; fixed-$q$ decision and unrestricted optimization have different relevant time comparisons.
\subhead{Algorithm}
\begin{algorithmindex}{10}
\item \href{https://arxiv.org/pdf/2311.08042}{Gaspers and Li, \emph{Quantum Algorithms for Graph Coloring and Other Partitioning, Covering, and Packing Problems} (ICALP 2024)}. Theorem 8 and the coloring construction following it; appendix Algorithms 1--3 give the preprocessing and recursive procedures. The preliminary covering formulation gives the classical inclusion--exclusion comparison.
\item \href{https://github.com/openai/math/blob/main/preprints/The-Exponential-Time-Hypothesis-is-False-October-25-2026/paper.pdf}{OpenAI, \emph{The Exponential Time Hypothesis Is False} (2026)}. Theorem 1.1 and Appendices D.1--D.4, applied to the one-hot coloring encoding.
\end{algorithmindex}

\entrytitle{11}{Minimum Steiner tree}{strengthPolynomial}{evidencePartial}{hardwareExponential}
\setcounter{problem}{10}
\begin{problem}[Minimum Steiner tree]
Given an undirected graph $G=(V,E)$, nonnegative edge weights $w$, and terminals $R\subseteq V$, $|R|=k$, find a connected subgraph containing $R$ of minimum total weight. An optimum may be chosen as a tree. Steiner vertices permit cheaper network connections than restricting the tree to terminals.
The quantum bound uses QRAM time, while classical bounds count RAM operations with the stated weight-arithmetic costs.
\end{problem}
\speedupheading
\complexity{Quantum}{$O^*(1.812^k)$, quantum Steiner-tree dynamic programming.}
The exponential parameter is the number $k$ of terminals, while polynomial factors include graph size and weight arithmetic. The construction uses a recurrence based on splitting a Steiner tree, together with partial preprocessing and nested quantum minimum finding. Its table values must be coherently accessible.
\complexity{Classical}{$(2+\delta)^k n^{f(1/\delta)}$ for fixed $\delta>0$, M\"olle--Richter--Rossmanith; $O^*(2^k)$ for bounded integer weights, subset convolution.}
The subset-convolution bound includes the dependence on the weight range. These improve the $O^*(3^k)$ Dreyfus--Wagner recurrence and the $O^*(2.684^k)$ splitting algorithm. For unit-weight sparse graphs, Gaspers' path-decomposition methods give vertex-parameter bounds $O^*(1.1661^n)$ on degree-three graphs and $O^*(1.3100^n)$ on degree-four graphs, using exponential memory. On such inputs the appropriate comparator is the minimum of the applicable terminal- and vertex-parameter bounds.
\dimensionlabel{Strength}{strengthPolynomial} Polynomial speedup through a smaller exponential constant in terminal count in the general comparison, using amplitude amplification. Sparse unit-weight regimes also have the stated classical vertex-parameter bounds.
\dimensionlabel{Evidence}{evidencePartial} Rigorous QRAM upper bound and partial evidence against the strongest displayed classical algorithms in the relevant weight and graph regimes.
\dimensionlabel{Hardware}{hardwareExponential} Polynomial active registers encode terminal subsets, vertices, costs and reconstruction data, together with exponential coherent preprocessing memory. Useful terminal counts require fault tolerance; graph and weight parameters affect the polynomial register and arithmetic costs.
\subhead{Algorithm}
\begin{algorithmindex}{11}
\item \href{https://arxiv.org/pdf/1904.03581}{Miyamoto, Iwamura, Kise and Le Gall, \emph{Quantum Speedup for the Minimum Steiner Tree Problem} (2019/2020)}. Theorem 1; Section 3.1 defines Dreyfus--Wagner, Section 3.2 gives the splitting recurrence, and Section 4 gives the quantum algorithm. Table 1 and its cited classical papers specify the weight-dependent comparisons.
\item \href{https://serggasp.github.io/pathwidth-site/papers/sparse-pathwidth/14-apps-connectivity.html}{Gaspers, \emph{Path Decompositions of Sparse Graphs via Gaussian Sweeps and Faster Exponential Algorithms} (2026)}. Theorem 14.1(a): sparse unit-weight Steiner tree.
\end{algorithmindex}

\entrytitle{12}{Graph bandwidth}{strengthPolynomial}{evidencePartial}{hardwareExponential}
\setcounter{problem}{11}
\begin{problem}[Graph bandwidth]
For $G=(V,E)$ with $n$ vertices, compute
\[\operatorname{bw}(G)=\min_{\pi:V\overset{\sim}{\longrightarrow}[n]}\max_{uv\in E}|\pi(u)-\pi(v)|.\]
Bandwidth minimizes the largest separation of adjacent vertices, equivalently the bandwidth of a symmetrically permuted adjacency matrix.
The quantum bound counts QRAM time and the classical bound counts RAM operations on the explicit graph.
\end{problem}
\speedupheading
\complexity{Quantum}{$O^*(2.946^n)$, quantum bandwidth dynamic programming.}
A decomposition into constrained vertex-ordering instances permits variable-time quantum search and partial subset preprocessing. The quoted time includes the search over decomposition choices and assumes coherent access to stored dynamic-programming values. Polynomial factors include graph processing and the objective threshold.
\complexity{Classical}{$O^*(4.383^n)$, Cygan--Pilipczuk exact bandwidth algorithm.}
The algorithm uses structural restrictions on feasible layouts to reduce the search space. The quantum construction also accelerates a simpler $O^*(5^n)$ formulation, while the $4.383^n$ bound supplies the stronger overall classical comparison. Both bounds use the number of vertices as their exponential parameter.
\dimensionlabel{Strength}{strengthPolynomial} Polynomial speedup from $4.383^n$ to $2.946^n$; quantum minimum finding and variable-time amplitude amplification are the principal acceleration tools.
\dimensionlabel{Evidence}{evidencePartial} The upper bound is rigorous in the stated memory model. The best-algorithm comparison has partial evidence from the established classical exact-algorithm literature.
\dimensionlabel{Hardware}{hardwareExponential} Polynomial active subset, layout and objective registers require exponential QRAM for preprocessing. The nested searches and stored table make useful large instances a fault-tolerant target.
\subhead{Algorithm}
\begin{algorithmindex}{12}
\item \href{https://arxiv.org/pdf/1807.05209}{Ambainis et al., \emph{Quantum Speedups for Exponential-Time Dynamic Programming Algorithms} (SODA 2019)}. Section 4, Theorem 6: bandwidth algorithm; the section specifies the classical recurrence and cites the Cygan--Pilipczuk $4.383^n$ algorithm.
\end{algorithmindex}

\entrytitle{13}{Exact treewidth}{strengthPolynomial}{evidencePartial}{hardwareExponential}
\setcounter{problem}{12}
\begin{problem}[Treewidth]
For a graph $G$, find a tree decomposition $(T,\{B_t\})$ minimizing $\max_t|B_t|-1$. Every vertex belongs to a bag, every edge has both endpoints in a bag, and the bags containing each vertex induce a connected subtree. Treewidth governs the complexity of many graph and constraint algorithms.
The polynomial-space quantum bound counts gate time, the exponential-memory quantum bounds use QRAM time, and classical bounds use RAM time.
\end{problem}
\speedupheading
\complexity{Quantum}{$O^*(1.538^n)$ with exponential QRAM, quantum treewidth dynamic programming.}
K\c{l}evickis--Pr\=usis--Vihrovs also give $O^*(1.554^n)$ time with $O^*(1.452^n)$ space, and $O^*(1.618^n)$ time with polynomial space. The fastest variant uses $O^*(1.538^n)$ space. These are distinct memory tradeoffs obtained from separator-based structural decompositions and quantum search.
\complexity{Classical}{$O^*(1.7347^n)$ with exponential space, potential-maximal-clique dynamic programming; $O^*(2.6151^n)$ with polynomial space, Fomin--Villanger.}
The first bound enumerates potential maximal cliques and optimizes compatible decomposition blocks. The polynomial-space method enumerates connected sets with constrained neighborhoods and recursively assembles decompositions. The comparison with the polynomial-space quantum algorithm therefore uses the $2.6151^n$ bound; the $1.538^n$ quantum algorithm uses the exponential-space comparison.
\dimensionlabel{Strength}{strengthPolynomial} Polynomial speedups through smaller exponential constants in both displayed memory regimes. The quantum constructions combine structural enumeration with amplitude amplification and minimum finding.
\dimensionlabel{Evidence}{evidencePartial} Rigorous time-space bounds and partial evidence relative to the best classical algorithms for the same space regime.
\dimensionlabel{Hardware}{hardwareExponential} The fastest quantum construction needs exponential coherent memory and polynomial active graph registers. The $1.618^n$ variant reduces storage to polynomial logical workspace but retains exponential coherent running time. Fault tolerance is needed in either useful regime.
\subhead{Algorithm}
\begin{algorithmindex}{13}
\item \href{https://arxiv.org/pdf/2202.08186}{K\c{l}evickis, Pr\=usis and Vihrovs, \emph{Quantum Speedups for Treewidth} (TQC 2022)}. Theorems 21, 22 and 24 give the three tradeoffs; Algorithm 4 gives the fastest construction. Algorithm 1 states the classical polynomial-space procedure.
\item \href{https://arxiv.org/pdf/1309.1559}{Fomin, Todinca and Villanger, \emph{Large Induced Subgraphs via Triangulations and CMSO} (SIAM Journal on Computing 2015)}. Proposition 5 states the $O(1.7347^n)$ potential-maximal-clique enumeration bound of Fomin--Villanger; the preceding triangulation preliminaries define the blocks used in classical treewidth dynamic programming.
\end{algorithmindex}

\entrytitle{14}{Pathwidth and vertex separation}{strengthPolynomial}{evidencePartial}{hardwareExponential}
\setcounter{problem}{13}
\begin{problem}[Pathwidth]
Minimize $\max_i|B_i|-1$ over path decompositions: bags $B_1,\ldots,B_r$ cover all vertices and edges, and each vertex occurs in consecutive bags. Equivalently, minimize over vertex orders the largest number of already placed vertices with a neighbor still unplaced. Pathwidth captures the live frontier in linear processing and memory allocation.
The quantum algorithm is measured in QRAM time and classical ordering algorithms in RAM operations.
\end{problem}
\speedupheading
\complexity{Quantum}{$O^*(1.817^n)$, quantum vertex-ordering dynamic programming.}
The vertex-separation objective admits a subset recurrence: each state records an already placed prefix and its boundary. Partial preprocessing and quantum minimum finding optimize the largest boundary along a full ordering, using exponential coherent storage.
\complexity{Classical}{$O^*(1.89^n)$, Kitsunai et al. directed-pathwidth algorithm applied to an undirected graph.}
Replace each edge by its two opposite arcs; this preserves the relevant vertex-separation objective. The algorithm exploits structural restrictions on orderings and improves the generic $O^*(2^n)$ subset recurrence. Its exponential parameter remains the original number of vertices.
\dimensionlabel{Strength}{strengthPolynomial} Polynomial speedup through a smaller exponential constant relative to the $1.89^n$ classical bound. The quantum acceleration comes from recursive minimum finding and amplitude amplification.
\dimensionlabel{Evidence}{evidencePartial} Rigorous QRAM upper bound with partial evidence from the classical exact-algorithm comparison.
\dimensionlabel{Hardware}{hardwareExponential} The active subset, boundary-count and reconstruction registers use polynomially many logical qubits. Exponential QRAM and long coherent recursive searches dominate useful-instance requirements.
\subhead{Algorithm}
\begin{algorithmindex}{14}
\item \href{https://arxiv.org/pdf/1807.05209}{Ambainis et al., \emph{Quantum Speedups for Exponential-Time Dynamic Programming Algorithms} (SODA 2019)}. Corollary 5 and the vertex-ordering recurrence, equation (3): pathwidth application.
\item \href{https://doi.org/10.1007/s00453-015-0015-9}{Kitsunai et al., \emph{Computing Directed Pathwidth in $O(1.89^n)$ Time} (Algorithmica)}. The directed-pathwidth algorithm and main running-time theorem; the symmetric-digraph specialization gives the undirected comparison.
\end{algorithmindex}

\entrytitle{15}{Cutwidth}{strengthPolynomial}{evidencePartial}{hardwareExponential}
\setcounter{problem}{14}
\begin{problem}[Cutwidth]
Find a vertex order $\pi$ minimizing
\[\max_{1\leq i<n}|\{uv\in E:\pi(u)\leq i<\pi(v)\text{ or }\pi(v)\leq i<\pi(u)\}|.\]
The objective is the largest number of edges crossing a prefix cut. It represents the largest simultaneous communication or wiring requirement along a linear layout.
The comparison uses quantum QRAM time and classical RAM time on an explicit graph.
\end{problem}
\speedupheading
\complexity{Quantum}{$O^*(1.817^n)$, quantum vertex-ordering dynamic programming.}
A prefix subset determines the number of crossing edges. The recurrence minimizes the maximum of this local cut size and the value of the preceding prefix. Its quantum implementation uses partial subset preprocessing and recursive minimum finding, with exponential coherent memory.
\complexity{Classical}{$O^*(2^n)$ time and space, subset dynamic programming; $O^*(4^n)$ time with polynomial space, divide and conquer.}
The subset algorithm processes prefixes in order of their cardinality and retains the best maximum-cut value for each prefix. These bounds concern arbitrary cutwidth as a function of vertex count. Algorithms parameterized by a small target cutwidth can instead be used on instances in that regime.
\dimensionlabel{Strength}{strengthPolynomial} Polynomial speedup from $2^n$ to $1.817^n$ in the exponential-memory comparison, based on amplitude amplification.
\dimensionlabel{Evidence}{evidencePartial} Rigorous template and problem upper bounds, with partial evidence against the displayed general classical bound.
\dimensionlabel{Hardware}{hardwareExponential} Polynomial active registers represent a subset and its cut size. Exponentially many stored dynamic-programming values must be coherently accessible, and useful searches require fault tolerance.
\subhead{Algorithm}
\begin{algorithmindex}{15}
\item \href{https://arxiv.org/pdf/1807.05209}{Ambainis et al., \emph{Quantum Speedups for Exponential-Time Dynamic Programming Algorithms} (SODA 2019)}. Corollary 5 and equation (3): quantum ordering recurrence and cutwidth application.
\item \href{https://doi.org/10.1007/s00224-011-9312-0}{Bodlaender et al., \emph{A Note on Exact Algorithms for Vertex Ordering Problems on Graphs} (2012)}. The dynamic-programming and polynomial-space divide-and-conquer constructions for vertex-ordering objectives.
\end{algorithmindex}

\entrytitle{16}{Minimum linear arrangement}{strengthPolynomial}{evidencePartial}{hardwareExponential}
\setcounter{problem}{15}
\begin{problem}[Minimum linear arrangement]
For $G=(V,E)$, compute $\min_\pi\sum_{uv\in E}|\pi(u)-\pi(v)|$ over bijections $\pi:V\to[n]$. The sum measures total edge length in a linear layout and models total wiring length or aggregate communication distance.
The quantum comparison uses QRAM time and the classical comparison uses RAM operations on the graph.
\end{problem}
\speedupheading
\complexity{Quantum}{$O^*(1.817^n)$, quantum vertex-ordering dynamic programming.}
The sum of edge lengths equals the sum of cut sizes over the prefixes of the ordering. This additive objective fits the ordering recurrence, with polynomial-time evaluation of each transition. The quantum algorithm uses nested minimum finding and coherent access to partially precomputed subset values.
\complexity{Classical}{$O^*(2^n)$, subset dynamic programming.}
For each subset, the recurrence minimizes its accumulated prefix-cut cost over choices of the last vertex. Straightforward implementations include an $O(m2^n)$ bound, with polynomial factors depending on the graph representation, and use exponential memory. Specialized graph classes, including trees, admit faster classical algorithms and are separate input regimes.
\dimensionlabel{Strength}{strengthPolynomial} Polynomial speedup through a smaller exponential constant, obtained from amplitude amplification and quantum minimum finding.
\dimensionlabel{Evidence}{evidencePartial} Rigorous QRAM upper bound; partial evidence relative to the best general classical exact-algorithm comparison.
\dimensionlabel{Hardware}{hardwareExponential} A subset needs $n$ active qubits and the objective needs $O(\log(nm))$ bits, with reversible workspace and exponential coherent table storage. Fault tolerance is needed for useful depths.
\subhead{Algorithm}
\begin{algorithmindex}{16}
\item \href{https://arxiv.org/pdf/1807.05209}{Ambainis et al., \emph{Quantum Speedups for Exponential-Time Dynamic Programming Algorithms} (SODA 2019)}. Corollary 5 and equation (3): additive vertex-ordering recurrence and minimum-linear-arrangement application.
\item \href{https://doi.org/10.1007/s00224-011-9312-0}{Bodlaender et al., \emph{A Note on Exact Algorithms for Vertex Ordering Problems on Graphs} (2012)}. The subset dynamic-programming construction for additive ordering objectives.
\end{algorithmindex}

\entrytitle{17}{Directed minimum feedback arc set}{strengthPolynomial}{evidencePartial}{hardwareExponential}
\setcounter{problem}{16}
\begin{problem}[Directed minimum feedback arc set]
For a directed graph $G=(V,A)$, find a minimum-cardinality $F\subseteq A$ such that $(V,A\setminus F)$ is acyclic. Equivalently, choose a total order of the vertices minimizing the number of arcs directed backward. Weighted variants minimize the total deleted weight. Applications include ranking with inconsistent pairwise preferences.
The quantum bound counts QRAM time and the classical bound counts RAM time, including polynomial arithmetic costs for binary-encoded weights.
\end{problem}
\speedupheading
\complexity{Quantum}{$O^*(1.728^n)$, quantum feedback-arc-set dynamic programming.}
The ordering formulation admits the stronger split recurrence used for travelling salesman. A partition of the vertices fixes the direction of every arc crossing the two ordered parts; the remaining costs are internal to the parts. Quantum minimum finding combines partially precomputed values over balanced splits.
\complexity{Classical}{$O^*(2^n)$, subset dynamic programming.}
An ordering prefix determines the cost of appending a vertex, and the recurrence minimizes this cost over all subsets. It applies to arbitrary directed graphs and to nonnegative weighted variants. Structural restrictions such as tournament input select more specialized classical algorithms and hence a different comparison.
\dimensionlabel{Strength}{strengthPolynomial} Polynomial speedup through a smaller exponential constant. The acceleration comes from quantum minimum finding and amplitude amplification over subset decompositions.
\dimensionlabel{Evidence}{evidencePartial} Rigorous QRAM upper bound and partial evidence against the general classical subset algorithm.
\dimensionlabel{Hardware}{hardwareExponential} Polynomial active subset and cost registers require exponential coherently accessible preprocessing memory. Weight arithmetic and long nested searches set the fault-tolerant implementation cost.
\subhead{Algorithm}
\begin{algorithmindex}{17}
\item \href{https://arxiv.org/pdf/1807.05209}{Ambainis et al., \emph{Quantum Speedups for Exponential-Time Dynamic Programming Algorithms} (SODA 2019)}. Section 7 and Corollary 9: directed feedback arc set and its split recurrence; the construction uses the quantum dynamic-programming routine of Algorithm 2.
\end{algorithmindex}

\entrytitle{18}{Worst-case subset sum}{strengthPolynomial}{evidencePartial}{hardwareExponential}
\setcounter{problem}{17}
\begin{problem}[Subset sum]
Given integers $a_1,\ldots,a_n,t$ in binary, find $x\in\{0,1\}^n$ with $\sum_i a_ix_i=t$, or decide that none exists. This entry concerns arbitrary inputs; the random density-one regime has separate, faster heuristic algorithms.
Quantum time is measured with coherent random-access memory and polynomial integer-arithmetic costs; classical bounds specify word-RAM operations and coefficient bit length.
\end{problem}
\speedupheading
\complexity{Quantum}{$O^*(2^{2n/7})$, Chukhin--Kulikov--Levitskii--Mihajlin; $\widetilde O(2^{n/3})$, Arvind--Schuler.}
The $2n/7$ bound is a 2026 preprint result obtained from a worst-case seven-list sum algorithm, using quantum walks and coherent data structures. The $n/3$ construction uses a smaller precomputed sum list and quantum search over the remaining variables. Both bounds suppress polynomial arithmetic factors and require exponential accessible memory.
\complexity{Classical}{$O(2^{0.49n})$ word operations for polynomial-bit inputs, OpenAI's worst-case subset-sum algorithm.}
For each fixed $c$, the bound holds on all inputs of bit length $b\le n^c$ for sufficiently large $n$, with success probability at least $2/3$ and a running-time cap on every random execution. The word size is $O(n+b+\log n)$. A lower-space alternative uses $O^*(2^{n/2})$ time and $O^*(2^{0.199n})$ writable words, in particular $O(2^{n/5})$ words under a fixed polynomial-bit convention. For nonnegative inputs and small target $t$, pseudopolynomial algorithms with near-linear dependence on $n+t$ can be faster.
\dimensionlabel{Strength}{strengthPolynomial} Polynomial speedup from $2^{0.49n}$ to $2^{2n/7}$. The fastest quantum result uses quantum walks; the $2^{n/3}$ algorithm uses amplitude amplification.
\dimensionlabel{Evidence}{evidencePartial} Worst-case upper bounds with partial evidence for a speedup over classical algorithms. The strongest displayed quantum and classical bounds are 2026 preprints with their stated memory and bit-length assumptions.
\dimensionlabel{Hardware}{hardwareExponential} Polynomial active registers hold subset indices and integer sums; exponentially large coherent list and search structures dominate storage. Long walk or amplification sequences require fault tolerance. The classical low-space tradeoff provides a separate memory-sensitive comparison.
\subhead{Algorithm}
\begin{algorithmindex}{18}
\item \href{https://arxiv.org/pdf/2608.07309}{Chukhin, Kulikov, Levitskii and Mihajlin, \emph{Improved Quantum Algorithms for Subset Sum and $k$-SUM} (2026)}. Corollary 1.3; Theorem 1.2 and Algorithm 1 give the underlying $k$-SUM routine; the Subset Sum application gives the block reduction.
\item \href{https://arxiv.org/pdf/quant-ph/0212048}{Arvind and Schuler, \emph{The Quantum Query Complexity of 0-1 Knapsack and Associated Claw Problems} (ISAAC 2003)}. Section 2, Theorem 1 and Corollary 2; steps 1--6 define the search construction.
\item \href{https://github.com/openai/math/tree/main/preprints/Subset-Sum-in-Time-2-power-0-49n-October-4-2026}{OpenAI, \emph{Subset Sum in Time $O(2^{0.49n})$} (2026)}. Theorem 1.1 and Section 6, especially Sections 6.1--6.3: full algorithm and bounded running time.
\item \href{https://github.com/openai/math/tree/main/preprints/A-Low-Space-Algorithm-for-Worst-Case-Subset-Sum-September-26-2026}{OpenAI, \emph{A Low-Space Algorithm for Worst-Case Subset Sum} (2026)}. Theorem 1.1, Corollary 1.2 and Section 7: uniform low-space algorithm.
\end{algorithmindex}

\entrytitle{19}{Zero-one knapsack}{strengthPolynomial}{evidencePartial}{hardwareExponential}
\setcounter{problem}{18}
\begin{problem}[Zero-one knapsack]
Given positive integer weights $w_i$, nonnegative profits $p_i$, and capacity $W$, maximize $\sum_i p_ix_i$ subject to $x_i\in\{0,1\}$ and $\sum_iw_ix_i\leq W$. The decision version also supplies a profit target. Knapsack models budgeted selection and resource allocation.
The meet-in-the-middle comparison uses quantum coherent-memory time and classical RAM time, with polynomial bit-arithmetic costs for weights and profits.
\end{problem}
\speedupheading
\complexity{Quantum}{$\widetilde O(2^{n/3})$, Arvind--Schuler integer-programming search.}
The capacity and profit threshold give a constant number of inequalities. Binary search over the profit range recovers the optimum with a polynomial bit-length factor. The construction requires coherent access to precomputed partial-sum and dominance-search data. QuantumTreeGenerator, a separate branch-and-bound construction, gives instance-dependent search estimates based on classical solver trees.
\complexity{Classical}{$O^*(2^{n/2})$, Horowitz--Sahni meet in the middle; $O(nW)$, capacity dynamic programming.}
The first bound enumerates partial weight-profit pairs and removes dominated pairs, using exponential memory. The second stores the best profit for each achievable capacity and is pseudopolynomial in the binary-encoded input. The smaller applicable expression supplies the comparison. If the numerical data have $B$ bits, a direct Boolean arithmetic encoding has size $O(n(B+\log n))$; any tensor-contraction bound applied to that encoding retains this dependence on $B$.
\dimensionlabel{Strength}{strengthPolynomial} Polynomial speedup through a smaller exponential constant in item count in the large-coefficient meet-in-the-middle regime, based on Grover search and amplitude amplification. Small-capacity instances have the pseudopolynomial classical comparator.
\dimensionlabel{Evidence}{evidencePartial} The few-inequality upper bound is rigorous in its access model. Evidence for practical speedup from branch-and-bound estimates depends on the solver trees and implementation costs.
\dimensionlabel{Hardware}{hardwareExponential} Active registers include $n$ assignment qubits, $O(\log W)$ weight bits and $O(\log\sum_i p_i)$ profit bits, plus reversible workspace. The asymptotic search construction additionally needs exponential coherent memory. Useful numerical instances require fault tolerance.
\subhead{Algorithm}
\begin{algorithmindex}{19}
\item \href{https://arxiv.org/pdf/quant-ph/0212048}{Arvind and Schuler, \emph{The Quantum Query Complexity of 0-1 Knapsack and Associated Claw Problems} (ISAAC 2003)}. Section 2, Theorem 1, steps 1--6, and Corollary 2: fixed-inequality feasibility and optimization by thresholds.
\item \href{https://www.nature.com/articles/s41534-025-01097-8}{\emph{A Quantum Algorithm for Solving 0-1 Knapsack Problems} (npj Quantum Information 2025)}. The \textsc{QuantumTreeGenerator} construction and the numerical comparison with COMBO search trees.
\end{algorithmindex}

\entrytitle{20}{Independent set, clique and vertex cover}{white}{white}{hardwareExponential}
\setcounter{problem}{19}
\begin{problem}[Maximum independent set, maximum clique and minimum vertex cover]
Find $I\subseteq V$ of maximum size with no edge inside $I$. Equivalent exact objectives are maximum clique in the complement graph and minimum vertex cover $V\setminus I$ in the original graph. These equivalences preserve vertex count and translate witnesses directly.
The comparison uses quantum gate complexity and classical bit complexity on explicit graph inputs.
\end{problem}
\speedupheading
\complexity{Quantum}{$O^*(2^{n/5})=O^*(1.1488^n)$, D\"orn's amplified branching algorithm.}
This bound applies to maximum independent set on an arbitrary $n$-vertex graph. Complementation transfers it to maximum clique, and taking the complementary vertex set gives minimum vertex cover. Running the classical counting algorithm below on a quantum computer also supplies its subexponential upper bound.
\complexity{Classical}{$\operatorname{poly}(n+|E|)2^{O(n\log\log n/\sqrt{\log n})}$, Boolean tensor contraction with cardinality constraints.}
Use one Boolean variable per vertex and a width-two clause for each forbidden selected edge. A linear-size Boolean cardinality circuit tests the target size, so the encoding uses $O(n)$ variables and fixed-width clauses even for a dense graph. Exact counting and self-reduction recover an optimum in polynomial space. For sparse graphs, maximum independent set and minimum vertex cover have the sharper exponent $O(n/\sqrt{\log n})$; maximum clique has this sharper bound when its complement is sparse.
\dimensionlabel{Strength}{white} The subexponential classical bound is asymptotically smaller than the displayed amplified-branching quantum bound. The classical exponent is $n^{1-o(1)}$, so its exponential-to-subexponential reduction has no power saving.
\dimensionlabel{Evidence}{white} The stated algorithms have rigorous upper bounds. D\"orn's displayed bound supplies no asymptotic quantum-speedup claim against the current classical comparator.
\dimensionlabel{Hardware}{hardwareExponential} The amplified branching method uses polynomial assignment and graph-evaluation workspace, with exponentially many coherent evaluations. Useful instances would require fault tolerance. The polynomial-space classical algorithm changes the relevant time comparison without adding an exponential memory requirement.
{\small [Before ``ETH is false'', the known general classical bound was $O^*(1.1996^n)$, and D\"orn's quantum algorithm claimed an improvement to $O^*(1.1488^n)$.]\par}
\subhead{Algorithm}
\begin{algorithmindex}{20}
\item \href{https://arxiv.org/pdf/quant-ph/0510084}{D\"orn, \emph{Quantum Complexity Bounds for Independent Set Problems} (2005)}. Algorithm 2 and Theorem 6: maximum independent set.
\item \href{https://arxiv.org/pdf/1312.6260}{Xiao and Nagamochi, \emph{Exact Algorithms for Maximum Independent Set} (Information and Computation 2017)}. The main maximum-independent-set algorithm and its $1.1996^n$ running-time result.
\item \href{https://github.com/anthropics/formal-math/tree/main/eth}{Claude (Anthropic), \emph{A $2^{o(n)}$ Time Algorithm for 3-SAT} (2026)}. Corollary 7: graph cardinality applications. \href{https://github.com/openai/math/blob/main/preprints/The-Exponential-Time-Hypothesis-is-False-October-25-2026/paper.pdf}{OpenAI, \emph{The Exponential Time Hypothesis Is False} (2026)}, Theorem 1.1 and Appendix D, supplies the explicit fixed-width time and space bounds.
\end{algorithmindex}

\entrytitle{21}{Domatic number}{strengthPolynomial}{evidencePartial}{hardwareExponential}
\setcounter{problem}{20}
\begin{problem}[Domatic number]
Given $G=(V,E)$, maximize $k$ such that $V$ can be partitioned into $D_1,\ldots,D_k$ with $N[D_i]=V$ for every $i$. Each class is a dominating set. The problem models disjoint teams of resources, each capable of providing coverage to the entire network.
Quantum complexity is measured in QRAM time, classical inclusion--exclusion in RAM time, and the fixed-part-count Boolean algorithms in bit operations.
\end{problem}
\speedupheading
\complexity{Quantum}{$O^*((2-\epsilon)^n)$ for some constant $\epsilon>0$, Gaspers--Li domatic-number algorithm.}
The theorem obtains a strict improvement below base two through restricted enumeration of minimal dominating sets and the covering framework. The analysis establishes the existence of $\epsilon$ rather than an optimized numerical base. The number of parts is optimized as part of the input problem, and preprocessing values are coherently accessible.
\complexity{Classical}{$O^*(2^n)$, inclusion--exclusion over dominating sets.}
For any fixed target $q$, the stronger bound is $O_{q,\eta}(2^{\eta n}\operatorname{poly}(n))$ for every fixed $\eta>0$. One-hot membership variables and closed-neighborhood domination clauses produce $O(qn)$ variables and $O(qn)$ possibly long clauses; long-clause shortening followed by fixed-width counting gives this bound. For fixed maximum degree, the same encoding has bounded width and linear size, yielding $\operatorname{poly}(n)2^{O_q(n/\sqrt{\log n})}$. These fixed-$q$ bounds do not treat an unbounded number of parts as a constant.
\dimensionlabel{Strength}{strengthPolynomial} Polynomial speedup through a smaller exponential constant for unrestricted domatic number, based on amplitude amplification. Fixed-$q$ decision has the subexponential classical comparison above.
\dimensionlabel{Evidence}{evidencePartial} A rigorous quantum upper bound with partial evidence against the general classical exact algorithm. The fixed-$q$ classical conclusion follows the explicit domination encoding.
\dimensionlabel{Hardware}{hardwareExponential} Polynomial active subset and part-count registers accompany exponential coherent preprocessing memory. The small unspecified base improvement still entails exponentially long coherent searches on large graphs.
\subhead{Algorithm}
\begin{algorithmindex}{21}
\item \href{https://arxiv.org/pdf/2311.08042}{Gaspers and Li, \emph{Quantum Algorithms for Graph Coloring and Other Partitioning, Covering, and Packing Problems} (ICALP 2024)}. Theorem 9: restricted minimal-dominating-set enumeration; Corollary 2: domatic number; Theorem 4 defines the quantum covering routine.
\item \href{https://github.com/anthropics/formal-math/tree/main/eth}{Claude (Anthropic), \emph{A $2^{o(n)}$ Time Algorithm for 3-SAT} (2026)}. Long-clause shortening; together with \href{https://github.com/openai/math/blob/main/preprints/The-Exponential-Time-Hypothesis-is-False-October-25-2026/paper.pdf}{OpenAI, \emph{The Exponential Time Hypothesis Is False}, Theorem 1.1 and Appendix D}, applied to the domination clauses above.
\end{algorithmindex}

\entrytitle{22}{Set packing}{strengthPolynomial}{evidencePartial}{hardwareExponential}
\setcounter{problem}{21}
\begin{problem}[Set packing]
Given $U$, $|U|=n$, a family $\mathcal F$ of nonempty subsets, and $k$, decide whether there exist $S_1,\ldots,S_k\in\mathcal F$ with $S_i\cap S_j=\varnothing$ for $i\neq j$. Maximizing $k$ selects the largest collection of mutually compatible objects. The family may be implicit, for example objects defined by graph properties.
The quantum framework uses QRAM time and its supplied enumeration routine, while explicit-family classical bounds use RAM or Boolean bit complexity.
\end{problem}
\speedupheading
\complexity{Quantum}{$O^*((2+c)^{n/2})$ for eligible enumeration base $c<2$, Gaspers--Li packing framework.}
The hypothesis is a simple branching algorithm that, for every restricted universe $X$, enumerates a representative family in $O^*(c^{|X|})$ time and preserves a nonstraddling feasible packing as specified in Theorem 4. Refined bounds apply when $c\le1.147899$; near $c=1$ the base approaches $1.7274$. These conditions concern the representation and enumeration of the set family, not just the number of its members.
\complexity{Classical}{$O^*(2^n)$, subset methods for an efficiently represented family.}
The bound includes the required family counting or membership costs. For an explicit family of $m$ sets, construct its conflict graph, with one vertex per set and an edge for each intersecting pair. Maximum independent set gives the stronger bound $\operatorname{poly}(n,m)2^{O(m\log\log m/\sqrt{\log m})}$. Thus $m=O(n)$ yields a subexponential comparator in universe size, while a large implicit family retains the stated enumeration-dependent framework comparison.
\dimensionlabel{Strength}{strengthPolynomial} Polynomial speedup through a smaller exponential constant when the representative-family hypothesis holds in the unrestricted-family regime. The quantum mechanism is amplitude amplification over partial packings.
\dimensionlabel{Evidence}{evidencePartial} Rigorous conditional-on-enumeration upper bound, with partial evidence against the classical universe-size algorithm. Explicit small families have the stronger conflict-graph comparator.
\dimensionlabel{Hardware}{hardwareExponential} Polynomial subset and feasibility registers require exponential coherently accessible preprocessing data. The reversible implementation of the enumeration and compatibility tests contributes to the gate cost.
\subhead{Algorithm}
\begin{algorithmindex}{22}
\item \href{https://arxiv.org/pdf/2311.08042}{Gaspers and Li, \emph{Quantum Algorithms for Graph Coloring and Other Partitioning, Covering, and Packing Problems} (ICALP 2024)}. Theorem 4 states the representative-family conditions and the general algorithm; Theorems 5--6 and Corollary 1 give the refined small-$c$ bounds.
\item \href{https://github.com/openai/math/blob/main/preprints/The-Exponential-Time-Hypothesis-is-False-October-25-2026/paper.pdf}{OpenAI, \emph{The Exponential Time Hypothesis Is False} (2026)}. Theorem 1.1 and Appendix D, composed with the conflict-graph and cardinality encoding of maximum independent set.
\end{algorithmindex}

\entrytitle{23}{Exact set partition}{strengthPolynomial}{evidencePartial}{hardwareExponential}
\setcounter{problem}{22}
\begin{problem}[Exact set partition]
Given $U$, a nonempty-set family $\mathcal F\subseteq2^U$, and integer $k$, find $S_1,\ldots,S_k\in\mathcal F$ such that $U=S_1\mathbin{\dot\cup}\cdots\mathbin{\dot\cup}S_k$. This is exact cover with a prescribed number of selected sets, and includes many clustering and assignment formulations.
The general comparison uses quantum QRAM time and classical subset-transform time; the explicit-family Boolean bounds count bit operations.
\end{problem}
\speedupheading
\complexity{Quantum}{$O^*((2+c)^{n/2})$ for eligible enumeration base $c<2$, Gaspers--Li partition framework.}
Here $n=|U|$. A restricted-family branching algorithm must enumerate representative members in $O^*(c^{|X|})$ time while preserving a nonstraddling feasible partition. Theorems 5--6 refine the base for small $c$. The complexity includes preprocessing and assumes coherent access to its results; polynomial factors account for membership and feasibility tests.
\complexity{Classical}{$O^*(2^n)$, inclusion--exclusion and subset transforms.}
For an explicit family of $m=\Theta(n)$ sets, selecting exactly one covering set per element and enforcing the chosen cardinality gives a stronger $O_\epsilon(2^{\epsilon n}\operatorname{poly}(n,m))$ bound for each fixed $\epsilon>0$. Long coverage clauses coexist with pairwise incompatibility constraints and a linear-size cardinality circuit. Long-clause shortening and bounded-width counting give the stated time. If the incidence count is linear, a linear-size encoding gives $\operatorname{poly}(n,m)2^{O(n/\sqrt{\log n})}$.
\dimensionlabel{Strength}{strengthPolynomial} Polynomial speedup through a smaller exponential constant for an eligible large or implicit family, from amplitude amplification. Linear-size explicit families have a subexponential classical comparison.
\dimensionlabel{Evidence}{evidencePartial} Rigorous framework theorem under its representative-family hypothesis, with partial evidence against the general classical subset-transform bound.
\dimensionlabel{Hardware}{hardwareExponential} Polynomial active subset and cardinality registers are supplemented by exponential coherent memory. Large-instance fault-tolerant costs also include checking disjointness, membership and total coverage.
\subhead{Algorithm}
\begin{algorithmindex}{23}
\item \href{https://arxiv.org/pdf/2311.08042}{Gaspers and Li, \emph{Quantum Algorithms for Graph Coloring and Other Partitioning, Covering, and Packing Problems} (ICALP 2024)}. Theorem 4: exact-partition framework; Theorems 5--6 and Corollary 1: improved small-$c$ implementation.
\item \href{https://github.com/anthropics/formal-math/tree/main/eth}{Claude (Anthropic), \emph{A $2^{o(n)}$ Time Algorithm for 3-SAT} (2026)}. Long-clause shortening; \href{https://github.com/openai/math/blob/main/preprints/The-Exponential-Time-Hypothesis-is-False-October-25-2026/paper.pdf}{OpenAI, \emph{The Exponential Time Hypothesis Is False}}, Theorem 1.1 and Appendix D, gives the fixed-width contraction routine used after shortening.
\end{algorithmindex}

\entrytitle{24}{Covering by maximal induced bicliques}{strengthPolynomial}{evidencePartial}{hardwareExponential}
\setcounter{problem}{23}
\begin{problem}[Vertex covering by maximal induced bicliques]
For a graph $G=(V,E)$ and $k$, decide whether $V$ is the union of $k$ vertex sets, each inducing a maximal complete bipartite subgraph. Maximality means that no extra vertex can be added while preserving this property. The covering is on vertices; overlapping bicliques are permitted. Packing and exact partition variants impose disjointness.
The covering framework uses quantum QRAM time and classical RAM time, while fixed-cover-count algorithms use Boolean bit complexity.
\end{problem}
\speedupheading
\complexity{Quantum}{$O^*(1.8554^n)$, Gaspers--Li biclique covering, packing and partition algorithms.}
Maximal induced bicliques admit the restricted enumeration used by the framework, with base $3^{1/3}$. The same theorem handles the three set-family variants once their covering, packing or partition feasibility condition is specified. Coherent preprocessing memory is assumed.
\complexity{Classical}{$O^*(2^n)$, inclusion--exclusion over the biclique family.}
For covering by a fixed number $q$ of bicliques, a stronger bound is $\operatorname{poly}(n)2^{O_q(n\log\log n/\sqrt{\log n})}$. Two side-membership bits per vertex per biclique impose induced complete-bipartite adjacency through fixed-width clauses, together with a covering clause per vertex. Each covering biclique may be extended to a maximal one, so the encoding decides the formal covering problem. This extension preserves covering and the fixed-$q$ comparison is stated for that variant.
\dimensionlabel{Strength}{strengthPolynomial} Polynomial speedup through a smaller exponential constant for the general family problems with unrestricted part count, using amplitude amplification. Fixed-$q$ covering has the subexponential classical bound.
\dimensionlabel{Evidence}{evidencePartial} Rigorous enumeration-based quantum upper bound and partial evidence against the general classical framework.
\dimensionlabel{Hardware}{hardwareExponential} Polynomial registers encode subsets, biclique sides and part counts. Exponential coherent table storage and long searches make useful cases fault-tolerant targets.
\subhead{Algorithm}
\begin{algorithmindex}{24}
\item \href{https://arxiv.org/pdf/2311.08042}{Gaspers and Li, \emph{Quantum Algorithms for Graph Coloring and Other Partitioning, Covering, and Packing Problems} (ICALP 2024)}. Theorem 7, maximal-biclique application; Theorem 4 defines the covering, packing and partition routines and their representative-family condition.
\item \href{https://github.com/openai/math/blob/main/preprints/The-Exponential-Time-Hypothesis-is-False-October-25-2026/paper.pdf}{OpenAI, \emph{The Exponential Time Hypothesis Is False} (2026)}. Theorem 1.1 and Appendix D, applied to the two-side fixed-$q$ covering encoding.
\end{algorithmindex}

\entrytitle{25}{Packing maximal induced matchings}{strengthPolynomial}{evidencePartial}{hardwareExponential}
\setcounter{problem}{24}
\begin{problem}[Packing and partitioning maximal induced matchings]
Let $\mathcal F$ contain all vertex sets $S\subseteq V$ such that every vertex of $G[S]$ has degree one and $S$ is inclusion-maximal with that property. Given $k$, find $k$ disjoint members of $\mathcal F$, or a partition of $V$ into $k$ such members. Each member is the endpoint set of an induced matching.
The quantum comparison uses QRAM time and classical subset-transform time, with bit complexity for the fixed-part-count Boolean encoding.
\end{problem}
\speedupheading
\complexity{Quantum}{$O^*(1.8934^n)$, Gaspers--Li maximal-induced-matching framework.}
Restricted maximal-induced-matching enumeration has base $10^{1/5}$. Combining it with partial subset preprocessing and quantum search gives the bound for packing or exact partition into an unrestricted number of members. Membership includes both induced matching and maximality.
\complexity{Classical}{$O^*(2^n)$, inclusion--exclusion and subset transforms.}
For every fixed number $q$ of members, a stronger bound is $O_{q,\epsilon}(2^{\epsilon n}\operatorname{poly}(n))$ for each fixed $\epsilon>0$. Vertex-membership and neighborhood-dominance bits encode degree one inside each member and exclusion of any further compatible edge. This uses $O(qn)$ variables and $O(qn)$ long clauses, together with polynomially many fixed-width clauses. Long-clause shortening followed by tensor contraction gives the bound; disjointness and complete coverage specify packing and partition respectively.
\dimensionlabel{Strength}{strengthPolynomial} Polynomial speedup through a smaller exponential constant for unrestricted part count, arising from amplitude amplification. A fixed part count has the stronger subexponential classical comparator.
\dimensionlabel{Evidence}{evidencePartial} The quantum upper bound follows the framework theorem and the maximal-induced-matching enumeration bound. The classical fixed-$q$ conclusion follows the explicit Boolean encoding.
\dimensionlabel{Hardware}{hardwareExponential} Polynomial active registers represent selected endpoints, local degree and maximality conditions, with exponential coherent preprocessing memory. Large useful instances require fault tolerance and long coherent searches.
\subhead{Algorithm}
\begin{algorithmindex}{25}
\item \href{https://arxiv.org/pdf/2311.08042}{Gaspers and Li, \emph{Quantum Algorithms for Graph Coloring and Other Partitioning, Covering, and Packing Problems} (ICALP 2024)}. Theorem 7, maximal-induced-matching application, using the Gupta--Raman--Saurabh enumeration bound; Theorem 4 gives the packing and partition procedures.
\item \href{https://github.com/anthropics/formal-math/tree/main/eth}{Claude (Anthropic), \emph{A $2^{o(n)}$ Time Algorithm for 3-SAT} (2026)}. Long-clause shortening; \href{https://github.com/openai/math/blob/main/preprints/The-Exponential-Time-Hypothesis-is-False-October-25-2026/paper.pdf}{OpenAI, \emph{The Exponential Time Hypothesis Is False}}, Theorem 1.1 and Appendix D: fixed-width counting after shortening.
\end{algorithmindex}

\entrytitle{26}{Vertex arboricity}{strengthPolynomial}{evidencePartial}{hardwareExponential}
\setcounter{problem}{25}
\begin{problem}[Vertex arboricity]
Find the smallest $k$ for which $V(G)=V_1\mathbin{\dot\cup}\cdots\mathbin{\dot\cup}V_k$ and every $G[V_i]$ is a forest. This invariant is vertex arboricity. Equivalently, cover the vertex set by $k$ maximal induced forests: enlarge each forest for the forward implication and assign overlapping vertices to one covering forest for the reverse implication.
The quantum algorithm is measured in QRAM time and the classical subset algorithm in RAM time.
\end{problem}
\speedupheading
\complexity{Quantum}{$O^*(1.9629^n)$, Gaspers--Li induced-forest covering algorithm.}
The framework uses a maximal-induced-forest enumeration procedure with base $1.8527$, combined with partial preprocessing and quantum search. Heredity permits deletion of overlaps from a forest covering to obtain a partition with the same number of parts. The input graph is explicit, and the number of parts is optimized.
\complexity{Classical}{$O^*(2^n)$, inclusion--exclusion over induced forests.}
Forest membership is polynomial-time decidable. Subset transforms aggregate admissible families and test covering feasibility; the covering-partition equivalence then gives the vertex-arboricity optimum. Classical memory is exponential at this time bound. Forest acyclicity is the defining global condition in this comparison.
\dimensionlabel{Strength}{strengthPolynomial} Polynomial speedup from $2^n$ to $1.9629^n$, based on amplitude amplification.
\dimensionlabel{Evidence}{evidencePartial} Rigorous quantum framework upper bound with partial evidence relative to the strongest displayed general classical subset method.
\dimensionlabel{Hardware}{hardwareExponential} Polynomial active graph and subset registers require exponential coherent preprocessing storage. The modest reduction in the exponential base still leaves exponentially long coherent computation; useful implementations require fault tolerance.
\subhead{Algorithm}
\begin{algorithmindex}{26}
\item \href{https://arxiv.org/pdf/2311.08042}{Gaspers and Li, \emph{Quantum Algorithms for Graph Coloring and Other Partitioning, Covering, and Packing Problems} (ICALP 2024)}. Theorem 7, maximal-induced-forest application; the preceding applications discussion gives the enumeration base and covering-partition equivalence, and Theorem 4 gives the quantum framework.
\end{algorithmindex}

\entrytitle{27}{Partition into hereditary graph classes}{strengthPolynomial}{evidencePartial}{hardwareExponential}
\setcounter{problem}{26}
\begin{problem}[Partition into a hereditary graph class]
For a fixed hereditary graph class $\mathcal H$, minimize $k$ such that $V=V_1\mathbin{\dot\cup}\cdots\mathbin{\dot\cup}V_k$ and $G[V_i]\in\mathcal H$. For cluster graphs, every connected component of $G[V_i]$ must be a clique, equivalently $G[V_i]$ contains no induced three-vertex path. This is a structured clustering or generalized coloring problem.
The unrestricted framework uses quantum QRAM time and classical RAM time, and fixed-part-count Boolean encodings use bit complexity.
\end{problem}
\speedupheading
\complexity{Quantum}{$O^*(1.9149^n)$ for cluster or triangle-free graphs, Gaspers--Li partition framework.}
The corresponding bases are $1.9365$ for cographs, claw-free graphs, trivially perfect graphs or threshold graphs; $1.9494$ for split graphs; and $1.9579$ for line graphs. Restricted enumeration of maximal induced members supplies the family routine. Heredity converts a covering into a partition by deleting overlaps. These bounds optimize the number of parts without treating it as a fixed constant.
\complexity{Classical}{$O^*(2^n)$, inclusion--exclusion and subset transforms.}
If the number of parts $q$ is fixed and the class has a finite forbidden-induced-subgraph description of maximum size $r$, a stronger bound is $\operatorname{poly}(n)2^{O_{q,r}(n\log\log n/\sqrt{\log n})}$. One-hot variables assign a part to each vertex, and a clause forbids every monochromatic forbidden copy. This uses $O(qn)$ variables, fixed width and polynomially many clauses. Exact counting and self-reduction recover the partition in polynomial space. This applies to the listed finite-forbidden-subgraph classes in the fixed-$q$ regime.
\dimensionlabel{Strength}{strengthPolynomial} Polynomial speedups through smaller exponential constants for unrestricted part count, from amplitude amplification. Fixed-$q$ decision has the subexponential classical comparison above.
\dimensionlabel{Evidence}{evidencePartial} Rigorous enumeration-based quantum upper bounds, with partial evidence against the classical general partition framework.
\dimensionlabel{Hardware}{hardwareExponential} Polynomial subset, part and forbidden-pattern registers accompany exponential coherent preprocessing memory. The largest forbidden-pattern size affects both enumeration cost and reversible membership testing.
\subhead{Algorithm}
\begin{algorithmindex}{27}
\item \href{https://arxiv.org/pdf/2311.08042}{Gaspers and Li, \emph{Quantum Algorithms for Graph Coloring and Other Partitioning, Covering, and Packing Problems} (ICALP 2024)}. Theorem 7 and its preceding list of forbidden-induced-subgraph classes; Theorem 4 defines the partition routine.
\item \href{https://github.com/openai/math/blob/main/preprints/The-Exponential-Time-Hypothesis-is-False-October-25-2026/paper.pdf}{OpenAI, \emph{The Exponential Time Hypothesis Is False} (2026)}. Theorem 1.1 and Appendix D, applied to the fixed-$q$, bounded-forbidden-pattern encoding.
\end{algorithmindex}

\entrytitle{28}{Minimum fill-in}{white}{white}{hardwareExponential}
\setcounter{problem}{27}
\begin{problem}[Minimum fill-in]
Given $G=(V,E)$, minimize $|F|$ over added-edge sets $F\subseteq\binom V2\setminus E$ such that $(V,E\cup F)$ is chordal, meaning every cycle of length at least four has a chord. Minimum fill-in controls additional nonzero entries created by sparse elimination.
The quantum ordering bound uses QRAM time, while the classical structural algorithm uses RAM time on the explicit graph.
\end{problem}
\speedupheading
\complexity{Quantum}{$O^*(1.817^n)$, quantum vertex-ordering dynamic programming.}
An elimination order determines a chordal completion by making each eliminated vertex's remaining neighborhood a clique. The ordering recurrence charges the newly inserted edges and optimizes the total cost. Partial preprocessing and quantum minimum finding provide the stated bound with exponential coherent memory.
\complexity{Classical}{$O^*(1.7347^n)$, potential-maximal-clique dynamic programming.}
The algorithm enumerates potential maximal cliques and minimal-separator blocks, then selects compatible triangulations with minimum fill cost. Its time improves on the generic $O^*(2^n)$ elimination-order recurrence, and its space is exponential. The graph structure exploited by this method gives the smaller base in the overall comparison.
\dimensionlabel{Strength}{white} The classical $1.7347^n$ bound is asymptotically smaller than the displayed quantum $1.817^n$ bound. The quantum result accelerates the subset-ordering recurrence.
\dimensionlabel{Evidence}{white} The ordering acceleration has a rigorous upper bound; the strongest displayed overall complexity bound is classical.
\dimensionlabel{Hardware}{hardwareExponential} The quantum construction requires polynomial active subset and fill-cost registers with exponential QRAM. Useful depths require fault tolerance. The stronger classical structural method also uses exponential storage.
\subhead{Algorithm}
\begin{algorithmindex}{28}
\item \href{https://arxiv.org/pdf/1807.05209}{Ambainis et al., \emph{Quantum Speedups for Exponential-Time Dynamic Programming Algorithms} (SODA 2019)}. Corollary 5 and equation (3): minimum-fill-in ordering application.
\item \href{https://arxiv.org/pdf/1309.1559}{Fomin, Todinca and Villanger, \emph{Large Induced Subgraphs via Triangulations and CMSO} (SIAM Journal on Computing 2015)}. Proposition 5 states the $1.7347^n$ potential-maximal-clique enumeration bound; the triangulation preliminaries define minimal separators, full blocks and potential maximal cliques used in minimum-fill dynamic programming.
\end{algorithmindex}

\entrytitle{29}{Minimum interval completion}{strengthPolynomial}{evidencePartial}{hardwareExponential}
\setcounter{problem}{28}
\begin{problem}[Minimum interval graph completion]
Given $G=(V,E)$, minimize $|F|$ such that $(V,E\cup F)$ is an interval graph: there must exist real intervals $I_v$ with $uv\in E\cup F$ exactly when $I_u\cap I_v\neq\varnothing$. The objective adds as few missing overlaps as possible and occurs in interval representations and ordering problems.
The quantum ordering bound uses QRAM time and the classical exact and parameterized bounds use RAM time.
\end{problem}
\speedupheading
\complexity{Quantum}{$O^*(1.817^n)$, quantum vertex-ordering dynamic programming.}
The ordering characterization assigns a local interval-completion cost to each vertex appended after a prefix. Recursive minimum finding combines partial subset values to optimize the total cost. The resulting ordering determines the completion, and the stated bound assumes exponential coherent preprocessing memory.
\complexity{Classical}{$O^*(2^n)$, subset dynamic programming; $k^{O(\sqrt k)}n^{O(1)}$, parameterized interval completion.}
The first bound uses vertex count $n$ as its exponential parameter. The second decides whether at most $k$ edges suffice and recovers a completion by exploiting the structure of a small fill set. For an instance with small optimum $k$, the parameterized expression supplies the stronger comparator. Threshold search can be used to obtain the exact optimum.
\dimensionlabel{Strength}{strengthPolynomial} Polynomial speedup through a smaller exponential constant in the unrestricted vertex-count comparison, based on amplitude amplification. Small-fill instances have the stated classical parameterized alternative.
\dimensionlabel{Evidence}{evidencePartial} Rigorous QRAM upper bound and partial evidence relative to the best displayed general classical bound.
\dimensionlabel{Hardware}{hardwareExponential} Polynomial active subset, objective and reconstruction registers accompany exponential coherent storage. Large-instance quantum computation requires fault tolerance; the fill parameter changes which classical bound should be used.
\subhead{Algorithm}
\begin{algorithmindex}{29}
\item \href{https://arxiv.org/pdf/1807.05209}{Ambainis et al., \emph{Quantum Speedups for Exponential-Time Dynamic Programming Algorithms} (SODA 2019)}. Corollary 5 and equation (3): interval-completion ordering application.
\item \href{https://arxiv.org/pdf/1402.3473}{Bliznets et al., \emph{A Subexponential Parameterized Algorithm for Interval Completion}}. Theorem 1.1 states the $k^{O(\sqrt k)}n^{O(1)}$ result; the algorithmic construction combines the structural enumeration and dynamic-programming sections.
\end{algorithmindex}

\entrytitle{30}{Binary integer programming with few inequalities}{strengthPolynomial}{evidencePartial}{hardwareExponential}
\setcounter{problem}{29}
\begin{problem}[Binary integer programming]
Given $A\in\mathbb Z^{d\times n}$ and $b\in\mathbb Z^d$, find $x\in\{0,1\}^n$ satisfying $Ax\leq b$, or decide infeasibility. An additional profit threshold converts a linear optimization problem into this form. The regime of interest has fixed or slowly growing $d$ and binary-encoded coefficients.
The few-inequality quantum bound uses coherent-memory time with bit-arithmetic costs, classical meet in the middle uses RAM time, and bounded-row-arity encodings use bit complexity.
\end{problem}
\speedupheading
\complexity{Quantum}{$\widetilde O(2^{n/3}n^d)$, Arvind--Schuler integer-programming algorithm.}
The polynomial arithmetic factors depend on coefficient bit length. For $d=o(n/\log n)$, the time is $2^{(1/3+o(1))n}$. Precomputed partial-sum vectors support coherent multidimensional range queries; both preprocessing and searching are included in the bound. With $B$-bit coefficients, direct arithmetic and comparisons have reversible cost polynomial in $dn(B+\log n)$.
\complexity{Classical}{$O^*(2^{n/2})$ for fixed $d$, meet in the middle with dominance searching.}
The two lists of partial-sum vectors are combined by coordinatewise inequalities and use exponential memory. For a different regime in which every row has support at most a fixed $r$, encode each forbidden local assignment as a clause using the original $n$ Boolean variables. Polynomially many rows then give $\operatorname{poly}(\text{input size})2^{O_r(n\log\log n/\sqrt{\log n})}$ deterministic time and polynomial space; linearly many rows improve the exponent to $O_r(n/\sqrt{\log n})$. A fixed number of dense, large-coefficient inequalities retains the few-constraint meet-in-the-middle comparison.
\dimensionlabel{Strength}{strengthPolynomial} Polynomial speedup through a smaller exponential constant for fixed $d$ and unrestricted row support, using Grover search and amplitude amplification. Bounded-row-arity instances have the stronger subexponential classical comparator.
\dimensionlabel{Evidence}{evidencePartial} Rigorous quantum upper bounds in the specified access model, with partial evidence against the best displayed classical algorithms in the few-dense-inequality regime.
\dimensionlabel{Hardware}{hardwareExponential} The active assignment register has $n$ qubits; $d$ arithmetic accumulators require input-dependent bit width. Exponential coherent range-search data and long searches dominate large-instance implementation. Increasing $d$ raises both the $n^d$ factor and arithmetic cost.
\subhead{Algorithm}
\begin{algorithmindex}{30}
\item \href{https://arxiv.org/pdf/quant-ph/0212048}{Arvind and Schuler, \emph{The Quantum Query Complexity of 0-1 Knapsack and Associated Claw Problems} (ISAAC 2003)}. Section 2, Theorem 1 and steps 1--6, pp.~3--4: partial-sum preprocessing, range-tree queries and quantum search; Corollary 2 gives threshold optimization.
\item \href{https://github.com/openai/math/blob/main/preprints/The-Exponential-Time-Hypothesis-is-False-October-25-2026/paper.pdf}{OpenAI, \emph{The Exponential Time Hypothesis Is False} (2026)}. Theorem 1.1 and Appendices D.1--D.4, applied to the bounded-row-arity clauses.
\end{algorithmindex}

\entrytitle{102}{Recursive Fourier sampling}{strengthPolynomialTime}{evidenceProved}{hardwarePolynomial}
\setcounter{problem}{101}
\begin{problem}[Recursive Fourier sampling]
Let $n\ge 6$ and $h\ge1$ be integers. A rooted tree of height $h$ has children $vx$ indexed by $x\in\{0,1\}^n$. Each internal vertex $v$ has an unknown string $s_v\in\{0,1\}^n$ and a Boolean function $g_v:\{0,1\}^n\to\{0,1\}$. Its answer bit is $b_v=g_v(s_v)$, and the promise is
\[
b_{vx}=s_v\cdot x\pmod 2
\qquad\text{for every internal vertex $v$ and every $x\in\{0,1\}^n$.}
\]
Given oracle access to the leaf bits and all the functions $g_v$, output the root bit $b_{\varnothing}$ with probability at least $2/3$. A vertex is identified by its path from the root; the oracle address length is $\ell=nh+O(\log h)$. This promise problem was introduced to study separations between quantum and classical computation. Complexity is measured in oracle queries, including queries to the functions $g_v$.
\end{problem}
\speedupheading
\complexity{Quantum}{$O(2^h)$ queries, recursive Bernstein--Vazirani Fourier sampling.}
The exact algorithm recursively recovers $s_v$ with Hadamard transforms, queries $g_v$, and uncomputes $s_v$. Johnson's clean implementation uses $2^{h+1}-1$ queries. For the restriction in which all $g_v$ equal a fixed nonparity function $g$, Aaronson proves the lower bound $\Omega((1-\mu(g))^{-h/2})$, where $\mu(g)$ is its nonparity coefficient. The function $g(s)=\mathbf 1[|s|\not\equiv0\pmod 3]$ has $\mu(g)=3/4-O(1/n)$, yielding $2^{\Omega(h)}$ queries and a matching $\Theta(2^h)$ bound when $h=o(n)$. The non-oracle gate count has an additional factor polynomial in $n$ and $h$.

\complexity{Classical}{$O(n^h)$ queries, recursive classical Fourier sampling.}
At each vertex, recursively evaluating $n$ linearly independent children recovers $s_v$; one query to $g_v$ then determines its answer. This uses $\sum_{j=0}^{h}n^j$ queries. For success probability $1/2+\epsilon$, Johnson proves the lower bound $\epsilon\prod_{j=1}^{h}(n-j\log_2 n)$ in the positive-factor range $h\log_2 n<n$. For constant $\epsilon>0$ and $h=o(n/\log n)$, these bounds give randomized query complexity $2^{\Theta(h\log n)}$.

\dimensionlabel{Strength}{strengthPolynomialTime} Subexponential to polynomial in the strongest comparison: throughout $h=\omega(1)$ and $h=O(\log n)$, the quantum cost is polynomial in the oracle address length $\ell$, while the classical cost is superpolynomial. At larger depths within the stated lower-bound range, both costs become superpolynomial and the general bounds describe an intermediate improvement. Fixed depth gives a polynomial query advantage.

\dimensionlabel{Evidence}{evidenceProved} Unconditional randomized and quantum query lower bounds for the same oracle problem family establish the separation.

\dimensionlabel{Hardware}{hardwarePolynomial} The recursive circuit uses $O(nh)=O(\ell)$ logical qubits outside the oracle implementation and $O(2^h)$ coherent oracle calls, with polynomial non-oracle gate overhead. Oracle implementation determines the additional workspace and gate cost.

\subhead{Algorithm}
\begin{algorithmindex}{102}
\item \href{https://doi.org/10.1137/S0097539796300921}{Bernstein and Vazirani, \emph{Quantum Complexity Theory} (1997)}. Section 8.4 introduces recursive Fourier sampling and the quantum--classical oracle separation.
\item \href{https://citeseerx.ist.psu.edu/document?doi=ee55b07d326ac3cb422fa755aae4484ab3b2ad2d\&repid=rep1\&type=pdf}{Johnson, \emph{Upper and Lower Bounds for Recursive Fourier Sampling}, Ph.D. thesis, University of California, Berkeley (2008)}. Section 1.2.3, starting on page 5, defines the generalized oracle problem; Algorithm 3, page 42, gives the clean quantum recursion; Algorithm 9, page 51, gives the classical recursion; Section 3.4, starting on page 59, gives the randomized query lower bound.
\item \href{https://arxiv.org/pdf/quant-ph/0209060}{Aaronson, \emph{Quantum Lower Bound for Recursive Fourier Sampling} (2003)}. Section 1 and Lemma 1: arbitrary-depth problem and quantum recursion; Theorem 5 and Proposition 6: quantum lower bound.
\end{algorithmindex}

\clearpage
\section{Graph algorithms and discrete optimization}

\entrytitle{31}{Minimum spanning tree}{strengthPolynomial}{evidenceProved}{hardwarePolynomial}
\setcounter{problem}{30}
\begin{problem}[Minimum spanning tree]
Given a connected undirected graph $G=(V,E)$, $|V|=n$, $|E|=m$, and weights $w:E\to\mathbb R$, output a spanning tree $T$ minimizing $\sum_{e\in T}w(e)$. This is the minimum-cost network connecting all vertices.
The comparison counts weight-oracle queries; explicit-input classical time counts comparisons and RAM operations, and quantum gate cost includes reversible weight access.
\end{problem}

\speedupheading
\complexity{Quantum}{$\Theta(n^{3/2})$ matrix queries or $\Theta(\sqrt{nm})$ array queries, D\"urr--Heiligman--H{\o}yer--Mhalla.}
The adjacency array supplies each vertex degree and coherent access to its indexed weighted neighbors. Both bounds allow bounded error and have polylogarithmic computational overhead when graph access and weight comparison are efficient. The output consists of $n-1$ classical edges.

\complexity{Classical}{$O(m)$ expected time, Karger--Klein--Tarjan; $\Theta(n^2)$ matrix queries.}
The randomized linear-time algorithm uses an explicit edge list and weight comparisons. In the oracle representation, inspecting the edge weights requires $\Omega(m)$ queries in the worst case. Thus the matrix comparison is $n^2$ against $n^{3/2}$; the array advantage is a factor $\sqrt{m/n}$. For $b$-bit weights, arithmetic and comparison costs must be included when converting either query bound to bit or gate complexity.

\dimensionlabel{Strength}{strengthPolynomial} Polynomial: a factor $\sqrt n$ for matrix input and $\sqrt{m/n}$ for array input. The quantum improvement comes from amplitude amplification and minimum finding inside Bor\r uvka phases.

\dimensionlabel{Evidence}{evidenceProved} The classical and quantum query bounds are tight for their respective input representations.

\dimensionlabel{Hardware}{hardwarePolynomial} Fault-tolerant graph access is required at useful scale. Active search and weight registers use $O(\log n+b)$ logical qubits, in addition to reversible-oracle workspace and $O(n(\log n+b))$ classical component records.

\subhead{Algorithm}
\begin{algorithmindex}{31}
\item \href{https://arxiv.org/pdf/quant-ph/0401091}{D\"urr, Heiligman, H{\o}yer and Mhalla, \emph{Quantum Query Complexity of Some Graph Problems} (2004/2006)}. Section 4, ``Minimum Spanning Tree'', Theorem 8; Section 8, Theorem 24, for the lower bound.
\item \href{https://doi.org/10.1145/201019.201022}{Karger, Klein and Tarjan, \emph{A Randomized Linear-Time Algorithm to Find Minimum Spanning Trees} (1995)}. The minimum-spanning-tree algorithm and its expected linear-time analysis.
\end{algorithmindex}


\entrytitle{32}{Undirected graph connectivity}{strengthPolynomial}{evidenceProved}{hardwarePolynomial}
\setcounter{problem}{31}
\begin{problem}[Undirected graph connectivity]
Given an undirected graph $G$, decide whether $\forall u,v\in V$ there is a path from $u$ to $v$; optionally return its connected components and a spanning forest. Connectivity is a feasibility condition in network design.
The quantum and classical bounds count adjacency-matrix or neighbor-array queries; explicit graph traversal time is also stated.
\end{problem}

\speedupheading
\complexity{Quantum}{$\Theta(n^{3/2})$ matrix queries or $\Theta(n)$ array queries, quantum graph exploration.}
The array model includes the vertex degrees and indexed coherent access to each neighbor list. The algorithm can also output connected components and a spanning forest with the same asymptotic query bound. The forest has $O(n)$ edges.

\complexity{Classical}{$\Theta(n^2)$ matrix queries; $O(n+m)$ explicit-input time, breadth-first or depth-first search.}
The matrix lower bound follows from distinguishing two disconnected components from a graph with one hidden connecting edge. Array traversal visits every discovered vertex and its incident edges. The quantum linear-query result uses the symmetry of undirected neighbor access and the supplied degrees.

\dimensionlabel{Strength}{strengthPolynomial} Polynomial: $n^2\to n^{3/2}$ in the matrix model and $n+m\to n$ in the array comparison. Quantum search accelerates structured graph exploration.

\dimensionlabel{Evidence}{evidenceProved} A rigorous query separation, with matching quantum upper and lower bounds.

\dimensionlabel{Hardware}{hardwarePolynomial} Fault-tolerant implementation uses $O(\log n)$ active vertex registers, coherent neighbor-oracle workspace and $O(n\log n)$ classical forest records. Constructing or accessing the graph oracle contributes to the gate count.

\subhead{Algorithm}
\begin{algorithmindex}{32}
\item \href{https://arxiv.org/pdf/quant-ph/0401091}{D\"urr, Heiligman, H{\o}yer and Mhalla, \emph{Quantum Query Complexity of Some Graph Problems} (2004/2006)}. Section 5, ``Connectivity'', Theorems 9 and 11; Section 8, Theorem 23, for the matrix lower bound.
\end{algorithmindex}


\entrytitle{33}{Strong connectivity of directed graphs}{strengthPolynomial}{evidenceProved}{hardwarePolynomial}
\setcounter{problem}{32}
\begin{problem}[Strong connectivity of directed graphs]
For a directed graph $G=(V,E)$, decide whether every ordered pair $(u,v)$ is joined by a directed path. A constructive version returns the strongly connected components. Applications include reachability and recurrent states in transition systems.
The comparison counts adjacency queries, with classical RAM time stated for explicit adjacency lists.
\end{problem}

\speedupheading
\complexity{Quantum}{$\Theta(n^{3/2})$ matrix queries and $O(\sqrt{nm}\log n)$ array queries, quantum directed graph exploration.}
The array bound has lower bound $\Omega(\sqrt{nm})$. Classical discovery orders and component records support the quantum searches, and the output can list the strongly connected components. The graph representation supplies outgoing neighbors; this access model determines the array bound.

\complexity{Classical}{$O(n+m)$ time, Tarjan or Kosaraju; $\Theta(n^2)$ matrix queries.}
Tarjan uses depth-first discovery indices and low-link information, while Kosaraju uses two directed traversals. Both produce an explicit component partition. With matrix input, deciding whether a hidden arc closes a one-way cut gives the quadratic classical query requirement.

\dimensionlabel{Strength}{strengthPolynomial} Polynomial, from Grover search and quantum minimum finding. The dense matrix improvement is a factor $\sqrt n$.

\dimensionlabel{Evidence}{evidenceProved} Rigorous oracle advantage. The matrix bound is tight, and the array bound is tight up to logarithmic factors.

\dimensionlabel{Hardware}{hardwarePolynomial} Fault-tolerant search uses logarithmic vertex-address registers and oracle workspace, with $O(n\log n)$ classical component data. The resource requirement is dominated by repeated coherent access to the directed graph.

\subhead{Algorithm}
\begin{algorithmindex}{33}
\item \href{https://arxiv.org/pdf/quant-ph/0401091}{D\"urr, Heiligman, H{\o}yer and Mhalla, \emph{Quantum Query Complexity of Some Graph Problems} (2004/2006)}. Section 6, ``Strong Connectivity'', including Theorem 13; Section 8, Theorem 22.
\item \href{https://doi.org/10.1137/0201010}{Tarjan, \emph{Depth-First Search and Linear Graph Algorithms} (1972)}. The strongly connected components algorithm.
\end{algorithmindex}


\entrytitle{34}{Single-source shortest paths}{strengthPolynomial}{evidenceProved}{hardwarePolynomial}
\setcounter{problem}{33}
\begin{problem}[Single-source shortest paths]
Given a directed graph with nonnegative edge lengths and a source $s$, output $d(v)=\min_{P:s\leadsto v}\sum_{e\in P}w(e)$ for all vertices, together with optional predecessor edges. Unreachable vertices receive distance $+\infty$.
The oracle comparison counts weighted adjacency queries; explicit-input times use the comparison-addition model, with bit or gate costs depending on weight precision.
\end{problem}

\speedupheading
\complexity{Quantum}{$O(n^{3/2}\log^2 n)$ matrix queries or $O(\sqrt{nm}\log^2 n)$ array queries, quantum shortest-path search.}
The corresponding lower bounds are $\Omega(n^{3/2})$ and $\Omega(\sqrt{nm})$. The algorithms return all distances and can retain predecessor edges. Their computational overhead includes arithmetic on distances and reversible evaluation of the weighted graph oracle.

\complexity{Classical}{$O(m\log^{2/3}n)$ time, Duan--Mao--Mao--Shu--Yin; $\Theta(n^2)$ dense-matrix queries.}
The deterministic 2025 algorithm treats directed nonnegative real weights in the comparison-addition model. Taking its minimum with Dijkstra using Fibonacci heaps, $O(m+n\log n)$, gives a useful bound across graph densities. For matrix-oracle input, the tight classical query complexity remains quadratic, so the dense oracle comparison is $n^2$ versus $\widetilde O(n^{3/2})$. Weight precision determines the bit complexity of additions and comparisons.

\dimensionlabel{Strength}{strengthPolynomial} Polynomial oracle speedup, originating in amplitude amplification and minimum finding. The explicit sparse-graph comparison uses the smaller of the two stated classical time bounds.

\dimensionlabel{Evidence}{evidenceProved} Rigorous query separation, with quantum lower bounds matching the main powers of $n$ and $m$.

\dimensionlabel{Hardware}{hardwarePolynomial} Fault tolerance is required for long coherent searches. Active registers use $O(\log n+b)$ logical qubits for $b$-bit distance arithmetic, plus reversible arithmetic and oracle workspace; classical distance records occupy $O(n(\log n+b))$ bits.

\subhead{Algorithm}
\begin{algorithmindex}{34}
\item \href{https://arxiv.org/pdf/quant-ph/0401091}{D\"urr, Heiligman, H{\o}yer and Mhalla, \emph{Quantum Query Complexity of Some Graph Problems} (2004/2006)}. Section 7, ``Single Source Shortest Paths'', Theorem 18 and Table 1; Section 8, Theorem 24.
\item \href{https://arxiv.org/pdf/2504.17033}{Duan, Mao, Mao, Shu and Yin, \emph{Breaking the Sorting Barrier for Directed Single-Source Shortest Paths} (2025)}. Theorem 1.1.
\end{algorithmindex}


\entrytitle{35}{Maximum-cardinality matching}{white}{white}{hardwarePolynomial}
\setcounter{problem}{34}
\begin{problem}[Maximum-cardinality matching]
Output $M\subseteq E$ maximizing $|M|$, subject to every vertex occurring in at most one edge of $M$. The input may be bipartite or general. Matching models assignment and resource allocation.
The comparison uses quantum time with coherent adjacency-array access and classical RAM or field-arithmetic time on explicit graphs.
\end{problem}

\speedupheading
\complexity{Quantum}{$\widetilde O(n\sqrt{m+n})$ for bipartite graphs; $\widetilde O(n^{3/2}\sqrt m+n^2)$ for general graphs, Ambainis--\v Spalek.}
The matching is stored classically, and quantum access supports augmenting-path and blossom operations. These are time upper bounds for explicit classical output under coherent adjacency access.

\complexity{Classical}{$m^{1+o(1)}$ for bipartite graphs, almost-linear maximum flow; $\min\{O(m\sqrt n),O(n^{9/4+\varepsilon})\}$ for general graphs.}
Bipartite cardinality matching reduces to unit-capacity flow. For general graphs, Micali--Vazirani supplies the sparse combinatorial bound and randomized algebraic matching uses $O(n^\omega)$ field operations. The bound $\omega\le9/4$ yields the displayed algebraic time for every fixed $\varepsilon>0$. The minimum of these classical bounds is asymptotically at least as strong as the stated quantum bounds.

\dimensionlabel{Strength}{white} The displayed quantum algorithms have no current worst-case time advantage over these classical algorithms.

\dimensionlabel{Evidence}{white} The runtime upper bounds are rigorous; the present comparison assigns no quantum-speedup claim.

\dimensionlabel{Hardware}{hardwarePolynomial} Fault-tolerant coherent adjacency access is needed. Search registers are logarithmic in $n$, while matching, blossom and layered-graph records require polynomial classical storage with reversible query support.

\subhead{Algorithm}
\begin{algorithmindex}{35}
\item \href{https://arxiv.org/pdf/quant-ph/0508205}{Ambainis and \v Spalek, \emph{Quantum Algorithms for Matching and Network Flows} (2005/2006)}. Theorems 3 and 4: bipartite and general matching.
\item \href{https://arxiv.org/pdf/2203.00671}{Chen et al., \emph{Maximum Flow and Minimum-Cost Flow in Almost-Linear Time} (2022)}. Main maximum-flow and minimum-cost-flow theorem.
\item \href{https://github.com/openai/math/blob/main/preprints/Matrix-Multiplication-Nine-Fourths-October-2-2026/paper.pdf}{OpenAI, \emph{An Upper Bound of $9/4$ for the Matrix Multiplication Exponent} (2026)}. Main matrix-multiplication theorem; \href{https://github.com/jzuiddam/omega-nine-quarters-all-fields}{Zuiddam, extension to every field (2026)}.
\end{algorithmindex}


\entrytitle{36}{Integral maximum flow}{white}{white}{hardwarePolynomial}
\setcounter{problem}{35}
\begin{problem}[Integral maximum flow]
Given directed capacities $c_e\in\{0,\ldots,U\}$ and terminals $s,t$, maximize the net outflow from $s$ over $0\le f_e\le c_e$, subject to flow conservation at $V\setminus\{s,t\}$. This is a central network optimization primitive.
Complexity counts quantum time with coherent residual-graph access and classical RAM time for polynomially bounded integral capacities.
\end{problem}

\speedupheading
\complexity{Quantum}{$\widetilde O(\min\{n^{7/6}\sqrt m\,U^{1/3},m\sqrt{nU}\})$, Ambainis--\v Spalek.}
This bound assumes $U\le n^{1/4}$. The flow is returned as an explicit classical edge array. Residual-capacity updates and subsequent coherent queries are part of the implementation.

\complexity{Classical}{$m^{1+o(1)}$, Chen--Kyng--Liu--Peng--Probst Gutenberg--Sachdeva and deterministic van den Brand et al.}
The almost-linear bound applies to polynomially bounded integer capacities. Interior-point optimization is combined with dynamic graph data structures for finding improving flow directions. This bound is asymptotically stronger than the displayed quantum blocking-flow bound for the same capacity regime.

\dimensionlabel{Strength}{white} The current comparison supplies no general quantum time advantage.

\dimensionlabel{Evidence}{white} Both algorithmic bounds are rigorous; the quantum theorem is an upper bound for a specific flow method.

\dimensionlabel{Hardware}{hardwarePolynomial} Fault-tolerant implementation uses $O(\log n+\log U)$ search and arithmetic registers plus oracle workspace. Residual-flow data occupy $O(m\log U)$ classical bits and must support coherent queries after updates.

\subhead{Algorithm}
\begin{algorithmindex}{36}
\item \href{https://arxiv.org/pdf/quant-ph/0508205}{Ambainis and \v Spalek, \emph{Quantum Algorithms for Matching and Network Flows} (2005/2006)}. Theorem 6 and the blocking-flow procedure following it.
\item \href{https://arxiv.org/pdf/2203.00671}{Chen et al., \emph{Maximum Flow and Minimum-Cost Flow in Almost-Linear Time} (2022)}. Main algorithmic theorem.
\item \href{https://arxiv.org/pdf/2309.16629}{van den Brand et al., \emph{A Deterministic Almost-Linear Time Algorithm for Minimum-Cost Flow} (2023)}. Main minimum-cost-flow theorem.
\end{algorithmindex}


\entrytitle{37}{Unweighted triangle finding}{strengthPolynomial}{evidenceProved}{hardwarePolynomial}
\setcounter{problem}{36}
\begin{problem}[Unweighted triangle finding]
Given oracle access to $A\in\{0,1\}^{n\times n}$ for a simple undirected graph, find distinct $u,v,w$ with $A_{uv}A_{vw}A_{wu}=1$, or report that none exist. Triangles encode local clustering and small constraint witnesses.
The primary comparison counts adjacency-oracle queries; the explicit-input classical comparator counts arithmetic operations.
\end{problem}

\speedupheading
\complexity{Quantum}{$\widetilde O(n^{5/4})$ queries, Le Gall\textquotesingle s combinatorial quantum-walk algorithm.}
The algorithm outputs a triangle with bounded error whenever one exists. The query count includes sampled neighborhoods and nested quantum walks. A circuit implementation additionally pays for updating and querying the stored subsets and neighborhoods.

\complexity{Classical}{$\Theta(n^2)$ entry queries; $O(n^{9/4+\varepsilon})$ arithmetic operations, algebraic triangle detection.}
Reading the matrix attains the classical query upper bound, and a hidden isolated triangle gives a quadratic worst-case query lower bound. Explicit-input triangle detection uses fast matrix multiplication with exponent $\omega\le9/4$, for every fixed $\varepsilon>0$; sparse inputs permit additional combinatorial refinements.

\dimensionlabel{Strength}{strengthPolynomial} Polynomial query improvement, $n^2\to\widetilde O(n^{5/4})$. The mechanism combines quantum walks and combinatorial preprocessing.

\dimensionlabel{Evidence}{evidenceProved} Rigorous separation from optimal classical entry-query complexity.

\dimensionlabel{Hardware}{hardwarePolynomial} Fault-tolerant quantum walks require coherent subset registers and neighborhood data of polynomial size. A logical-qubit estimate depends on the stored-data representation and graph oracle; the query theorem itself fixes their query count.

\subhead{Algorithm}
\begin{algorithmindex}{37}
\item \href{https://arxiv.org/pdf/1407.0085}{Le Gall, \emph{Improved Quantum Algorithm for Triangle Finding via Combinatorial Arguments} (2014)}. Theorem 1.1 and the algorithm in Sections 3--4.
\item \href{https://github.com/openai/math/blob/main/preprints/Matrix-Multiplication-Nine-Fourths-October-2-2026/paper.pdf}{OpenAI, \emph{An Upper Bound of $9/4$ for the Matrix Multiplication Exponent} (2026)}. Main matrix-multiplication theorem.
\end{algorithmindex}


\entrytitle{38}{Graph collision}{strengthPolynomial}{evidenceProved}{hardwarePolynomial}
\setcounter{problem}{37}
\begin{problem}[Graph collision]
A graph $G=([n],E)$ is known in advance; oracle bits $x_1,\ldots,x_n$ mark vertices. Decide whether $\exists\{u,v\}\in E$ with $x_u=x_v=1$, and find such an edge. The graph structure is free input, whereas marks are queried.
Complexity counts queries to the mark bits; the known graph is available independently of this oracle.
\end{problem}

\speedupheading
\complexity{Quantum}{$\widetilde O(n^{11/17})$ queries, Cornelissen et al.; $O(\sqrt n\,\alpha(G)^{1/6})$ for the independence-number refinement.}
The November 2026 algorithm splits high- and low-degree neighborhoods and combines quantum search with sampled graph structure. Here $\alpha(G)$ is the independence number. The earlier general quantum-walk construction uses $O(n^{2/3})$ queries. The quantum lower bound is $\Omega(\sqrt{n\log n})$ over an appropriate family of known graphs.

\complexity{Classical}{$\Theta(n)$ mark queries, exhaustive reading and randomized search.}
Reading all marks permits direct checking of the known edge set. A star with its center marked embeds an unstructured search over the other marks, giving the linear classical lower bound. Graph preprocessing and checking a stored edge relation have additional time costs outside the mark-query measure.

\dimensionlabel{Strength}{strengthPolynomial} Polynomial: $n\to\widetilde O(n^{11/17})$ in the worst case. The algorithms use quantum search, quantum walks and graph structure.

\dimensionlabel{Evidence}{evidenceProved} The quantum upper bound and classical lower bound give a rigorous query separation; the strongest cited bound is a November 2026 preprint.

\dimensionlabel{Hardware}{hardwarePolynomial} Fault-tolerant coherent processing is required. The $O(n^{2/3})$ walk uses roughly $\widetilde O(n^{2/3})$ coherent bits for stored labels and marks. The $11/17$ query exponent has its own implementation-dependent workspace costs.

\subhead{Algorithm}
\begin{algorithmindex}{38}
\item \href{https://arxiv.org/pdf/2611.05356}{Cornelissen et al., \emph{Breaking the Quantum Barriers for Graph Collision} (2026)}. Theorems 1.1--1.2; Section 4, ``An Improved Algorithm''; Section 6 for the lower bound.
\item \href{https://arxiv.org/pdf/1105.4024}{Belovs, \emph{Span Programs for Functions with Constant-Sized 1-Certificates} (2012)}. The graph-collision application.
\end{algorithmindex}


\entrytitle{39}{Testing bipartiteness in bounded-degree graphs}{strengthPolynomial}{evidenceProved}{hardwarePolynomial}
\setcounter{problem}{38}
\begin{problem}[Testing bipartiteness in bounded-degree graphs]
For a graph of maximum degree $d=O(1)$, distinguish a bipartite graph from one requiring more than $\varepsilon dn$ edge changes to become bipartite. Neighbor-list queries give the $j$th neighbor of a specified vertex. The promise permits sublinear property testing.
Complexity counts neighbor-oracle queries and running time for fixed maximum degree and fixed distance parameter.
\end{problem}

\speedupheading
\complexity{Quantum}{$\widetilde O(n^{1/3})$, Ambainis--Childs--Liu bipartiteness testing.}
The theorem assumes fixed $d$ and $\varepsilon$. Reversible random-walk generation and quantum collision finding locate opposite-parity walks with a common endpoint. The decision has bounded error, with one-sided acceptance for bipartite graphs.

\complexity{Classical}{$\widetilde\Theta(\sqrt n)$, Goldreich--Ron random-walk testing.}
The tester samples walks from chosen starting vertices and detects an odd-cycle certificate from conflicting parities. The lower bound applies to the same bounded-degree property-testing promise. Dependence on $d$, $\varepsilon$ and confidence is restored when these parameters vary with $n$.

\dimensionlabel{Strength}{strengthPolynomial} Polynomial: $\sqrt n\to\widetilde O(n^{1/3})$, through collision-finding quantum walks.

\dimensionlabel{Evidence}{evidenceProved} Rigorous advantage over a classical property-testing lower bound.

\dimensionlabel{Hardware}{hardwarePolynomial} Fault-tolerant neighbor queries and reversible walk generation are needed. A direct collision-finding implementation stores $\widetilde O(n^{1/3})$ bits of walk seeds and endpoints in the fixed-parameter regime, in addition to oracle workspace.

\subhead{Algorithm}
\begin{algorithmindex}{39}
\item \href{https://arxiv.org/pdf/1012.3174}{Ambainis, Childs and Liu, \emph{Quantum Property Testing for Bounded-Degree Graphs} (2011)}. Algorithm 2 and Theorem 10; Algorithm 1 and Theorem 7 give the classical comparator.
\item \href{https://doi.org/10.1007/s004930050059}{Goldreich and Ron, \emph{A Sublinear Bipartiteness Tester for Bounded Degree Graphs} (1999)}. Random-walk bipartiteness tester and lower-bound construction.
\end{algorithmindex}


\entrytitle{40}{Testing graph expansion}{strengthPolynomial}{evidenceProved}{hardwarePolynomial}
\setcounter{problem}{39}
\begin{problem}[Testing graph expansion]
For a bounded-degree graph, distinguish expansion at least $\Phi$ from being $\varepsilon$-far from every graph with expansion at least a specified smaller threshold $\Phi'$. The gap between $\Phi$ and $\Phi'$ is part of the testing promise. Expansion controls mixing and the robustness of cuts.
Complexity counts neighbor queries and quantum or classical operations implementing the prescribed random-walk tester.
\end{problem}

\speedupheading
\complexity{Quantum}{$\widetilde O(n^{1/3}/\Phi)$, Apers\textquotesingle s fast-forwarding and seed-set tester.}
The proximity parameter and degree are fixed. The promise relates the accepting expansion threshold $\Phi$ to the smaller rejecting threshold as specified in the tester theorem. The bound combines an evolving-set construction with quantum fast-forwarding and estimates of the resulting walk distribution.

\complexity{Classical}{$\widetilde O(\sqrt n/\Phi^2)$, classical random-walk expansion testing; $\Omega(\sqrt n)$ queries at fixed parameters.}
The classical algorithm checks whether sampled walk endpoints exhibit the concentration expected in poorly expanding regions. Upper and lower bounds refer to bounded-degree oracle access with the stated gap between the two expansion thresholds.

\dimensionlabel{Strength}{strengthPolynomial} Polynomial in $n$, and quadratic improvement in the walk-time dependence $\Phi^{-2}\to\Phi^{-1}$. The mechanism is quantum fast-forwarding and amplitude estimation.

\dimensionlabel{Evidence}{evidenceProved} Rigorous oracle advantage for the testing promise.

\dimensionlabel{Hardware}{hardwarePolynomial} Fault-tolerant coherent walk steps require logarithmic vertex registers, seed-set storage and estimation ancillas. The seed-set data structure supports coherent access; total depth depends on $\Phi^{-1}$ and the target confidence.

\subhead{Algorithm}
\begin{algorithmindex}{40}
\item \href{https://arxiv.org/pdf/1907.02369}{Apers, \emph{Expansion Testing Using Quantum Fast-Forwarding and Seed Sets} (2019)}. Algorithm 2 and Theorem 5; Theorem 1 states the main complexity bound.
\end{algorithmindex}


\entrytitle{41}{AND--OR formula and game-tree evaluation}{strengthPolynomial}{evidenceProved}{hardwarePolynomial}
\setcounter{problem}{40}
\begin{problem}[AND--OR formula and game-tree evaluation]
Given a fixed read-once formula $F$ with $N$ input leaves and gates AND, OR and NOT, query input bits and compute $F(x)$. For alternating AND--OR trees, this decides the winner of a finite perfect-information game with known tree and queried terminal outcomes.
Complexity counts queries to the leaf bits; gate complexity additionally includes access to the preprocessed formula graph.
\end{problem}

\speedupheading
\complexity{Quantum}{$O(\sqrt N)$ queries, quantum-walk and span-program formula evaluation.}
The bound applies to fixed read-once AND--OR formulas, allowing negations on inputs. Efficient implementations have nearly matching time after preprocessing the formula. The output is the Boolean value at the root.

\complexity{Classical}{$\Theta(N^c)$ for balanced binary trees, randomized alpha--beta pruning, where $c=\log_2((1+\sqrt{33})/4)\approx0.753$.}
The randomized algorithm evaluates children in random order and stops as soon as the parent value is determined. This is the optimal randomized classical bound for the balanced alternating binary tree. Arbitrary formulas have structure-dependent classical complexity and the general $O(N)$ leaf-reading upper bound.

\dimensionlabel{Strength}{strengthPolynomial} Polynomial: $N^{0.753\ldots}\to N^{1/2}$ on balanced binary trees. The mechanism is a quantum walk or span-program reflection algorithm.

\dimensionlabel{Evidence}{evidenceProved} Rigorous: the balanced-tree classical and quantum exponents have matching lower bounds.

\dimensionlabel{Hardware}{hardwarePolynomial} Fault-tolerant evaluation uses logarithmic graph-address registers, phase-estimation ancillas and coherent access to the preprocessed formula. Regular balanced trees admit succinct graph descriptions.

\subhead{Algorithm}
\begin{algorithmindex}{41}
\item \href{https://theoryofcomputing.org/articles/v008a013/}{Reichardt and \v Spalek, \emph{Span-Program-Based Quantum Algorithm for Evaluating Formulas} (2012)}. Formula-evaluation construction and the main algorithmic theorem.
\item \href{https://arxiv.org/pdf/1005.1601}{Reichardt, \emph{Reflections for Quantum Query Algorithms} (2011)}. Theorem 1.3 and Algorithms 1--3: conversion of span-program witnesses to query algorithms.
\end{algorithmindex}


\entrytitle{42}{Matrix-product verification over a field}{strengthPolynomial}{evidenceProved}{hardwarePolynomial}
\setcounter{problem}{41}
\begin{problem}[Matrix-product verification over a field]
Given $A,B,C\in\mathbb F^{n\times n}$ through entry oracles, decide whether $AB=C$, with field operations charged explicitly or treated as unit-cost operations. Matrix products encode path counts and algebraic certificates used throughout combinatorial algorithms.
The comparison counts matrix-entry queries and field operations; bit and gate costs also charge the representation and implementation of field arithmetic.
\end{problem}

\speedupheading
\complexity{Quantum}{$O(n^{5/3})$ queries and field-arithmetic time, Buhrman--\v Spalek quantum verification.}
The quantum walk combines selected-row and selected-column subproblems with randomized field checks. It accepts correct products and detects incorrect products with bounded error. The known quantum query lower bound is $\Omega(n^{3/2})$.

\complexity{Classical}{$O(n^2)$ field operations and $\Theta(n^2)$ entry queries, Freivalds verification.}
A random vector $r$ tests the identity $A(Br)=Cr$. Repetition reduces the error probability logarithmically. The quadratic entry-query bound is optimal. For $b$-bit field elements, both algorithms incur the corresponding classical or reversible arithmetic cost.

\dimensionlabel{Strength}{strengthPolynomial} Polynomial: $n^2\to n^{5/3}$, using quantum walks and randomized algebraic checks.

\dimensionlabel{Evidence}{evidenceProved} Rigorous separation from the tight classical entry-query bound.

\dimensionlabel{Hardware}{hardwarePolynomial} The published quantum algorithm uses $O(n)$ stored field elements, or $O(nb)$ logical data qubits for $b$-bit elements, plus arithmetic and oracle workspace. Fault-tolerant coherent arithmetic is required at useful scale.

\subhead{Algorithm}
\begin{algorithmindex}{42}
\item \href{https://arxiv.org/pdf/quant-ph/0409035}{Buhrman and \v Spalek, \emph{Quantum Verification of Matrix Products} (2006)}. Figure 1, ``Product Verification'' and ``Verify Once''; Theorem 6 gives time and space; Theorem 11 gives the lower bound.
\item \href{https://www.math.uwaterloo.ca/~anayak/papers/Nayak08.pdf}{Nayak, \emph{Checking Matrix Identities} (2008)}. The exposition of Freivalds\textquotesingle s identity test.
\end{algorithmindex}


\entrytitle{43}{Boolean matrix-product verification}{strengthPolynomial}{evidenceProved}{hardwarePolynomial}
\setcounter{problem}{42}
\begin{problem}[Boolean matrix-product verification]
For $A,B,C\in\{0,1\}^{n\times n}$, decide whether $C_{ij}=\bigvee_k(A_{ik}\wedge B_{kj})$ for all $i,j$. This checks the claimed two-step reachability relation of a directed graph.
The main comparison counts Boolean matrix-entry queries; explicit classical computation also counts arithmetic or word operations.
\end{problem}

\speedupheading
\complexity{Quantum}{$\widetilde O(n^{17/12})$ queries, Gilani--Le Gall--Zhou.}
The 2026 algorithm reduces verification to Orthogonal Vectors and uses quantum search on structured subinstances. Its quantum lower bound is $\Omega(n^{5/4})$. Nested quantum search also supplies an $O(n^{3/2})$-query construction with a simpler implementation.

\complexity{Classical}{$\Theta(n^2)$ entry queries; $O(n^{9/4+\varepsilon})$ arithmetic time for computing and checking the Boolean product.}
The query upper bound reads all three matrices. Exact Boolean multiplication can use algebraic matrix multiplication over a suitable field, with exponent $\omega\le9/4$, for every fixed $\varepsilon>0$. The problem uses Boolean-semiring multiplication; a zero output entry and a one output entry have different witness conditions.

\dimensionlabel{Strength}{strengthPolynomial} Polynomial query improvement: $n^2\to\widetilde O(n^{17/12})$. The method combines reductions, quantum search and quantum walks.

\dimensionlabel{Evidence}{evidenceProved} Rigorous separation from classical query complexity; the strongest cited bounds are a 2026 preprint.

\dimensionlabel{Hardware}{hardwarePolynomial} Fault-tolerant implementation needs coherent access to the matrices and to the records used by the Orthogonal Vectors procedure. Logical memory depends on that implementation; the simpler nested-search algorithm uses small index registers with reversible query workspace.

\subhead{Algorithm}
\begin{algorithmindex}{43}
\item \href{https://arxiv.org/pdf/2609.40180}{Gilani, Le Gall and Zhou, \emph{Improved Quantum Query Bounds for Boolean Matrix Product Verification} (2026)}. Theorems 1.1--1.2 and 3.5; Algorithms 2 and 6 and Theorems 4.1 and 4.6.
\item \href{https://arxiv.org/pdf/quant-ph/0409035}{Buhrman and \v Spalek, \emph{Quantum Verification of Matrix Products} (2006)}. Theorem 13: Boolean verification by nested search.
\end{algorithmindex}


\entrytitle{44}{Promised short-path \texorpdfstring{$s$--$t$}{s--t} connectivity}{strengthPolynomial}{evidenceProved}{hardwarePolynomial}
\setcounter{problem}{43}
\begin{problem}[Promised short-path $s$--$t$ connectivity]
Given an undirected graph and vertices $s,t$, distinguish the existence of an $s$--$t$ path of length at most $d$ from complete disconnection. The positive-instance path-length promise exposes how certificate structure affects quantum complexity.
The comparison counts adjacency-matrix queries and the quantum working space used between queries.
\end{problem}

\speedupheading
\complexity{Quantum}{$O(n\sqrt d)$ queries and $\widetilde O(n\sqrt d)$ time, Belovs--Reichardt span-program connectivity.}
The promise is that connected terminals have a path of at most $d$ edges. The algorithm uses logarithmic quantum space. The case $d=\Theta(n)$ gives $O(n^{3/2})$ queries. For very small $d$, direct specialized tests can improve the displayed general upper bound; for $d=1$ a single adjacency query suffices.

\complexity{Classical}{$O(n^2)$ matrix time and queries, breadth-first search; $O(n+m)$ with explicit adjacency lists.}
The unrestricted matrix-oracle problem has classical complexity $\Theta(n^2)$. The promise parameter must be retained when comparing algorithms on shorter-path subclasses. A classical traversal records discovered vertices and stops once $t$ is reached or its component has been exhausted.

\dimensionlabel{Strength}{strengthPolynomial} Polynomial oracle speedup, including $n^2\to n^{3/2}$ for the unrestricted path-length regime. The mechanism is a span program and quantum reflections.

\dimensionlabel{Evidence}{evidenceProved} Rigorous query advantage in the regime where the classical quadratic lower bound applies.

\dimensionlabel{Hardware}{hardwarePolynomial} The algorithm uses $O(\log n)$ logical working qubits up to precision ancillas, with coherent adjacency access. This space-efficient construction still needs fault tolerance for its long sequence of queries.

\subhead{Algorithm}
\begin{algorithmindex}{44}
\item \href{https://arxiv.org/pdf/1203.2603}{Belovs and Reichardt, \emph{Span Programs and Quantum Algorithms for $st$-Connectivity and Claw Detection} (2012)}. Theorem 3: query bound; Theorem 9: time-efficient implementation.
\end{algorithmindex}


\entrytitle{46}{Global minimum cut}{strengthPolynomial}{evidenceProved}{hardwarePolynomial}
\setcounter{problem}{45}
\begin{problem}[Global minimum cut]
In a weighted undirected graph, output a nonempty proper subset $S\subset V$ minimizing $w(\delta(S))=\sum_{u\in S,v\notin S}w_{uv}$. This identifies the cheapest edge failure that disconnects a network.
The quantum comparison counts weighted adjacency queries and running time; classical bounds refer to explicit weighted graphs.
\end{problem}

\speedupheading
\complexity{Quantum}{$\widetilde O(n^{3/2}\sqrt\tau)$ matrix time and queries, quantum minimum cut.}
Here $\tau$ is the ratio of the largest to smallest positive weight; reading the matrix caps the bound at $O(n^2)$. In the adjacency-array model, the query bound is $\widetilde O(\sqrt{mn\tau})$ and the time bound is $\widetilde O(\sqrt{mn\tau}+n^{3/2})$. The algorithm returns the minimum-cut weight and shores. Approximate cuts can instead use quantum cut sparsification.

\complexity{Classical}{$\widetilde O(m)$ expected time, randomized classical minimum cut; $\Theta(n^2)$ matrix queries.}
The explicit-input comparison includes randomized contraction and near-linear cut algorithms. The matrix quantum lower bound is $\Omega(n^{3/2}\sqrt\tau)$ in the nontrivial range $1\le\tau\le n$, up to the precision and logarithmic conventions of the theorem.

\dimensionlabel{Strength}{strengthPolynomial} Polynomial oracle improvement; for bounded weight ratio and dense graphs it is $n^2\to\widetilde O(n^{3/2})$.

\dimensionlabel{Evidence}{evidenceProved} Rigorous matrix-query separation with nearly matching quantum bounds.

\dimensionlabel{Hardware}{hardwarePolynomial} Fault-tolerant search and sparsification require logarithmic vertex registers, weight registers and coherent graph-access workspace. A classical sparse graph and contracted partitions are stored during the computation.

\subhead{Algorithm}
\begin{algorithmindex}{46}
\item \href{https://arxiv.org/pdf/2011.09823}{Apers and Lee, \emph{Quantum Complexity of Minimum Cut} (2021)}. Theorems 3 and 5 for the matrix and adjacency-array algorithms; the named routines FindMinStar, Cut-Sparsifier, LearnCutAtoms and LearnContraction.
\end{algorithmindex}


\entrytitle{49}{MaxCut approximation with QAOA}{white}{white}{hardwareSmall}
\setcounter{problem}{48}
\begin{problem}[Maximum cut]
Given nonnegative weights $w_{uv}$, maximize $C(z)=\sum_{\{u,v\}\in E}w_{uv}(1-z_uz_v)/2$ over $z\in\{-1,1\}^n$, or return a cut meeting a specified approximation ratio. Sparse regular graphs are a central QAOA benchmark.
Complexity uses quantum gate and sample counts, classical arithmetic time for approximation, and classical bit complexity for exact finite-precision instances.
\end{problem}

\speedupheading
\complexity{Quantum}{$O(p(m+n))$ gates per sample, depth-$p$ QAOA; $O(KSp(m+n))$ for $K$ objective evaluations with $S$ shots each.}
The cost includes the problem-phase and mixer layers; synthesis precision contributes additional factors. On 3-regular graphs, the $p=1$ analysis guarantees an approximation ratio of at least $0.6924$. Results on higher-girth graphs and larger depth have their own graph and parameter assumptions. The complete variational optimization includes the evaluations and samples represented by $K$ and $S$.

\complexity{Classical}{Polynomial time with ratio $0.87856$, Goemans--Williamson; $\operatorname{poly}(n)2^{O(n/\sqrt{\log n})}$ for exact sparse unweighted MaxCut.}
The approximation guarantee applies to general nonnegative weights. For unweighted graphs with $m=O(n)$, threshold counting and self-reduction give the exact subexponential bound. For arbitrary unweighted graphs and fixed $\varepsilon>0$, sampling $O_\varepsilon(n)$ edges and optimizing the sample gives value at least $\operatorname{OPT}-\varepsilon m$ with the same exponent. For constant $p$, bounded degree and a fixed finite algebraic gate set, the complete QAOA output distribution can be sampled classically in expected $\operatorname{poly}(n)2^{O(n/\sqrt{\log n})}$ time; compiled gate count controls the corresponding statement for synthesized rotations.

\dimensionlabel{Strength}{white} The general QAOA runtime and approximation theorems establish no asymptotic advantage over the applicable classical comparator.

\dimensionlabel{Evidence}{white} Approximation guarantees are rigorous for specified graph classes and depths. Empirical comparisons are tied to their tested instances and classical solvers.

\dimensionlabel{Hardware}{hardwareSmall} Small superconducting-processor demonstrations exist. A circuit uses $n$ data qubits plus ancillas; useful large-depth operation needs accurate two-qubit gates, repeated sampling and parameter optimization. Fault-tolerant estimates also include rotation synthesis.

\subhead{Algorithm}
\begin{algorithmindex}{49}
\item \href{https://arxiv.org/pdf/1411.4028}{Farhi, Goldstone and Gutmann, \emph{A Quantum Approximate Optimization Algorithm} (2014)}. Section II, ``Fixed $p$ Algorithm''; Section V, ``MaxCut on 3-Regular Graphs''.
\item \href{https://math.mit.edu/~goemans/PAPERS/maxcut-jacm.pdf}{Goemans and Williamson, \emph{Improved Approximation Algorithms for Maximum Cut and Satisfiability Problems Using Semidefinite Programming} (1995)}. The MaxCut semidefinite relaxation and randomized hyperplane rounding.
\item \href{https://github.com/openai/math/blob/main/preprints/The-Exponential-Time-Hypothesis-is-False-October-25-2026/paper.pdf}{OpenAI, \emph{The Exponential Time Hypothesis Is False} (25 October 2026)}. Theorem 1.1 and the sparse-instance bound; witness recovery by conditional counting.
\item \href{https://drive.google.com/file/d/1RzhP9ZMy1xkC22oqKmMHEOh7gG1cwkfA/view}{Gheorghiu and Gupte, \emph{Subexponential-Time Sampling of Quantum Circuits from Exact Counting} (2026)}. Main exact-sampling theorem for fixed finite algebraic gate sets.
\item \href{https://arxiv.org/pdf/2004.04197}{Harrigan et al., \emph{Quantum Approximate Optimization of Non-Planar Graph Problems on a Planar Superconducting Processor} (2021)}. Experimental QAOA circuits and performance figures.
\end{algorithmindex}


\entrytitle{50}{Bounded-occurrence Max-E3LIN2}{white}{white}{hardwarePolynomial}
\setcounter{problem}{49}
\begin{problem}[Bounded-occurrence Max-E3LIN2]
Given $m$ equations $x_a\oplus x_b\oplus x_c=b_j$ over $\mathbb F_2$, with every variable appearing in at most $D$ equations, maximize the number satisfied. Random assignment satisfies half in expectation; the approximation question measures the additive improvement over $1/2$.
Complexity counts quantum gates for QAOA, classical arithmetic time for approximation, and bit operations for exact bounded-degree instances.
\end{problem}

\speedupheading
\complexity{Quantum}{$O(m+n)$ gates per sample, $p=1$ QAOA for bounded-occurrence Max-E3LIN2.}
For sufficiently large $D$, the guaranteed satisfied fraction is at least $1/2+1/(101\sqrt D\ln D)$. The theorem chooses parameters from the objective analysis; shot counts for estimating the objective and success amplification contribute additional costs.

\complexity{Classical}{Polynomial time with fraction $1/2+\Omega(D^{-1/2})$, Barak et al.; $\operatorname{poly}(n)2^{O_D(n/\sqrt{\log n})}$ for exact optimization at fixed $D$.}
The classical approximation algorithm applies to odd-arity XOR constraints and improves the dependence on $D$ in the QAOA guarantee. At fixed occurrence bound, $m=O(n)$ and the exact objective threshold has a linear-size Boolean circuit. Exact counting, binary search and witness recovery supply the displayed bit-complexity bound.

\dimensionlabel{Strength}{white} The approximation guarantee has a polynomial-time classical comparator with a stronger dependence on $D$. The stated quantum algorithm has no established runtime advantage.

\dimensionlabel{Evidence}{white} The approximation theorems are rigorous; they yield no asymptotic quantum-speedup claim in this comparison.

\dimensionlabel{Hardware}{hardwarePolynomial} A QAOA circuit uses $n$ data qubits and coherent three-body parity phases, typically decomposed into two-qubit gates. Modest instances are feasible on circuit hardware; scaling depends on connectivity, noise and sampling cost.

\subhead{Algorithm}
\begin{algorithmindex}{50}
\item \href{https://arxiv.org/pdf/1412.6062}{Farhi, Goldstone and Gutmann, \emph{A Quantum Approximate Optimization Algorithm Applied to a Bounded Occurrence Constraint Problem} (2014)}. Section II, ``The General Case'', for the worst-case guarantee; Section III for typical instances.
\item \href{https://arxiv.org/pdf/1505.03424}{Barak et al., \emph{Beating the Random Assignment on Constraint Satisfaction Problems of Bounded Degree} (2015)}. Theorem 1.1: odd-arity Max-$k$XOR approximation.
\item \href{https://github.com/openai/math/blob/main/preprints/The-Exponential-Time-Hypothesis-is-False-October-25-2026/paper.pdf}{OpenAI, \emph{The Exponential Time Hypothesis Is False} (25 October 2026)}. Theorem 1.1 and the sparse-instance bound; witness recovery by conditional counting.
\end{algorithmindex}


\entrytitle{51}{Maximum independent set in Rydberg geometric graphs}{white}{evidencePartial}{hardwareSmall}
\setcounter{problem}{50}
\begin{problem}[Maximum independent set in unit-disk graphs]
Given points $r_1,\ldots,r_n\in\mathbb R^2$ and a blockade radius $R$, let $\{i,j\}\in E$ when $\|r_i-r_j\|\le R$. Find the largest $S\subseteq[n]$ with $\|r_i-r_j\|>R$ for distinct $i,j\in S$. This geometric promise makes independent-set constraints directly implementable with atom interactions.
The experiment measures analog evolution time and repetitions to success; the exact classical comparison uses time on explicitly represented geometric graphs.
\end{problem}

\speedupheading
\complexity{Quantum}{$t\lceil\log\delta/\log(1-p)\rceil$ evolution time to success probability $1-\delta$, Rydberg adiabatic and variational optimization.}
Here $p$ is the measured single-run success probability and $t$ the evolution duration. State preparation, readout and parameter tuning add to the total time. Experiments with up to 289 atoms reported favorable scaling against an optimized simulated-annealing comparator on the tested graph families.

\complexity{Classical}{$2^{O(\sqrt n)}\operatorname{poly}(n)$ for exact unit-disk independent set, geometric separators and dynamic programming.}
The geometric promise is essential: weighted separators and tree decompositions for similarly sized fat objects yield the subexponential exact algorithm in the plane. A polynomial-time approximation scheme is available for unit-disk graphs at fixed approximation error. Simulated annealing is the experimental comparator, while the geometric separator algorithm supplies the exact worst-case comparison.

\dimensionlabel{Strength}{white} Experimental improvement over a specified annealing solver; the cited results establish no general asymptotic advantage over exact geometric algorithms.

\dimensionlabel{Evidence}{evidencePartial} Partial evidence from measured time-to-solution scaling on the tested instances and solver settings.

\dimensionlabel{Hardware}{hardwareSmall} The 289-atom demonstration uses physical analog degrees of freedom. One atom represents one vertex; blockade geometry, calibration, coherence time and preparation/readout govern scaling.

\subhead{Algorithm}
\begin{algorithmindex}{51}
\item \href{https://arxiv.org/pdf/2202.09372}{Ebadi et al., \emph{Quantum Optimization of Maximum Independent Set Using Rydberg Atom Arrays} (Science 2022)}. Main-text Figures 3--4; Supplementary Sections 4--5 define QAOA and variational adiabatic protocols; Sections 8--9 give the annealing comparison and success scaling.
\item \href{https://arxiv.org/pdf/1803.10633}{de Berg et al., \emph{A Framework for ETH-Tight Algorithms and Lower Bounds in Geometric Intersection Graphs} (2020)}. Theorem 14: Independent Set on a weighted tree decomposition; the geometric separator and decomposition theorems specialize to $d=2$.
\end{algorithmindex}


\entrytitle{52}{Sherrington--Kirkpatrick spin-glass optimization}{white}{white}{hardwareSmall}
\setcounter{problem}{51}
\begin{problem}[Sherrington--Kirkpatrick spin-glass optimization]
For independent standard Gaussian couplings $J_{ij}$, minimize $H(z)=n^{-1/2}\sum_{i<j}J_{ij}z_i z_j$ over $z\in\{-1,1\}^n$. One may request the exact ground state or an energy within $\varepsilon n$ of it. Random fully connected spin glasses are an average-case optimization benchmark.
Quantum circuits are measured in gates and samples; the Gaussian classical approximation theorem counts real-arithmetic operations, while exact discretized instances use bit complexity.
\end{problem}

\speedupheading
\complexity{Quantum}{$O(pn^2)$ gates per sample, depth-$p$ QAOA for the SK Hamiltonian.}
The infinite-size expected-energy analysis can be evaluated classically in $O(16^p)$ time for fixed parameters. It predicts the achieved energy density and guides angle selection. Exact optimization by generalized short-path algorithms uses additional overlap, spectral and landscape hypotheses, with runtime relative to an associated classical Markov-chain search.

\complexity{Classical}{$C(\varepsilon)n^2$ arithmetic operations for an $\varepsilon n$-near optimum, Montanari\textquotesingle s incremental approximate message passing.}
The guarantee assumes full replica symmetry breaking and holds with high probability over Gaussian couplings. For exact optimization, enumeration supplies $O(n^2 2^n)$ arithmetic time, with instance-specific branch-and-bound improvements. Discretizing dense couplings to $b$ bits gives an energy-threshold circuit of size $O(n^2b)$; sparse exact counting on that circuit costs $\operatorname{poly}(n,b)2^{O(n^2b/\sqrt{\log(n^2b)})}$. The smaller exact bound is chosen for the represented instance.

\dimensionlabel{Strength}{white} The QAOA energy-density results have a polynomial-time classical near-optimal comparator. A general asymptotic QAOA advantage is unestablished.

\dimensionlabel{Evidence}{white} The limiting QAOA formulas and conditional classical approximation theorem are analytical results; empirical comparisons concern fixed depths and finite instances.

\dimensionlabel{Hardware}{hardwareSmall} QAOA uses $n$ data qubits but $\Theta(n^2)$ interactions per layer before compilation. Small processor experiments exist; dense routing, two-qubit errors and repeated shots determine practical cost.

\subhead{Algorithm}
\begin{algorithmindex}{52}
\item \href{https://arxiv.org/pdf/1910.08187}{Farhi et al., \emph{The Quantum Approximate Optimization Algorithm and the Sherrington--Kirkpatrick Model at Infinite Size} (2022)}. Section 2 defines QAOA; Section 4 states the limiting-energy result; Section 6 gives the general-$p$ calculation.
\item \href{https://arxiv.org/pdf/1812.10897}{Montanari, \emph{Optimization of the Sherrington--Kirkpatrick Hamiltonian} (2021)}. Assumption 1, Theorem 2, and Algorithm 1, ``IAMP Algorithm to Optimize SK Hamiltonian''.
\item \href{https://arxiv.org/pdf/2410.23270}{Chakrabarti et al., \emph{Generalized Short Path Algorithms} (2024/2025)}. Algorithms 1--2 and Theorem 3.5.
\item \href{https://github.com/openai/math/blob/main/preprints/The-Exponential-Time-Hypothesis-is-False-October-25-2026/paper.pdf}{OpenAI, \emph{The Exponential Time Hypothesis Is False} (25 October 2026)}. Theorem 1.1 and the sparse-instance bound; witness recovery by conditional counting.
\end{algorithmindex}


\entrytitle{53}{Exact bounded-locality spin optimization by short paths}{white}{white}{hardwareExponential}
\setcounter{problem}{52}
\begin{problem}[Exact fixed-arity spin optimization]
Given $H(z)=\sum_{S:|S|\le k}J_S\prod_{i\in S}z_i$, output $z^*\in\arg\min H$; $k$ is fixed and coefficients are efficiently represented. QUBO is the case $k=2$ after changing binary variables to spins.
Complexity compares quantum gate time with classical bit time for explicitly represented fixed-arity objective functions.
\end{problem}

\speedupheading
\complexity{Quantum}{$2^{(1-c_{\rm cl})n/2+o(n)}$, Le Gall--Tamaki quantum conditioning-and-search under the low-energy-density promise.}
For a centered homogeneous degree-$k$ objective, let $H_{\min}=\min_z H(z)<0$ and suppose
\[
\bigl|\{z:H(z)\le(1-\eta)H_{\min}\}\bigr|\le2^{(1-\gamma)n}.
\]
For fixed $k$, $0<\eta<1$ and $\gamma>0$, the exponent is
\[
c_{\rm cl}=\min\!\left\{\gamma,\ h\!\left(\frac{1-(1-\eta)^{1/k}}{2}\right)\right\},
\]
where $h$ is binary entropy. The algorithm combines amplitude amplification with minimum finding; the weighted formulation removes the known-optimum assumption with a $2^{o(n)}$ overhead. Dalzell et al.'s short-path bound under the corresponding promise has exponent $(1-c_{\rm q})n/2$, where $c_{\rm q}=\gamma\eta/[2(2+\ln2)k]$, for $k\le n/6$ and $\gamma\ge(1+4\log_2n)/n$. The generalized short-path framework gives $O^*(T(n)^{1/2-c(n)})$ for its specified Markov-chain-search cost $T(n)$ under its overlap and spectral hypotheses. The subexponential factors in the conditioning-and-search bounds are displayed explicitly.

\complexity{Classical}{$2^{(1-c_{\rm cl})n+o(n)}$ under the same density promise, Le Gall--Tamaki; $\operatorname{poly}(N)2^{O_k(N/\sqrt{\log N})}$ with $N=O_k(n+mb)$, exact threshold counting.}
The first bound uses the same centered homogeneous objective and parameters as the quantum conditioning-and-search algorithm. For an objective with $m$ integer terms and coefficients of at most $b$ bits, a threshold circuit and its Boolean encoding have $O_k(n+mb)$ variables and gates. Binary search and conditional counting recover an optimum at the displayed counting cost. The classical comparison takes the minimum of the applicable bounds, $\operatorname{poly}(n,m,b)2^n$ enumeration and specialized exact solvers. The density promise, term count and coefficient precision determine which bound is smallest.

\dimensionlabel{Strength}{white} The paired conditioning-and-search bounds have a quadratic time relationship up to subexponential factors. The overall comparison uses the smallest applicable classical bound. Generalized short paths give a stronger-than-quadratic improvement over their specified Markov-chain-search cost $T(n)$.

\dimensionlabel{Evidence}{white} The low-energy-density, short-path and exact-counting theorems give rigorous parameter-dependent upper bounds. The paired algorithmic comparison has the same objective and density promise on both sides.

\dimensionlabel{Hardware}{hardwareExponential} Fault-tolerant implementation needs $n$ assignment qubits, polynomial reversible workspace and exponentially many coherent operations for the displayed fixed-positive-exponent bound. Cost precision and spectral resolution contribute to depth.

{\small [Before ``ETH is false'', the cited CSP comparisons had classical bounds $2^{(1-c\prime)n}$ and quantum bounds $2^{(1-c)n/2}$ for problem-dependent constants, and claimed an improvement in the exponential constant.]\par}

\subhead{Algorithm}
\begin{algorithmindex}{53}
\item \href{https://arxiv.org/pdf/2212.01513}{Dalzell et al., \emph{Mind the Gap: Achieving a Super-Grover Quantum Speedup by Jumping to the End} (2023/2026)}. Algorithm 1 and Theorems 1, 4, 6--8.
\item \href{https://arxiv.org/pdf/2410.23270}{Chakrabarti et al., \emph{Generalized Short Path Algorithms} (2024/2025)}. Algorithm 2, Theorems 3.3--3.5, and Table 1.
\item \href{https://arxiv.org/pdf/2604.12131}{Le Gall and Tamaki, \emph{Dequantizing Short-Path Quantum Algorithms} (2026)}. Theorems 1.1--1.4 and 4.1: matched classical and quantum bounds; Section 4.3: amplitude amplification and minimum finding; Appendices A--B: removing the known-optimum assumption.
\item \href{https://github.com/openai/math/blob/main/preprints/The-Exponential-Time-Hypothesis-is-False-October-25-2026/paper.pdf}{OpenAI, \emph{The Exponential Time Hypothesis Is False} (25 October 2026)}. Theorem 1.1 and the sparse-instance bound; witness recovery by conditional counting.
\end{algorithmindex}


\entrytitle{54}{Partition functions and approximate combinatorial counting}{strengthPolynomial}{evidencePartial}{hardwarePolynomial}
\setcounter{problem}{53}
\begin{problem}[Partition-function estimation]
Given an energy $H:\Omega\to\{0,\ldots,B\}$ and inverse temperature $\beta$, output $\widehat Z$ with $|\widehat Z/Z-1|\le\varepsilon$, where $Z(\beta)=\sum_{x\in\Omega}e^{-\beta H(x)}$. Suitable choices count independent sets, matchings and colorings; the zero-temperature limit counts optimal configurations.
Complexity counts coherent or classical Markov-chain steps and energy evaluations, with arithmetic costs restored for numerical output.
\end{problem}

\speedupheading
\complexity{Quantum}{$\widetilde O((\log|\Omega|)^{3/4}/(\varepsilon\sqrt\delta))$, Cornelissen--Hamoudi partition-function estimation.}
Here $\delta$ bounds the spectral gaps of the supplied sampling chains. Montanaro\textquotesingle s estimator uses $\widetilde O(\ell^2/(\varepsilon\sqrt\delta))$ for a supplied length-$\ell$ constant-relative-variance cooling schedule. In the independent-set uniqueness regime, appropriate Glauber dynamics yield a representative $\widetilde O(n^{5/4}/\varepsilon)$ bound. Schedule preparation and initial-state assumptions are part of these theorems.

\complexity{Classical}{$\widetilde O(\log|\Omega|/(\varepsilon^2\delta))$, \v Stefankovi\v c--Vempala--Vigoda adaptive cooling and classical sampling.}
For explicit sparse graphs, independent-set, fixed-$q$ coloring and perfect-matching counts also admit exact $\operatorname{poly}(n)2^{O(n/\sqrt{\log n})}$ time through sparse Boolean counting. For sparse constant-arity, constant-bit integer energies, counting each of the $O(n)$ possible energies gives the density of states and hence $Z$ with that exponent, plus numerical precision costs. These exact methods are useful when the supplied chain mixes slowly; in uniqueness regimes, polynomial-time correlation-decay and sampling algorithms provide classical alternatives.

\dimensionlabel{Strength}{strengthPolynomial} Polynomial improvement for the stated sampling model, including $\varepsilon^{-2}\to\varepsilon^{-1}$ and $\delta^{-1}\to\delta^{-1/2}$. Quantum amplitude estimation and quantum walks supply these factors.

\dimensionlabel{Evidence}{evidencePartial} Rigorous algorithmic comparison for the supplied chains and cooling assumptions. The explicit-instance comparator depends on the counting model and mixing regime.

\dimensionlabel{Hardware}{hardwarePolynomial} A configuration on $n$ bits needs at least $n$ logical data qubits, plus reversible transition, energy and estimation workspace. Fault tolerance and coherent Markov-chain steps are required; depth depends on $\delta^{-1/2}$.

\subhead{Algorithm}
\begin{algorithmindex}{54}
\item \href{https://arxiv.org/pdf/2207.08643}{Cornelissen and Hamoudi, \emph{A Sublinear-Time Quantum Algorithm for Approximating Partition Functions} (2023)}. Section 4 and Theorem 4.1; Algorithm 6 and Theorem 3.3 define the product estimator.
\item \href{https://pmc.ncbi.nlm.nih.gov/articles/PMC4614442/}{Montanaro, \emph{Quantum Speedup of Monte Carlo Methods} (2015)}. The partition-function approximation application and its cooling-schedule theorem.
\item \href{https://github.com/openai/math/blob/main/preprints/The-Exponential-Time-Hypothesis-is-False-October-25-2026/paper.pdf}{OpenAI, \emph{The Exponential Time Hypothesis Is False} (25 October 2026)}. Theorem 1.1 and the sparse-instance bound; witness recovery by conditional counting.
\end{algorithmindex}


\entrytitle{55}{Traversing a randomly welded pair of trees}{strengthPolynomialTime}{evidenceProved}{hardwarePolynomial}
\setcounter{problem}{54}
\begin{problem}[Welded-tree traversal]
Two depth-$h$ binary trees have their leaves joined by a random alternating cycle; vertex names are random strings. Given the entrance root and an oracle returning neighbors, output the other root. The graph has $N=\Theta(2^h)$ vertices and constant degree. This is a graph-search oracle problem exhibiting structural exponential advantage.
Complexity counts neighbor-oracle queries and quantum gates in the random-label oracle model.
\end{problem}

\speedupheading
\complexity{Quantum}{$\operatorname{poly}(h)$ queries and gates, the Childs et al. continuous-time quantum walk.}
The graph has $\Theta(2^h)$ vertices, but its column symmetry confines the relevant evolution to a polynomial-dimensional effective path. The oracle implements sparse-graph evolution from the entrance, and measurement finds the exit with probability amplified to a constant.

\complexity{Classical}{$2^{\Omega(h)}$ queries in the worst case, classical randomized oracle exploration.}
The original lower bound shows that an algorithm making at most $2^{h/6}$ queries has exponentially small success probability. Breadth-first exploration gives an $O(2^h)$ upper bound. Random labels and the hidden leaf welding prevent a classical query algorithm from reconstructing the useful column structure with polynomially many queries.

\dimensionlabel{Strength}{strengthPolynomialTime} Exponential to polynomial, through continuous-time quantum transport and interference.

\dimensionlabel{Evidence}{evidenceProved} Fully rigorous quantum--classical separation in the specified graph-oracle model.

\dimensionlabel{Hardware}{hardwarePolynomial} Fault-tolerant digital simulation uses $O(h)$ vertex-label qubits plus sparse-Hamiltonian simulation and oracle workspace. A physical analog realization with one site per graph vertex instead needs $\Theta(2^h)$ sites, a distinct resource model.

\subhead{Algorithm}
\begin{algorithmindex}{55}
\item \href{https://arxiv.org/pdf/quant-ph/0209131}{Childs et al., \emph{Exponential Algorithmic Speedup by a Quantum Walk} (2003)}. Theorem 3: quantum traversal; Theorem 9: classical query lower bound.
\end{algorithmindex}


\entrytitle{56}{Gibbs sampling and quantum simulated annealing}{strengthPolynomial}{evidencePartial}{hardwareExponential}
\setcounter{problem}{55}
\begin{problem}[Gibbs sampling]
Given $H:\Omega\to\mathbb R$ and $\beta$, generate a sample within total-variation error $\varepsilon$ of $\pi_\beta(x)=e^{-\beta H(x)}/Z(\beta)$. For an integer-valued objective with unit gap, choosing $\beta\ge\log(|\Omega|/\eta)$ makes an optimum occur with probability at least $1-\eta$. Thus sampling can yield discrete optimization.
The chain-based comparison counts coherent or classical transition steps; explicit constraint-model sampling also has a classical bit-time bound.
\end{problem}

\speedupheading
\complexity{Quantum}{$\widetilde O(L/\sqrt\delta)$ coherent walk steps, quantum simulated annealing and quantum sampling.}
The theorem assumes a supplied sequence of $L$ reversible chains with gap at least $\delta$, an efficiently prepared initial distribution and suitable overlap between adjacent stationary states. Accuracy and overlap parameters contribute logarithmic or explicit factors. Measuring the final coherent Gibbs state produces a sample.

\complexity{Classical}{$\widetilde O(L/\delta)$ transition steps, sequential classical mixing; subexponential conditional counting for sparse explicit models.}
Warm starts and the smallest relevant stationary probability enter the mixing bound. Sparse fixed-width hard constraints admit exact uniform sampling in expected $\operatorname{poly}(n)2^{O(n/\sqrt{\log n})}$ time by conditional counts. For sparse constant-arity, constant-bit integer energies, conditional counts at the $O(n)$ possible energies yield Gibbs samples to total-variation error $\varepsilon$ with the same exponent and polynomial numerical overhead. The smaller applicable sampler determines the explicit-instance comparison.

\dimensionlabel{Strength}{strengthPolynomial} Quadratic improvement in inverse spectral gap for the specified chain sequence, from quantum-walk phase detection and state transformations. An exponentially small gap gives an improvement in the exponential constant of that chain-based bound.

\dimensionlabel{Evidence}{evidencePartial} Rigorous relative-to-chain algorithmic speedup. Explicit-instance comparisons also include samplers obtained from counting.

\dimensionlabel{Hardware}{hardwareExponential} A configuration requires $\log_2|\Omega|$ logical data qubits plus reversible transition, energy and phase-estimation workspace. Fault-tolerant depth scales with $L/\sqrt\delta$ and can remain exponential for slowly mixing chains.

\subhead{Algorithm}
\begin{algorithmindex}{56}
\item \href{https://arxiv.org/pdf/0804.1571}{Somma, Boixo, Barnum and Knill, \emph{Quantum Simulations of Classical Annealing Processes} (2008)}. The quantum simulated annealing construction and its spectral-gap runtime analysis.
\item \href{https://arxiv.org/pdf/0804.4259}{Wocjan and Abeyesinghe, \emph{Speedup via Quantum Sampling} (2008)}. Theorem 2: preparation along a sequence of reversible Markov chains.
\item \href{https://github.com/openai/math/blob/main/preprints/The-Exponential-Time-Hypothesis-is-False-October-25-2026/paper.pdf}{OpenAI, \emph{The Exponential Time Hypothesis Is False} (25 October 2026)}. Theorem 1.1 and the sparse-instance bound; witness recovery by conditional counting.
\end{algorithmindex}


\entrytitle{59}{All-pairs shortest paths}{strengthPolynomial}{evidencePartial}{hardwarePolynomial}
\setcounter{problem}{58}
\begin{problem}[All-pairs shortest paths]
Given a directed graph with real edge weights and no negative cycle, output $D_{uv}$, the minimum path length from $u$ to $v$, for every ordered pair. This full $n\times n$ distance table is central to routing and fine-grained complexity.
The general comparison uses quantum arithmetic time with coherent matrix access and classical real-RAM operations; bounded-integer and unweighted variants use their stated word-RAM models.
\end{problem}

\speedupheading
\complexity{Quantum}{$\widetilde O(n^{5/2})$, quantum minimum finding for min-plus products and repeated squaring.}
Each distance-product entry is a minimum over $n$ candidates, and the complete distance table is output classically. Arithmetic on real weights is treated as unit cost in this comparison; a finite-precision circuit charges the corresponding addition, comparison and access gates.

\complexity{Classical}{$O(n^{2.998})$ expected time, Alman--Vassilevska Williams Las Vegas APSP.}
The real-weight algorithm uses comparisons, additions and subtractions and also has a corresponding high-probability runtime guarantee. Polynomially bounded integer weights admit deterministic $O(n^{2.99942})$ word-RAM time. Directed unweighted APSP admits $n^{2.4999+o(1)}$ time via rectangular matrix multiplication. Each specialized comparator applies to its stated weight regime.

\dimensionlabel{Strength}{strengthPolynomial} Polynomial for arbitrary real weights: approximately a factor $n^{0.498}$ between the displayed upper bounds, with quantum logarithmic factors suppressed. The mechanism is minimum finding.

\dimensionlabel{Evidence}{evidencePartial} Partial evidence from the gap between known general real-weight algorithms. The unweighted specialization has a stronger classical bound.

\dimensionlabel{Hardware}{hardwarePolynomial} Fault-tolerant implementation requires coherent access to $n^2$ distance entries and reversible arithmetic. With $b$-bit distances, the direct stored matrix occupies $O(n^2b)$ bits accessible to quantum queries, plus active index and arithmetic registers.

\subhead{Algorithm}
\begin{algorithmindex}{59}
\item \href{https://arxiv.org/pdf/1410.6220}{Nayebi and Vassilevska Williams, \emph{Quantum Algorithms for Shortest Paths Problems in Structured Instances} (2014)}. Introduction: general $\widetilde O(n^{5/2})$ APSP baseline; structured algorithms in Theorems 1--5.
\item \href{https://arxiv.org/pdf/2610.06783}{Alman and Vassilevska Williams, \emph{Truly Subquadratic 3SUM and Truly Subcubic APSP via Triangles in Sparse Lopsided Graphs} (2026)}. Theorems 22 and 35: integer and real-weight APSP.
\item \href{https://arxiv.org/pdf/2610.10494v2}{Uzna\'nski, \emph{Rectangular Matrix Multiplication from Polynomial Products} (2026)}. Section 7: directed unweighted APSP.
\end{algorithmindex}


\entrytitle{60}{Optimal polynomial intersection by decoded quantum interferometry}{strengthPolynomialTime}{evidencePartial}{hardwarePolynomial}
\setcounter{problem}{59}
\begin{problem}[Optimal polynomial intersection]
Given prime $p$, integer $n<p-1$, and explicit sets $F_y\subseteq\mathbb F_p$ for $y\ne0$, find a polynomial $Q\in\mathbb F_p[y]$, $\deg Q<n$, maximizing $|\{y:Q(y)\in F_y\}|$. The approximation task asks for a specified fraction of the $p-1$ constraints. This connects finite-field fitting, list recovery and coding theory.
The explicit-input comparison counts quantum gates and classical finite-field operations; the folded-OPI separation counts membership-oracle queries.
\end{problem}

\speedupheading
\complexity{Quantum}{$\operatorname{poly}(p)$ gates, decoded quantum interferometry for the stated balanced OPI regime.}
With $|F_y|\approx p/2$ and constant $n/p$, DQI achieves fraction $1/2+\sqrt{(n/(2p))(1-n/(2p))}$ asymptotically, using efficient decoding of the relevant Reed--Solomon dual code. Horinaga--Yamakawa give a worst-case polynomial-time quantum algorithm with full satisfaction for balanced sets at rate above $3/4$. For the folded membership-oracle family at rate $0.3$, the separation gives expected quantum score $0.85$ while a classical score exceeding $0.65$ by a constant requires superpolynomially many queries.

\complexity{Classical}{Polynomial-time interpolation with expected fraction $1/2+n/(2p)$, Prange-type interpolation; exponential repetition to reach the studied DQI target.}
Complete enumeration costs $p^n\operatorname{poly}(p,n)$. The explicit balanced tables contain $\Theta(p^2)$ bits. A direct Boolean objective circuit has $O(pn\log^2p+p^2\log p)$ gates, so applying subexponential exact counting to this circuit gives a bound whose exponent, at constant $n/p$, exceeds the studied exponential-in-$p$ interpolation-repetition comparator. The oracle lower bound concerns unknown acceptance tables in the folded family.

\dimensionlabel{Strength}{strengthPolynomialTime} Candidate exponential-to-polynomial advantage for the explicit random-instance approximation target; a superpolynomial query separation for the folded oracle family. The mechanism is decoded Fourier interference.

\dimensionlabel{Evidence}{evidencePartial} Partial evidence for the explicit-input comparison, based on the analyzed classical algorithms. The folded oracle comparison has a rigorous classical query lower bound.

\dimensionlabel{Hardware}{hardwarePolynomial} Fault-tolerant DQI requires approximately $O((p+n)\log p)$ logical data qubits plus decoding workspace. A cited $p=521$ design uses about 8,850 logical qubits and $1.01\times10^8$ Toffoli gates; decoder construction and error correction determine physical scale.

\subhead{Algorithm}
\begin{algorithmindex}{60}
\item \href{https://arxiv.org/pdf/2408.08292}{Jordan et al., \emph{Optimization by Decoded Quantum Interferometry} (Nature 2025)}. Theorem 4.1; Section 5, ``Optimal Polynomial Intersection''; Section 8, ``Decoded Quantum Interferometry''; Section 16 for resources.
\item \href{https://arxiv.org/pdf/2607.14650}{Horinaga and Yamakawa, \emph{Worst-Case Quantum Algorithm for Optimal Polynomial Intersection Beyond Decoded Quantum Interferometry} (2026)}. Theorems 1.1 and 6.1.
\item \href{https://arxiv.org/pdf/2610.02145}{Kramer et al., \emph{A Provable Quantum Advantage for Approximate Optimization via Decoded Quantum Interferometry} (2026)}. Theorem 1: folded-OPI quantum--classical query separation.
\item \href{https://github.com/openai/math/blob/main/preprints/The-Exponential-Time-Hypothesis-is-False-October-25-2026/paper.pdf}{OpenAI, \emph{The Exponential Time Hypothesis Is False} (25 October 2026)}. Theorem 1.1 and the sparse-instance bound; witness recovery by conditional counting.
\end{algorithmindex}

\entrytitle{103}{Maximum \texorpdfstring{$k$}{k}-cut with variational quantum circuits}{white}{evidencePartial}{hardwarePolynomial}
\setcounter{problem}{102}
\begin{problem}[Maximum $k$-cut]
Given an undirected graph $G=(V,E)$, nonnegative rational edge weights $w_{uv}$ and an integer $k\ge3$, find $c:V\to\{0,\ldots,k-1\}$ maximizing
\[
C(c)=\sum_{\{u,v\}\in E}w_{uv}\,\mathbf1[c(u)\ne c(v)].
\]
The approximation version requests a labeling with a specified value relative to the optimum or total edge weight. Write $n=|V|$ and $m=|E|$. The problem models partitioning into several groups and extends the binary cut objective. Complexity counts quantum gates and measurement samples, classical arithmetic for approximation, and bit operations for exact finite-precision optimization.
\end{problem}
\speedupheading
\complexity{Quantum}{$O(KSpk(m+n))$ gates, one-hot QAOA with a parity-ring mixer; polynomial-size qudit QAOA with a Grover mixer.}
Here $p$ is the number of alternating layers, $K$ counts evaluated parameter settings and $S$ is the number of samples per setting. The one-hot encoding uses $nk$ data qubits; its cost layer has $O(km)$ two-qubit phases and its mixer has $O(kn)$ two-qubit rotations. Each sample is a classical $k$-labeling. Binary encodings use $n\lceil\log_2k\rceil$ data qubits and additional gates for unused labels. The classical parameter optimization and gate-synthesis precision contribute to total cost. For the qudit Grover-mixer algorithm on unweighted $d$-regular graphs of girth $g\ge2p+2$, Apte et al. compute the expected cut fraction and identify $k=3,4$ comparisons in which depths starting at $p=9$ exceed the analyzed classical Local Vector guarantees. Their expectation calculation costs $O(pk^{2(p+1)})$ classical time and $O(k^{2p})$ memory, with $d$ entering the numerical values.

\complexity{Classical}{Polynomial time, Frieze--Jerrum approximation and the Local Vector algorithm; $\operatorname{poly}(n,m,W)k^{\omega n/3}$ for exact integer-weighted 2-CSP.}
Frieze--Jerrum rounding gives approximation ratio $\alpha_k$, with $\alpha_3\simeq0.836008$ and $\alpha_k=1-1/k+(2+o(1))\log(k)/k^2$. Garc\'ia-Soriano and Miyauchi give certified ratios for $3\le k\le16$ and characterize the ratio for every $k\ge3$. On regular graphs with a girth promise, the Local Vector algorithm supplies the corresponding cut-fraction comparison. For exact optimization, $\omega$ is the matrix-multiplication exponent and $W$ bounds the total integer score after clearing denominators. A threshold circuit with $L$ gates also gives $\operatorname{poly}(L)2^{O(L/\sqrt{\log L})}$ time through exact 3-CNF counting. The smaller applicable exact bound is used; at fixed $k$, unit weights and $m=O(n)$, the latter becomes $\operatorname{poly}(n)2^{O(n/\sqrt{\log n})}$.

\dimensionlabel{Strength}{white} The established comparison concerns approximation quality at specified depth and graph promises. A runtime speedup for attaining the same target quality remains unestablished.
\dimensionlabel{Evidence}{evidencePartial} Analytical expectation formulas and the reported parameter comparisons support an advantage over the stated classical cut-fraction guarantees. The strongest general approximation and exact algorithms provide additional comparisons.
\dimensionlabel{Hardware}{hardwarePolynomial} A prospective implementation uses $nk$ one-hot logical qubits or $n\lceil\log_2k\rceil$ binary logical data qubits, plus encoding and synthesis workspace. Depth, graph routing and repeated sampling control the required fault-tolerant resources; the cited encoding and high-girth studies use classical simulations and analysis.
\subhead{Algorithm}
\begin{algorithmindex}{103}
\item \href{https://arxiv.org/abs/1709.03489}{Hadfield et al., \emph{From the Quantum Approximate Optimization Algorithm to a Quantum Alternating Operator Ansatz} (2019)}. Section 3 and Appendix A.3.1: feasible-state encodings, mixers and Max-$k$-Cut mapping.
\item \href{https://arxiv.org/abs/2009.01095}{Fuchs et al., \emph{Efficient Encoding of the Weighted MAX $k$-CUT on a Quantum Computer Using QAOA} (2021)}. Sections 3.1--3.3, Equations (27)--(28) and Table 1: one-hot/binary encodings and gate counts.
\item \href{https://arxiv.org/abs/2609.39042}{Apte et al., \emph{Provable Classical and Quantum Local Algorithms for Max-$k$-Cut and Quantum Advantage at Moderate Girth} (2026)}. Algorithm 1; Sections 4.1--4.2 define and evaluate qudit QAOA; Section 5 and Appendix A give the comparisons.
\item \href{https://arxiv.org/abs/2609.35453}{Garc\'ia-Soriano and Miyauchi, \emph{Hermite Brings a Laptop: Analyzing Frieze--Jerrum Rounding Yields Improved Approximations for Clustering Problems} (2026)}. Section 4.4, Theorems 4.8--4.12: the Frieze--Jerrum ratio and its certified evaluation.
\item \href{https://www.cs.helsinki.fi/u/mkhkoivi/publications/ipl-2005.pdf}{Koivisto, \emph{Optimal 2-Constraint Satisfaction via Sum-Product Algorithms} (2006)}. Sections 2--4: sum-product reduction and exact weighted 2-CSP optimization.
\item \href{https://github.com/openai/math/blob/main/preprints/The-Exponential-Time-Hypothesis-is-False-October-25-2026/paper.pdf}{OpenAI, \emph{The Exponential Time Hypothesis Is False} (2026)}. Theorem 1.1: sparse fixed-width exact counting; threshold circuits and witness recovery give the exact optimization bound.
\end{algorithmindex}

\entrytitle{104}{Maximum bisection with constraint-preserving QAOA}{white}{white}{hardwarePolynomial}
\setcounter{problem}{103}
\begin{problem}[Maximum bisection]
Given an undirected graph $G=(V,E)$ with even $n=|V|$, $m=|E|$ and nonnegative rational weights, find $S\subset V$ with $|S|=n/2$ maximizing $\sum_{u\in S,v\notin S}w_{uv}$. The balance constraint models an equal-size partition of a network. Complexity counts quantum gates and samples, classical arithmetic for approximation, and bit operations for exact optimization.
\end{problem}
\speedupheading
\complexity{Quantum}{$O(KSp(m+n))$ gates, the quantum alternating operator ansatz with an $XY$ parity-ring mixer.}
The initial computational basis state has Hamming weight $n/2$. Each of $p$ layers applies $m$ weighted $ZZ$ phases followed by $n$ two-qubit $XY$ rotations arranged in parity layers. These rotations preserve Hamming weight, so every measured bit string describes a balanced partition. Here $K$ is the number of parameter settings evaluated and $S$ the samples per setting; the classical optimizer and rotation-synthesis accuracy supply further costs.

\complexity{Classical}{Polynomial time with ratio $0.8776$, Austrin--Benabbas--Georgiou; $\operatorname{poly}(n,m,W)2^{\omega n/3}$ for exact weighted 2-CSP.}
The approximation algorithm rounds a semidefinite hierarchy solution with biased thresholds and restores exact balance. For exact finite-precision optimization, a quadratic cardinality penalty produces a weighted 2-CSP, where $W$ includes the integer penalty and score range and $\omega$ is the matrix-multiplication exponent. Exact threshold counting also gives $\operatorname{poly}(L)2^{O(L/\sqrt{\log L})}$ for a size-$L$ feasibility-and-score circuit. The smaller applicable exact bound is used. For sparse unweighted graphs, balance and the cut threshold admit $L=O(n)$, yielding $\operatorname{poly}(n)2^{O(n/\sqrt{\log n})}$.

\dimensionlabel{Strength}{white} The variational construction supplies feasible balanced solutions with the displayed evaluation cost. Its approximation and time-to-solution performance have no established asymptotic advantage over these classical algorithms.
\dimensionlabel{Evidence}{white} Constraint preservation and circuit-resource bounds are rigorous. The number of parameter evaluations and samples needed for a prescribed approximation guarantee is instance dependent.
\dimensionlabel{Hardware}{hardwarePolynomial} The circuit uses $n$ data qubits. The parity-ring mixer has constant depth before routing; the cost layer has depth at most the graph's maximum degree plus one with suitable connectivity. Protected implementations must also support parameter optimization, repeated measurements and phase synthesis.
\subhead{Algorithm}
\begin{algorithmindex}{104}
\item \href{https://arxiv.org/abs/1709.03489}{Hadfield et al., \emph{From the Quantum Approximate Optimization Algorithm to a Quantum Alternating Operator Ansatz} (2019)}. Appendix A.3.2, ``Graph Partitioning (Minimum Bisection)'', and A.3.3, ``Maximum Bisection'': complete mappings and resource counts.
\item \href{https://arxiv.org/abs/1205.0458}{Austrin, Benabbas and Georgiou, \emph{Better Balance by Being Biased: A 0.8776-Approximation for Max Bisection} (2013/2016)}. Theorem 1.1 and the rounding algorithm.
\item \href{https://www.cs.helsinki.fi/u/mkhkoivi/publications/ipl-2005.pdf}{Koivisto, \emph{Optimal 2-Constraint Satisfaction via Sum-Product Algorithms} (2006)}. Sections 2--4: exact weighted 2-CSP optimization.
\item \href{https://github.com/openai/math/blob/main/preprints/The-Exponential-Time-Hypothesis-is-False-October-25-2026/paper.pdf}{OpenAI, \emph{The Exponential Time Hypothesis Is False} (2026)}. Theorem 1.1; sparse threshold-circuit encoding and conditional counting.
\end{algorithmindex}

\entrytitle{105}{Job-shop scheduling with delivery and changeover costs}{white}{white}{hardwareSmall}
\setcounter{problem}{104}
\begin{problem}[Discrete production scheduling]
Given $J$ unit-time jobs, $M$ machines, due dates $d_j$, idle allowances $e_m$, production groups $P_{mj}$ and nonnegative costs $c_e,c_l,c_p$, choose distinct integer slots $t_{mj}\in\{1,\ldots,J+e_m\}$ for the jobs on each machine. Require $t_{m+1,j}\ge t_{mj}$ and that each machine's occupied slots form one consecutive interval. If $\pi_m(r)$ is its $r$th processed job, minimize
\[
c_e\sum_j(d_j-t_{Mj})_++c_l\sum_j(t_{Mj}-d_j)_+
+c_p\sum_{m=1}^{M}\sum_{r=1}^{J-1}\mathbf1[P_{m,\pi_m(r)}\ne P_{m,\pi_m(r+1)}],
\]
where $(a)_+=\max(a,0)$. This scheduling model captures delivery penalties and production-group changes in steel manufacturing. Complexity counts quantum gates and measurement samples and classical bit operations for the time-indexed binary formulation.
\end{problem}
\speedupheading
\complexity{Quantum}{$O(KSpN)$ gates for hardware-efficient VQE and F-VQE; $O(KSp(N+q))$ for QAOA.}
The time-indexed QUBO has $N=\sum_m J(J+e_m)+\sum_m e_m$ binary variables before presolving, and $q$ nonzero quadratic terms. Here $p$ counts circuit layers, $K$ counts primitive circuit evaluations and $S$ the shots per evaluation. A computational-basis sample gives a schedule assignment and its classically evaluated cost. VQE minimizes the sampled energy or a conditional-value-at-risk objective; filtering VQE (F-VQE) reweights sampled energies and uses parameter-shift updates. Variational imaginary-time evolution (VarQITE) uses the same state family and estimates a parameter metric: with $d$ parameters, a direct iteration needs $O(d^2+dq)$ measurement settings and an $O(d^3)$ classical linear-system solve. These settings contribute to $K$. QAOA uses phases for the QUBO terms and a transverse-field mixer. Parameter optimization and feasibility checks are part of the hybrid computation.

\complexity{Classical}{$\min\{O^*(2^N),\operatorname{poly}(N,q,W)2^{\omega N/3},\operatorname{poly}(L)2^{O(L/\sqrt{\log L})}\}$, enumeration, weighted 2-CSP and exact threshold counting.}
Here $W$ bounds the integer QUBO score range after scaling, $\omega$ is the matrix-multiplication exponent, and $L$ is the size of a Boolean circuit checking the score threshold and constraints. These bounds apply to the same time-indexed instance, including its horizon and coefficient precision. Classical branch-and-bound, constraint programming and mixed-integer solvers use scheduling-specific propagation and lower bounds; their instance-dependent running times supply practical comparisons. The raw QUBO variable count can be reduced by fixing assignments and simplifying constraints before either classical or quantum optimization.

\dimensionlabel{Strength}{white} Circuit and sampling costs are parameterized by the optimization effort. The published study establishes empirical convergence and solution-quality comparisons among the variational methods.
\dimensionlabel{Evidence}{white} The hardware experiments compare sampled energies and optimum frequencies on the tested instances; an asymptotic quantum speedup has not been established.
\dimensionlabel{Hardware}{hardwareSmall} Amaro et al. compared the four methods on a five-qubit instance and ran F-VQE on instances with up to 23 physical superconducting qubits after presolving. VQE, F-VQE and QAOA use $N$ data qubits; the Hadamard-test implementation of VarQITE adds an ancilla. Logical implementations must protect the complete repeated-measurement computation.
\subhead{Algorithm}
\begin{algorithmindex}{105}
\item \href{https://arxiv.org/abs/2109.03745}{Amaro et al., \emph{A Case Study of Variational Quantum Algorithms for a Job Shop Scheduling Problem} (2022)}. Sections II.A--II.B and Appendix A define the scheduling problem and QUBO; Sections II.C.1--II.C.4 and Table II specify VQE, QAOA, VarQITE and F-VQE; Section III gives the experiments.
\item \href{https://arxiv.org/abs/2106.10055}{Amaro et al., \emph{Filtering Variational Quantum Algorithms for Combinatorial Optimization} (2022)}. Filtering objective, parameter-shift update and adaptive-filter prescription.
\item \href{https://www.cs.helsinki.fi/u/mkhkoivi/publications/ipl-2005.pdf}{Koivisto, \emph{Optimal 2-Constraint Satisfaction via Sum-Product Algorithms} (2006)}. Sections 2--4: exact optimization of integer-weighted pairwise objectives.
\item \href{https://github.com/openai/math/blob/main/preprints/The-Exponential-Time-Hypothesis-is-False-October-25-2026/paper.pdf}{OpenAI, \emph{The Exponential Time Hypothesis Is False} (2026)}. Theorem 1.1; exact threshold counting and witness recovery.
\end{algorithmindex}

\entrytitle{106}{Discrete portfolio rebalancing with constrained mixers}{white}{white}{hardwarePolynomial}
\setcounter{problem}{105}
\begin{problem}[Discrete portfolio rebalancing]
Given $n$ assets, a positive-semidefinite rational covariance matrix $\Sigma$, expected returns $\mu$, previous positions $y\in\{-1,0,1\}^n$, a net investment $D$, $0\le\lambda\le1$ and transaction cost $T\ge0$, find $z\in\{-1,0,1\}^n$ satisfying $\sum_i z_i=D$ and minimizing
\[
\lambda z^T\Sigma z-(1-\lambda)\mu^Tz+T\sum_i\mathbf1[z_i\ne y_i].
\]
The positions represent short, absent and long holdings in discrete lots. Complexity counts circuit gates and samples, with classical bit complexity depending on the coefficient precision.
\end{problem}
\speedupheading
\complexity{Quantum}{$O(KSpn^2)$ gates for dense covariance, the quantum alternating operator ansatz with paired parity mixers.}
The two-spin encoding writes $z_i=x_i^+-x_i^-$ using $2n$ data qubits. An entangled initial state has the required difference between long and short occupations. Separate $XY$ parity-ring mixers preserve their respective occupations and hence the investment constraint. A cost layer implements the risk, return and trading-cost phases; sampling and decoding return discrete portfolios. The parameters $p$, $K$ and $S$ count alternating layers, evaluated parameter settings and shots per setting. Classical optimizer work and cost evaluation of the sampled portfolios add to the displayed gate count.

\complexity{Classical}{$\min\{O^*(3^n),\operatorname{poly}(n,W)3^{\omega n/3},\operatorname{poly}(L)2^{O(L/\sqrt{\log L})}\}$, enumeration, weighted 2-CSP and exact threshold counting.}
The investment equality can be imposed by a sufficiently large quadratic penalty. Unary transaction costs and pairwise risk terms then form a three-state 2-CSP. Here $W$ includes the scaled integer coefficient and penalty ranges; the polynomial dependence on $W$ makes the matrix-multiplication bound pseudopolynomial in coefficient bit length. The threshold-circuit bound instead uses the binary encoding size $L$ of feasibility and objective evaluation. Mixed-integer quadratic programming with branch-and-bound supplies the practical classical solver; covariance structure and valid bounds determine its instance-specific search.

\dimensionlabel{Strength}{white} The algorithm preserves the investment constraint and produces sampled discrete portfolios. The established resource formula is a circuit-evaluation cost, with optimization effort specified by $K$ and $S$.
\dimensionlabel{Evidence}{white} The reported solution-quality results concern small simulated portfolios. A complexity separation for this rebalancing problem remains open.
\dimensionlabel{Hardware}{hardwarePolynomial} The encoding requires $2n$ logical data qubits plus compilation workspace. The cited eight-asset study uses an ideal simulator. Dense covariance gives $\Theta(n^2)$ interactions per cost layer, and routing, angle accuracy and repeated sampling dominate scaling on physical devices.
\subhead{Algorithm}
\begin{algorithmindex}{106}
\item \href{https://arxiv.org/abs/1911.05296}{Hodson et al., \emph{Portfolio Rebalancing Experiments Using the Quantum Alternating Operator Ansatz} (2019)}. Equations (9)--(12): optimization problem; Section IV and Equations (13)--(24): encoding, constraints, initial state and paired parity mixers; Section V: simulated experiments.
\item \href{https://www.cs.helsinki.fi/u/mkhkoivi/publications/ipl-2005.pdf}{Koivisto, \emph{Optimal 2-Constraint Satisfaction via Sum-Product Algorithms} (2006)}. Sections 2--4, specialized to a three-element variable domain.
\item \href{https://github.com/openai/math/blob/main/preprints/The-Exponential-Time-Hypothesis-is-False-October-25-2026/paper.pdf}{OpenAI, \emph{The Exponential Time Hypothesis Is False} (2026)}. Theorem 1.1; binary threshold-circuit encoding and witness recovery.
\end{algorithmindex}

\entrytitle{107}{Two-community modularity maximization by variational local search}{white}{white}{hardwareSmall}
\setcounter{problem}{106}
\begin{problem}[Two-community modularity maximization]
Given a simple undirected graph with adjacency matrix $A$, $m>0$ edges and degrees $d_i$, find $s\in\{-1,1\}^n$ maximizing
\[
Q(s)=\frac{1}{4m}\sum_{i,j=1}^n\left(A_{ij}-\frac{d_i d_j}{2m}\right)s_i s_j.
\]
The signs specify at most two communities; modularity measures their excess internal connectivity relative to the degree-based null model. Complexity counts quantum gates and samples for each local subproblem and classical graph-processing and bit-operation costs for the complete search.
\end{problem}
\speedupheading
\complexity{Quantum}{$O(IKSpb^2)$ gates, variational quantum local search on $b$-vertex subproblems.}
In each of $I$ local-search rounds, a classical routine chooses $b$ vertices, fixes the remaining signs, and forms a quadratic objective with linear terms from the fixed vertices. A variational circuit samples candidate updates, and an improving update is accepted. The bound uses a dense depth-$p$ ansatz or cost circuit, $K$ parameter evaluations per round and $S$ shots per evaluation; the classical selection and update work is additional. Shaydulin et al. describe the gate-model solver as QAOA and implement an $R_yR_z$ variational ansatz with COBYLA parameter optimization. The hardware runs use angles optimized on a simulator, whose cost is part of that experiment's computation.

\complexity{Classical}{$O^*(2^{\omega n/3})$ exact global optimization, weighted 2-CSP; $O^*(2^{\omega b/3})$ per exact local subproblem.}
Multiplying the objective by $8m^2$ gives integer pair coefficients $2mA_{ij}-d_i d_j$ of polynomial magnitude. Shifting each pair's score to be nonnegative permits the Williams--Koivisto weighted 2-CSP algorithm, with $\omega$ the matrix-multiplication exponent. A size-$L$ Boolean threshold circuit also supplies $\operatorname{poly}(L)2^{O(L/\sqrt{\log L})}$ time; the smaller applicable bound is used. The published practical comparison uses Gurobi both for global optimization and for the same local subproblems. Local search with an exact subproblem solver optimizes the selected neighborhood at each round, while global search supplies an optimum certificate.

\dimensionlabel{Strength}{white} The hybrid method improves modularity by optimizing small neighborhoods. Its measured solution quality and iteration count provide an empirical comparison with classical local search.
\dimensionlabel{Evidence}{white} The study reports comparable solution quality across the tested solvers. A same-task asymptotic runtime advantage is unestablished.
\dimensionlabel{Hardware}{hardwareSmall} The gate-model experiment uses 16 physical superconducting qubits for local subproblems inside graphs with up to 410 vertices. A logical realization uses $b$ data qubits plus circuit workspace; the full graph and current partition reside in classical memory. The number of graph vertices is therefore distinct from the quantum register size.
\subhead{Algorithm}
\begin{algorithmindex}{107}
\item \href{https://arxiv.org/abs/1810.12484}{Shaydulin et al., \emph{Network Community Detection on Small Quantum Computers} (2019)}. Equation (1): two-community modularity; Algorithm 1 and Equation (2): quantum local search and its subproblem; Section 3.2: gate-model implementation; Section 4: experiments.
\item \href{https://people.csail.mit.edu/rrw/2-csp-final.pdf}{Williams, \emph{A New Algorithm for Optimal Constraint Satisfaction and Its Implications} (2005)}. Split-and-list reduction of 2-CSP to weighted triangles and matrix multiplication.
\item \href{https://www.cs.helsinki.fi/u/mkhkoivi/publications/ipl-2005.pdf}{Koivisto, \emph{Optimal 2-Constraint Satisfaction via Sum-Product Algorithms} (2006)}. Sections 2--4: improved weighted 2-CSP implementation.
\item \href{https://github.com/openai/math/blob/main/preprints/The-Exponential-Time-Hypothesis-is-False-October-25-2026/paper.pdf}{OpenAI, \emph{The Exponential Time Hypothesis Is False} (2026)}. Theorem 1.1; threshold-circuit encoding and conditional counting.
\end{algorithmindex}

\clearpage
\section{Cryptanalysis and post-quantum problems}

\entrytitle{61}{Integer factorization}{strengthPolynomialTime}{evidenceStrong}{hardwarePolynomial}
\setcounter{problem}{60}
\begin{problem}[Integer factorization]
Given an odd composite integer $N$ in binary, output primes $p_i$ and positive integers $e_i$ such that $N=\prod_i p_i^{e_i}$. Write $n=\lceil\log_2N\rceil$. The balanced semiprime case underlies RSA. The explicit-integer comparison uses quantum gate complexity and classical bit complexity.
\end{problem}

\speedupheading
\complexity{Quantum}{$\widetilde O(n^2)$ total gates, Shor and Regev factoring.}
Shor's order-finding algorithm uses $\widetilde O(n^3)$ gates with elementary arithmetic and $\widetilde O(n^2)$ with fast arithmetic; its logical space is $O(n)$. Regev uses $\widetilde O(n^{3/2})$ gates per circuit and $\lceil\sqrt n\rceil+4$ repetitions, followed by polynomial-time lattice postprocessing. The total quantum gate count is $\widetilde O(n^2)$, with shorter individual coherent computations. Pilatte establishes unconditional correctness for the recent factoring framework.

\complexity{Classical}{$L_N[1/3,(64/9)^{1/3}]$ heuristic time, general number field sieve.}
Here $L_N[\alpha,c]=\exp((c+o(1))(\log N)^\alpha(\log\log N)^{1-\alpha})$. This is the standard fastest heuristic comparator for general balanced RSA semiprimes. It includes smooth-relation collection and sparse linear algebra. Rudich--Rousseau also announce deterministic worst-case $N^{o(1)}=2^{o(n)}$ time using the Boolean tensor-contraction theorem. That result specifies a deterministic guarantee; its $o(n)$ exponent does not quantitatively rank it against heuristic GNFS.

\dimensionlabel{Strength}{strengthPolynomialTime} Subexponential to polynomial in the input bit length, through Fourier sampling.

\dimensionlabel{Evidence}{evidenceStrong} Strong evidence under the standard assumption that factoring is classically hard in polynomial time. The algorithms give upper bounds in the explicit-integer model.

\dimensionlabel{Hardware}{hardwarePolynomial} Small factoring circuits have been demonstrated. RSA-sized instances require fault tolerance and polynomial arithmetic workspace. Gidney's RSA-2048 estimate uses about 1,409 logical qubits and projects fewer than one million physical qubits and an approximately one-week runtime under the specified architecture assumptions.


\subhead{Algorithm}
\begin{algorithmindex}{61}
\item \href{https://arxiv.org/pdf/quant-ph/9508027}{Shor, \emph{Polynomial-Time Algorithms for Prime Factorization and Discrete Logarithms on a Quantum Computer} (1997)}. Section 5: order finding and factor extraction.
\item \href{https://arxiv.org/pdf/2308.06572}{Regev, \emph{An Efficient Quantum Factoring Algorithm} (2023)}. Theorem 1.1: circuit cost and repetition count; Section 3: quantum procedure.
\item \href{https://www.cambridge.org/core/journals/forum-of-mathematics-pi/article/unconditional-correctness-of-recent-quantum-algorithms-for-factoring-and-computing-discrete-logarithms/3077F5BD48D1F62000B1ED800EB86308}{Pilatte, \emph{Unconditional Correctness of Recent Quantum Algorithms for Factoring and Computing Discrete Logarithms} (2026)}. Main factoring correctness theorem.
\item \href{https://arxiv.org/pdf/2505.15917}{Gidney, \emph{How to Factor 2048 Bit RSA Integers with Less Than a Million Noisy Qubits} (2025)}. Resource-estimate summary and architecture assumptions.
\item \href{https://hexagonmath.org/2610.00600v1}{Rudich and Rousseau, \emph{Deterministic Factoring in Time $N^{o(1)}$} (2026)}. Main deterministic factoring statement.
\end{algorithmindex}


\entrytitle{62}{Discrete logarithms in finite fields}{strengthPolynomialTime}{evidenceStrong}{hardwarePolynomial}
\setcounter{problem}{61}
\begin{problem}[Discrete logarithms in finite fields]
Given a finite field $\mathbb F_q$, an element $g$ of order $r$, and $h\in\langle g\rangle$, find $x\in\mathbb Z/r\mathbb Z$ with $g^x=h$. The input length is $\Theta(\log q)$ plus the field representation. Prime-field Diffie--Hellman is the principal cryptographic application. The explicit-field comparison uses quantum gate complexity and classical bit complexity; the black-box comparison counts group operations.
\end{problem}

\speedupheading
\complexity{Quantum}{$O(n^3)$ gates and $O(n)$ logical qubits, Shor discrete logarithms.}
For $n=\lceil\log_2q\rceil$, this conventional bound includes reversible field exponentiation and Fourier sampling; fast multiplication improves the arithmetic cost. The group-operation query count is polynomial in $n$. The output is the classical exponent $x$, with constant success probability amplified by repetition.

\complexity{Classical}{$L_q[1/3,(64/9)^{1/3}]$ heuristic time in general large prime fields, number field sieve discrete logarithms.}
The bound includes relation collection, sparse linear algebra and individual-logarithm descent. Special field representations can improve the $L$-constant. For applicable families in small fixed characteristic, Barbulescu--Gaudry--Joux--Thom\'e and subsequent rigorous developments give quasi-polynomial time $\exp((\log\log q)^{O(1)})$. The characteristic and field representation therefore determine the matched classical comparator.

\dimensionlabel{Strength}{strengthPolynomialTime} Subexponential to polynomial for large prime fields; quasi-polynomial to polynomial in the stated small-characteristic families. The quantum mechanism is abelian Fourier sampling.

\dimensionlabel{Evidence}{evidenceStrong} Strong conditional evidence under the corresponding classical discrete-logarithm hardness assumption; generic-group lower bounds apply to the separately specified oracle model.

\dimensionlabel{Hardware}{hardwarePolynomial} Cryptographic fields require fault-tolerant reversible arithmetic. The data and arithmetic workspace use $O(n)$ logical qubits in space-efficient implementations, with polynomially many coherent field multiplications. Physical-qubit counts depend on the chosen arithmetic circuits and error-correction architecture.


\subhead{Algorithm}
\begin{algorithmindex}{62}
\item \href{https://arxiv.org/pdf/quant-ph/9508027}{Shor, \emph{Polynomial-Time Algorithms for Prime Factorization and Discrete Logarithms on a Quantum Computer} (1997)}. Section 6, equations (6.1)--(6.3): discrete-logarithm algorithm.
\item \href{https://arxiv.org/pdf/1306.4244}{Barbulescu, Gaudry, Joux and Thom\'e, \emph{A Heuristic Quasi-Polynomial Algorithm for Discrete Logarithm in Finite Fields of Small Characteristic} (2014)}. Theorem 3: the quasi-polynomial bound for sparse medium subfield representations; descent construction in the following sections.
\end{algorithmindex}


\entrytitle{63}{Elliptic-curve discrete logarithm}{strengthPolynomialTime}{evidenceStrong}{hardwarePolynomial}
\setcounter{problem}{62}
\begin{problem}[Elliptic-curve discrete logarithm]
Given an elliptic curve $E/\mathbb F_p$ with $p>3$ and prime-order group $E(\mathbb F_p)$ of order $r$, a nonzero point $P$, and a point $Q$, find
\[
 x\in\mathbb Z/r\mathbb Z\quad\text{such that}\quad Q=xP.
\]
The curve and points are supplied by field coordinates; let $n=\lceil\log_2p\rceil$. This is the discrete-logarithm problem underlying prime-field elliptic-curve public-key cryptography.
The explicit-coordinate comparison uses quantum gate complexity and expected classical bit complexity; the generic-group comparison counts group-oracle queries.
\end{problem}

\speedupheading
\complexity{Quantum}{$\widetilde O(n^3)$ gates and $O(n)$ logical qubits, Shor's discrete-logarithm algorithm.}
This conventional bound uses reversible elliptic-curve arithmetic; implementation-dependent logarithmic factors arise from modular inversion and circuit synthesis. The input and output are classical, and the success probability can be amplified by repetition. Fast arithmetic and optimized point-addition circuits improve the detailed gate counts.

\complexity{Classical}{$2^{O(n/(\log n)^{1/4})}$ expected time, the Akshima et al. index-calculus algorithm.}
The result uses the Boolean tensor-contraction theorem in \emph{The Exponential Time Hypothesis Is False}. It applies to explicit prime-field coordinates and extends to fixed cofactors. Its analysis balances relation-search and factor-base costs of $2^{O(mn/\sqrt{\log n})}$ and $2^{O(n/m)}$, with $m\asymp(\log n)^{1/4}$. In the generic-group oracle model, Pollard's rho and the matching generic lower bound give $\widetilde\Theta(\sqrt r)$ classical group operations.

\dimensionlabel{Strength}{strengthPolynomialTime} Subexponential to polynomial for explicit prime-field curves, through Fourier sampling. The generic-group comparison is exponential to polynomial.

\dimensionlabel{Evidence}{evidenceStrong} Strong conditional evidence for explicit curves under polynomial-time ECDLP hardness; a rigorous separation in the generic-group model. The subexponential coordinate algorithm is stated in a 2026 preprint and depends on the tensor-contraction theorem.

\dimensionlabel{Hardware}{hardwarePolynomial} Cryptographic instances require fault tolerance. The Roetteler--Naehrig--Svore--Lauter construction for a 256-bit prime field uses 2,330 logical qubits and approximately $1.26\times10^{11}$ Toffoli gates. Its arithmetic workspace is polynomial in the field bit length.

{\small [Before ``ETH is false'', the known general-curve classical bound was $\widetilde O(2^{n/2})$ from Pollard's rho, and Shor's algorithm claimed an exponential-to-polynomial speedup.]\par}

\subhead{Algorithm}
\begin{algorithmindex}{E}
\item \href{https://arxiv.org/pdf/quant-ph/9508027}{Shor, \emph{Polynomial-Time Algorithms for Prime Factorization and Discrete Logarithms on a Quantum Computer} (1997)}. Section 6, equations (6.1)--(6.3): discrete-logarithm algorithm.
\item \href{https://arxiv.org/pdf/1706.06752}{Roetteler, Naehrig, Svore and Lauter, \emph{Quantum Resource Estimates for Computing Elliptic Curve Discrete Logarithms} (ASIACRYPT 2017)}. Section 2.2 and Figures 1--2: ECDLP circuit; Section 4, Algorithm 1: point addition; Section 5, Table 2: resource counts.
\item \href{https://github.com/bochum-crypto/subexponential-ecdl/blob/main/subexponential-ecdl.pdf}{Akshima et al., \emph{A Subexponential Time Algorithm for Elliptic Curve Discrete Logarithm}, v0.2 (29 October 2026)}. Theorem 4, p.~11; Section 3.4, pp.~15--17, and Figure 2, p.~16: full algorithm; Section 3.5: running time.
\end{algorithmindex}


\entrytitle{64}{Order finding in a black-box cyclic group}{strengthPolynomialTime}{evidencePartial}{hardwarePolynomial}
\setcounter{problem}{63}
\begin{problem}[Order finding in a black-box cyclic group]
Given an element $g$ in a uniquely encoded finite group and an upper bound $R$ on its order, output $r=\min\{t>0:g^t=1\}$. Multiplication and inversion are supplied as reversible group-operation oracles. Order finding is the algebraic primitive connecting Fourier sampling to factorization and discrete logarithms. Complexity is measured in group-oracle operations together with quantum gates and classical bit operations polynomial in the encoding length.
\end{problem}

\speedupheading
\complexity{Quantum}{$\operatorname{poly}(\log R)$ group operations, Shor order finding.}
The auxiliary gate cost is polynomial in $\log R$ and the group-element encoding length. An $O(\log R)$-qubit Fourier register supplies the required precision. Classical continued-fraction and least-common-multiple computations are polynomial in these parameters. For modular multiplication, the gate bounds specialize to entry 61.

\complexity{Classical}{$\widetilde O(\sqrt R)$ group operations, baby-step--giant-step and collision methods.}
Baby-step--giant-step uses $\widetilde O(\sqrt R)$ storage; collision methods supply low-memory alternatives in their applicable order-search settings. A tighter interval or a known factored multiple of the order can improve the bound. In explicit arithmetic groups, factorization and index calculus give additional instance-specific comparators. The displayed comparison supplies only the group-operation access and an order bound.

\dimensionlabel{Strength}{strengthPolynomialTime} Exponential to polynomial in $\log R$ relative to generic order-search algorithms. Fourier sampling extracts global periodic information from the group operation.

\dimensionlabel{Evidence}{evidencePartial} Partial evidence for a separation in this general formulation; the input promises and access model determine applicable lower bounds. The concrete cryptographic consequences occur in factoring and discrete logarithms.

\dimensionlabel{Hardware}{hardwarePolynomial} The implementation needs $O(\log R)$ control qubits plus group-element and reversible-arithmetic workspace. Explicit cryptographic groups lead to long fault-tolerant multiplication circuits; the number of logical qubits is polynomial in the input length.


\subhead{Algorithm}
\begin{algorithmindex}{64}
\item \href{https://arxiv.org/pdf/quant-ph/9508027}{Shor, \emph{Polynomial-Time Algorithms for Prime Factorization and Discrete Logarithms on a Quantum Computer} (1997)}. Section 5: the order-finding computation and continued-fraction reconstruction.
\item \href{https://arxiv.org/pdf/quant-ph/0011023}{Watrous, \emph{Quantum Algorithms for Solvable Groups} (2001)}. Theorem 1: quantum computation of group order in the black-box group model; the cyclic case provides the order-finding setting.
\end{algorithmindex}


\entrytitle{65}{Pell's equation and the real quadratic regulator}{strengthPolynomialTime}{evidencePartial}{hardwarePolynomial}
\setcounter{problem}{64}
\begin{problem}[Pell's equation and the real quadratic regulator]
Given a nonsquare positive integer $d$, let $\varepsilon=x+y\sqrt d>1$ be the least positive unit with $x^2-dy^2=1$. Compute the regulator $R=\log\varepsilon$ to prescribed polynomial-bit precision, or a compact representation of the fundamental solution. Compact output is the natural representation because $x,y$ can have exponentially many bits in $\log d$. Complexity is measured in quantum gates and classical bit operations, including the required output precision.
\end{problem}

\speedupheading
\complexity{Quantum}{$\operatorname{poly}(\log d,\text{precision})$, Hallgren regulator and Pell algorithms.}
The bound uses compact output and includes ideal arithmetic, discretization of the real-periodic function and Fourier precision. A polynomial number of Fourier samples suffices to reconstruct the regulator to the specified accuracy. The compact representation describes the fundamental solution even when its fully expanded integers are exponentially long in $\log d$.

\complexity{Classical}{$L_d[1/2,O(1)]$ heuristic time, Buchmann infrastructure and index calculus.}
The comparison uses the same discriminant-scale parameter and compact-output convention. Smooth relations among small prime ideals and integer linear algebra determine the regulator. The complexity assumes the standard smoothness and relation-generation hypotheses. Continued-fraction enumeration provides an alternative whose time depends on the length of the period of $\sqrt d$.

\dimensionlabel{Strength}{strengthPolynomialTime} Subexponential to polynomial in $\log d$, through Fourier sampling of a periodic function over the reals.

\dimensionlabel{Evidence}{evidencePartial} Partial evidence from sustained study of classical regulator algorithms and real-quadratic cryptosystems. The quantum polynomial-time bound is rigorous with its precision and output conventions.

\dimensionlabel{Hardware}{hardwarePolynomial} Polynomially many logical qubits in the input and precision lengths suffice. Reduced-ideal arithmetic, coherent function evaluation and Fourier precision require fault tolerance at cryptographic scale. The asymptotic algorithm supplies a polynomial-space bound; concrete counts depend on the arithmetic implementation.


\subhead{Algorithm}
\begin{algorithmindex}{65}
\item \href{https://www.cse.psu.edu/~sjh26/pell.pdf}{Hallgren, \emph{Polynomial-Time Quantum Algorithms for Pell's Equation and the Principal Ideal Problem} (JACM 2007)}. Algorithm 2.1: reduced-ideal evaluation; Algorithm 3.1: approximate period finding for the regulator. The introduction gives the classical subexponential comparison.
\end{algorithmindex}


\entrytitle{66}{The principal ideal problem}{strengthPolynomialTime}{evidencePartial}{hardwarePolynomial}
\setcounter{problem}{65}
\begin{problem}[The principal ideal problem]
Given a number field $K$ by a defining polynomial and an integral ideal $I\subseteq\mathcal O_K$ by a basis, determine whether $I=\alpha\mathcal O_K$; if so, output a compact representation of a generator $\alpha$. The relevant size parameters are the degree, logarithmic discriminant, ideal encoding length, and output precision. Complexity is measured in quantum gates and classical bit operations on the number-field and ideal encodings, with compact output.
\end{problem}

\speedupheading
\complexity{Quantum}{Polynomial time in the input parameters, Hallgren and Biasse--Song principal-ideal algorithms.}
Hallgren treats real quadratic fields. Biasse--Song give the arbitrary-degree result through their quantum $S$-unit algorithm, using compact generators and precision sufficient for the associated logarithmic lattice. The input parameters include degree, logarithmic discriminant and ideal encoding size. The principal-ideal theorem is the relevant search statement; the related class-group theorem has a GRH hypothesis.

\complexity{Classical}{$L_{|\Delta_K|}[1/2,O(1)]$ under the standard index-calculus hypotheses in quadratic fields, relation-collection algorithms.}
The classical work combines factor-base relations, unit computation and integer linear algebra. The quadratic-field family supplies the direct subexponential comparison for identical inputs and compact generators. In arbitrary degree, the strongest applicable classical bounds depend on the relationship between degree and discriminant; specialized field families also admit distinct relation-generation tradeoffs.

\dimensionlabel{Strength}{strengthPolynomialTime} Subexponential to polynomial in the matched quadratic-field regime. Quantum Fourier methods supply the hidden-lattice information needed for the generator.

\dimensionlabel{Evidence}{evidencePartial} Partial evidence for a classical separation, based on the established number-field algorithmic frontier. Polynomial quantum solvability follows from the cited theorems with their representation and precision assumptions.

\dimensionlabel{Hardware}{hardwarePolynomial} Polynomial logical space and coherent number-field arithmetic are required. Cryptographic-size computations belong to the fault-tolerant regime, with logical-qubit counts determined by field degree, coefficient precision and the chosen reversible arithmetic.


\subhead{Algorithm}
\begin{algorithmindex}{66}
\item \href{https://www.cse.psu.edu/~sjh26/pell.pdf}{Hallgren, \emph{Polynomial-Time Quantum Algorithms for Pell's Equation and the Principal Ideal Problem} (JACM 2007)}. Algorithm 4.1: computing the distance of an ideal and the principal-ideal procedure.
\item \href{https://fangsong.info/files/pubs/BS_SODA16.pdf}{Biasse and Song, \emph{Efficient Quantum Algorithms for Computing Class Groups and Solving the Principal Ideal Problem in Arbitrary Degree Number Fields} (SODA 2016)}. Theorem 1.1: $S$-units; Theorem 1.3 and Algorithm 2: principal ideal problem.
\end{algorithmindex}


\entrytitle{67}{Computing ideal class groups}{strengthPolynomialTime}{evidencePartial}{hardwarePolynomial}
\setcounter{problem}{66}
\begin{problem}[Computing ideal class groups]
For a number field $K$, compute invariant factors and represented generators of the finite abelian ideal class group $\mathrm{Cl}(\mathcal O_K)$. Equivalently, produce an explicit decomposition into cyclic groups together with ideals representing their generators. The class group parametrizes ideal classes and important isogeny group actions. Complexity is measured in quantum gates and classical bit operations on explicit number-field and ideal representations.
\end{problem}

\speedupheading
\complexity{Quantum}{Polynomial time under GRH, Biasse--Song class-group computation.}
The parameters are the number-field degree, logarithmic discriminant and standard encoding lengths. The output gives invariant factors and represented ideal generators. The cost includes finding the relevant $S$-unit relation lattice and classical Smith normal form; GRH bounds the small-prime generating set used by the reduction.

\complexity{Classical}{$L_{|\Delta|}[1/2,O(1)]$ in imaginary quadratic fields, Hafner--McCurley and Buchmann index calculus.}
The stated bounds use the analytic and smoothness hypotheses of the corresponding relation-collection algorithm. Factor-base relations and integer linear algebra recover a presentation of the same finite abelian class group. For fields of arbitrary degree, the applicable bound varies with degree and discriminant; the imaginary-quadratic family is the explicit matched comparison here.

\dimensionlabel{Strength}{strengthPolynomialTime} Subexponential to polynomial in the quadratic-field regime, through hidden-lattice Fourier methods.

\dimensionlabel{Evidence}{evidencePartial} Partial evidence for separation. The quantum polynomial bound is conditional on GRH, and the classical subexponential algorithms have their stated analytic or smoothness assumptions.

\dimensionlabel{Hardware}{hardwarePolynomial} Polynomial logical space supports coherent number-field arithmetic and high-precision Fourier sampling. Cryptographic instances require fault tolerance; concrete workspace and gate counts depend on the field family and arithmetic circuits.


\subhead{Algorithm}
\begin{algorithmindex}{67}
\item \href{https://fangsong.info/files/pubs/BS_SODA16.pdf}{Biasse and Song, \emph{Efficient Quantum Algorithms for Computing Class Groups and Solving the Principal Ideal Problem in Arbitrary Degree Number Fields} (SODA 2016)}. Theorem 1.2: class groups under GRH; Algorithm 1: reduction of class-group computation to $S$-units; Theorem 1.1: the $S$-unit quantum subroutine.
\end{algorithmindex}


\entrytitle{68}{Vectorization for commutative isogeny actions}{strengthPolynomialTime}{evidencePartial}{hardwarePolynomial}
\setcounter{problem}{67}
\begin{problem}[Vectorization for commutative isogeny group actions]
Let a finite abelian group $G$ act freely and transitively on a set $X$. Given $E,E'\in X$, find $a\in G$ such that $a*E=E'$. For CRS and CSIDH, $G$ is an ideal class group and $X$ consists of appropriately oriented elliptic curves. Write $N=|G|$. The group-action comparison counts action-oracle queries and auxiliary operations; the explicit-isogeny comparison includes quantum gate and classical bit costs of action evaluation.
\end{problem}

\speedupheading
\complexity{Quantum}{$\operatorname{poly}(\log N)\,2^{O(\sqrt{\log|2^tp^kG|})}$ time, queries and QROM complexity, Castryck--Vander Meeren hidden-shift algorithm.}
Here $p$ is any fixed prime, $t,k$ are fixed nonnegative integers, and $2^tp^kG=\{2^tp^kg:g\in G\}$. The bound assumes efficient representations and action evaluation, and uses $\operatorname{poly}(\log N)$ qubits. It is polynomial when $2^tp^kG=\{0\}$. For general $G$, Kuperberg's collimation sieve takes $2^{O(\sqrt{\log N})}$ time with $O(\log N)$ quantum space and subexponential classical space; Regev's polynomial-space variant takes $2^{O(\sqrt{\log N\log\log N})}$ time. For explicit isogenies, Childs--Jao--Soukharev include class-group action evaluation and obtain overall $L_{|\Delta|}[1/2,O(1)]$ quantum complexity under GRH.

\complexity{Classical}{$\widetilde O(\sqrt N)$ action evaluations, baby-step--giant-step or Pollard-rho vectorization.}
This is the generic freely transitive action comparison. Each action evaluation carries its implementation cost when curves and ideals are supplied explicitly. The action model represents the hidden-shift functions by $b\mapsto b*E$ and $b\mapsto b*E'$. Efficient classical structure specific to a curve family belongs in its full arithmetic comparison.

\dimensionlabel{Strength}{strengthPolynomialTime} Exponential to polynomial in $\log N$ for generic actions with $2^tp^kG=\{0\}$ and fixed $p,t,k$. General abelian hidden shifts give exponential-to-subexponential improvement, with power saving in the exponent. These comparisons use the action-oracle model; explicit isogeny costs include arithmetic evaluation.

\dimensionlabel{Evidence}{evidencePartial} Rigorous quantum upper bounds and partial evidence relative to the best known classical vectorization algorithms. For explicit isogeny actions, classical arithmetic structure and action-evaluation cost enter the comparison.

\dimensionlabel{Hardware}{hardwarePolynomial} The bounded-torsion algorithm uses polynomially many qubits and polynomial QROM when $2^tp^kG=\{0\}$. General-group variants include $O(\log N)$ quantum space for collimation. Their subexponential state generation and classical storage, together with reversible isogeny arithmetic, make cryptographic-size execution a fault-tolerant workload.


\subhead{Algorithm}
\begin{algorithmindex}{68}
\item \href{https://eprint.iacr.org/2022/1366}{Castryck and Vander Meeren, \emph{Two Remarks on the Vectorization Problem} (2022)}. Theorem 1: the bound in terms of $|2^tp^kG|$ and polynomial quantum storage; Section 2: vectorization and hidden shift; Section 5: bounded-torsion hidden shifts.
\item \href{https://arxiv.org/pdf/1112.3333}{Kuperberg, \emph{Another Subexponential-Time Quantum Algorithm for the Dihedral Hidden Subgroup Problem} (TQC 2013)}. Algorithms 4.1 and 4.4: collimation and the collimation sieve.
\item \href{https://arxiv.org/pdf/quant-ph/0406151}{Regev, \emph{A Subexponential Time Algorithm for the Dihedral Hidden Subgroup Problem with Polynomial Space} (2004)}. Polynomial-space sieve construction and its recursive combination procedure.
\item \href{https://arxiv.org/pdf/1012.4019}{Childs, Jao and Soukharev, \emph{Constructing Elliptic Curve Isogenies in Quantum Subexponential Time} (2014)}. Algorithm 2 and Theorem 4.6: group-action evaluation; Algorithm 3 and Theorem 5.4: full isogeny algorithm; Theorems 5.1--5.2: hidden-shift subroutines.
\end{algorithmindex}


\entrytitle{69}{Supersingular isogeny path finding}{strengthPolynomial}{evidencePartial}{hardwareExponential}
\setcounter{problem}{68}
\begin{problem}[Supersingular isogeny path finding]
Given supersingular elliptic curves $E,E'/\mathbb F_{p^2}$, output an explicitly evaluable isogeny $E\to E'$. One may represent it as a composition of low-degree isogenies. This is the unrestricted path problem underlying several isogeny constructions. Complexity is measured in field-operation-scale time, with polynomial-bit arithmetic translated into quantum gates and classical bit operations.
\end{problem}

\speedupheading
\complexity{Quantum}{$p^{1/4+o(1)}$ time, Biasse--Jao--Sankar supersingular isogeny search.}
Amplitude amplification reaches the subgraph of curves defined over $\mathbb F_p$, after which the algorithm uses its abelian structure. Reversible isogeny evaluation and the endpoint-recognition predicate determine the arithmetic factors. The exponent absorbs the subexponential costs required by the second stage.

\complexity{Classical}{$p^{1/3+o(1)}$ time and memory, Wesolowski supersingular isogeny algorithm.}
The result uses a heuristic assumption on smoothness of the relevant random integers. Non-scalar endomorphism computations and reductions between endomorphism and isogeny problems produce an explicitly evaluable path. The $o(1)$ exponent includes superpolynomial factors in $\log p$. For regimes using the random-walk approach, Delfs--Galbraith provides a $p^{1/2+o(1)}$ alternative with different memory requirements.

\dimensionlabel{Strength}{strengthPolynomial} Polynomial speedup through a smaller exponential constant in the input bit length $\log p$, from $1/3$ to $1/4$. Amplitude amplification is an essential quantum component.

\dimensionlabel{Evidence}{evidencePartial} Partial evidence for explicit curves. The classical $p^{1/3+o(1)}$ comparison is a 2026 preprint with a stated smoothness heuristic.

\dimensionlabel{Hardware}{hardwareExponential} Curve registers have polynomial size in $\log p$. Reversible field arithmetic and exponentially many coherent operations require fault tolerance; runtime is the principal obstacle at cryptographic sizes.


\subhead{Algorithm}
\begin{algorithmindex}{69}
\item \href{https://cacr.uwaterloo.ca/techreports/2014/cacr2014-24.pdf}{Biasse, Jao and Sankar, \emph{A Quantum Algorithm for Computing Isogenies Between Supersingular Elliptic Curves} (INDOCRYPT 2014)}. Theorem 1: main complexity result; Algorithm 1: complete isogeny procedure; Algorithms 2--3: reaching and navigating the $\mathbb F_p$ subgraph.
\item \href{https://eprint.iacr.org/2026/1486}{Wesolowski, \emph{The Supersingular Isogeny Problem in Time and Memory $p^{1/3+o(1)}$} (2026)}. Main complexity statement and smoothness assumption in the abstract.
\end{algorithmindex}


\entrytitle{70}{SIDH secret-isogeny recovery}{white}{white}{hardwareExponential}
\setcounter{problem}{69}
\begin{problem}[SIDH secret-isogeny recovery with torsion images]
Given a supersingular starting curve $E_0$, a secret isogeny image $E_A$, and the published images $\phi_A(P_B),\phi_A(Q_B)$ of a basis of coprime torsion, recover the kernel of the secret isogeny $\phi_A$ of prescribed smooth degree. This input is the public-key recovery problem of SIDH and SIKE. The comparison uses quantum isogeny-oracle calls and the resulting gate cost, and classical bit complexity on the full SIDH public key.
\end{problem}

\speedupheading
\complexity{Quantum}{$\widetilde O(p^{1/6})$ isogeny-oracle calls, Tani claw finding applied to split isogenies.}
In the balanced SIDH parameter regime, the two half-isogeny lists contain about $p^{1/4}$ candidates each. Quantum claw search identifies matching intermediate endpoints; the stated exponent counts isogeny evaluations with the required coherent list access. Arithmetic and memory-access factors enter the full gate cost.

\complexity{Classical}{$\operatorname{poly}(\log p)$, Castryck--Decru and Robert SIDH key recovery.}
The full public key includes torsion images, which supply the extra information used by product-isogeny attacks. Robert's higher-dimensional abelian-variety framework gives polynomial-time key recovery in the specified SIDH setting. The polynomial algorithm recovers the prescribed secret isogeny or equivalent secret-key information from the same explicit public input.

\dimensionlabel{Strength}{white} The problem has a polynomial-time classical solution; the exponential quantum claw bound has no speedup over this comparator.

\dimensionlabel{Evidence}{white} The classical key-recovery theorem determines the comparison for SIDH public keys with their auxiliary torsion data.

\dimensionlabel{Hardware}{hardwareExponential} The quantum claw approach requires exponentially large coherent data structures and deep reversible isogeny arithmetic. The polynomial-time attack uses classical computation. For the quantum construction, both coherent memory and total circuit depth determine the required fault-tolerant resources.


\subhead{Algorithm}
\begin{algorithmindex}{70}
\item \href{https://arxiv.org/pdf/0708.2584}{Tani, \emph{Claw Finding Algorithms Using Quantum Walk} (2009)}. Theorem 8 and the \textsc{ClawSearch} procedure: balanced claw-finding bound used for the split-isogeny search.
\item \href{https://eprint.iacr.org/2022/975}{Castryck and Decru, \emph{An Efficient Key Recovery Attack on SIDH} (2022)}. Main key-recovery construction using auxiliary torsion images.
\item \href{https://www.math.u-bordeaux.fr/~damienrobert/pro/publications/articles/breaking_sidh.pdf}{Robert, \emph{Breaking SIDH in Polynomial Time} (EUROCRYPT 2023)}. Theorem 1.1: polynomial-time recovery through higher-dimensional isogenies.
\end{algorithmindex}


\entrytitle{71}{The Euclidean shortest vector problem}{strengthPolynomial}{evidencePartial}{hardwareExponential}
\setcounter{problem}{70}
\begin{problem}[The Euclidean shortest vector problem]
Given a basis $B\in\mathbb Q^{m\times n}$ of a rank-$n$ lattice $L=B\mathbb Z^n$, output $v\in L\setminus\{0\}$ minimizing $\|v\|_2$. Exact SVP is the central subroutine in block lattice reduction and lattice cryptanalysis. The comparison uses arithmetic or bit-operation time with polynomial factors in input precision suppressed, and quantum gate time in the stated coherent-memory model; enumeration additionally counts tree-oracle operations.
\end{problem}

\speedupheading
\complexity{Quantum}{$2^{0.247n+o(n)}$ rigorous time, Hhan--Kim--Kim--Laarhoven; $2^{0.2563n+o(n)}$ heuristic time, reusable quantum walks.}
The superlattice bound is stated in the November 2026 preprint and has a $2^{0.247n+o(n)}$ quantum time-space bound. The heuristic Bonnetain et al. sieve uses coherent access to exponential lists and the standard random-vector assumptions of lattice sieving. Aono--Nguyen--Shen quantum enumeration instead takes roughly $\sqrt T\,\operatorname{poly}(n)$ operations for a search tree with $T$ nodes, using polynomial quantum workspace. Aggarwal--Chen--Kumar--Shen provide a separate rigorous memory tradeoff: $2^{0.950n+o(n)}$ time with polynomial qubits and exponential classical memory, or $2^{0.835n+o(n)}$ with QRAM.

\complexity{Classical}{$2^{0.409n+o(n)}$ rigorous time, Hhan--Kim--Kim--Laarhoven; $2^{0.2925n+o(n)}$ heuristic time, Becker--Ducas--Gama--Laarhoven sieving.}
The superlattice algorithm uses $2^{0.259n+o(n)}$ space. The heuristic sieve uses $2^{0.2075n+o(n)}$ space. These are distinct comparisons: the rigorous quantum exponent is paired with the rigorous classical exponent, and the heuristic pair shares its distribution assumptions. Classical enumeration costs $T$ node operations for the specified tree. Input precision contributes polynomial factors to the displayed dimension exponents.

\dimensionlabel{Strength}{strengthPolynomial} Polynomial speedups through smaller exponential constants in $n$. Enumeration gives a quadratic improvement in tree size through quantum backtracking; quantum walks and amplification accelerate sieve subroutines.

\dimensionlabel{Evidence}{evidencePartial} Partial evidence for a separation. Rigorous upper bounds and heuristic sieve estimates have separate assumptions; the superlattice results are stated in a 2026 preprint.

\dimensionlabel{Hardware}{hardwareExponential} Fast sieves need exponential coherent memory and fault-tolerant access to vector lists. Enumeration uses polynomial logical qubits at greater runtime. The polynomial-qubit BDD construction explicitly separates its quantum workspace from exponential classical memory and QRAM.


\subhead{Algorithm}
\begin{algorithmindex}{71}
\item \href{https://eprint.iacr.org/2026/2761}{Hhan, Kim, Kim and Laarhoven, \emph{Superlattice Methods for the Shortest Vector Problem and Discrete Gaussian Sampling} (2026)}. Main classical and quantum SVP complexity statements.
\item \href{https://arxiv.org/pdf/2205.14023}{Bonnetain et al., \emph{Finding Many Collisions via Reusable Quantum Walks: Application to Lattice Sieving} (EUROCRYPT 2023)}. Algorithm 2: reusable collision walk; Section 7.1: the $0.2563$ SVP exponent and the classical sieve comparator.
\item \href{https://eprint.iacr.org/2018/546}{Aono, Nguyen and Shen, \emph{Quantum Lattice Enumeration and Tweaking Discrete Pruning} (ASIACRYPT 2018)}. Algorithm 1 and Theorem 6: tree minimization; Section 4.2, Theorems 7--9: cylinder pruning; Section 7: quantum discrete pruning.
\item \href{https://ir.cwi.nl/pub/35122/35122.pdf}{Aggarwal, Chen, Kumar and Shen, \emph{Improved Classical and Quantum Algorithms for the Shortest Vector Problem via Bounded Distance Decoding} (SIAM J. Comput. 2025)}. Theorem 1.3: quantum time-space bounds; Section 6: the SVP algorithms.
\end{algorithmindex}


\entrytitle{72}{Isomorphism to the standard integer lattice}{strengthPolynomial}{evidencePartial}{hardwareExponential}
\setcounter{problem}{71}
\begin{problem}[Isomorphism to the standard integer lattice]
Given the Gram matrix of a lattice $L$ promised isometric to $\mathbb Z^n$, recover an orthonormal integral lattice basis, equivalently an explicit isometry $L\to\mathbb Z^n$. This is the standard-integer-lattice case of lattice isomorphism, relevant to signature constructions such as HAWK. Complexity is measured in quantum gate time with the stated QRAM access and in classical bit-operation time, suppressing polynomial input-precision factors.
\end{problem}

\speedupheading
\complexity{Quantum}{$2^{0.1235n+o(n)}$ time, Ducas reduction with the Hhan--Kim--Kim--Laarhoven SVP subroutine.}
The exponent is obtained by substituting the $2^{0.247d+o(d)}$ quantum SVP bound into the reduction to dimension $d\simeq n/2$. It is a direct consequence of the two results. The Aggarwal--Chen--Kumar--Shen implementation gives a separately quantified tradeoff: $2^{0.417n+o(n)}$ time, $2^{0.147n+o(n)}$ QRAM, $2^{0.25n+o(n)}$ classical storage and polynomial additional qubits. Both computations recover norm-one directions sufficient to reconstruct the integral isometry.

\complexity{Classical}{$2^{0.2045n+o(n)}$ time, the same reduction with classical superlattice SVP.}
Substitution of the $2^{0.409d+o(d)}$ classical SVP theorem gives the displayed exponent. Classical discrete-Gaussian SVP with $2^{d+o(d)}$ time yields the alternative $2^{n/2+o(n)}$ bound. The reduction and output task are identical in both comparisons; the dimension-halving exponents retain the assumptions and storage costs of their chosen SVP subroutines.

\dimensionlabel{Strength}{strengthPolynomial} Polynomial speedup through a smaller exponential constant in the lattice rank. The quantitative comparison follows from the matched classical and quantum shortest-vector algorithms.

\dimensionlabel{Evidence}{evidencePartial} Partial evidence for separation. The two fastest displayed exponents depend on the 2026 superlattice preprint; the $0.417$ quantum implementation has a published rigorous resource analysis.

\dimensionlabel{Hardware}{hardwareExponential} Fault-tolerant computation with exponential storage is required for these algorithms. The published construction distinguishes polynomial logical workspace, $2^{0.147n+o(n)}$ QRAM cells and $2^{0.25n+o(n)}$ classical memory; the superlattice substitution inherits its own exponential time-space requirement.


\subhead{Algorithm}
\begin{algorithmindex}{72}
\item \href{https://ir.cwi.nl/pub/35122/35122.pdf}{Aggarwal, Chen, Kumar and Shen, \emph{Improved Classical and Quantum Algorithms for the Shortest Vector Problem via Bounded Distance Decoding} (2025)}. Section 8: reduction and quantum algorithm for isomorphism to $\mathbb Z^n$; Theorem 1.3: its SVP subroutine.
\item \href{https://eprint.iacr.org/2026/2761}{Hhan, Kim, Kim and Laarhoven, \emph{Superlattice Methods for the Shortest Vector Problem and Discrete Gaussian Sampling} (2026)}. Main classical and quantum SVP results, substituted into the dimension-halving reduction.
\end{algorithmindex}


\entrytitle{73}{Learning parity with noisy quantum examples}{strengthPolynomialTime}{evidenceStrong}{hardwareSmall}
\setcounter{problem}{72}
\begin{problem}[Learning parity with noisy quantum examples]
Recover $a\in\mathbb F_2^n$ from a noisy quantum example oracle preparing $2^{-n/2}\sum_x|x,a\cdot x\rangle$, followed by a specified noise channel. The classical learner receives the same examples measured in the computational basis. For classification noise, the final bit is flipped independently with constant probability $\eta<1/2$, producing ordinary learning-parity-with-noise samples $(x,a\cdot x+e)$. The comparison measures quantum example queries and elementary gates, and classical sample complexity and bit-operation time for measured examples.
\end{problem}

\speedupheading
\complexity{Quantum}{$O(\log(n/\delta))$ examples and $O(n\log(n/\delta))$ processing, Cross--Smith--Smolin noisy parity learning.}
These bounds hold for fixed noise rate in the analyzed noisy quantum-example models, with failure probability at most $\delta$. Under classification noise alone, $O(\log(1/\delta))$ coherent examples suffice: the Hadamard-basis parity information survives a flip of the result qubit. The logarithmic dependence on $n$ covers bitwise reconstruction in the depolarizing-noise treatment. Noise parameters affect the constants.

\complexity{Classical}{$2^{O(n/\log n)}$ time and samples at fixed noise, BKW-type LPN algorithms.}
The classical input consists of computational-basis measurements of the same examples, giving uniformly random labeled samples with independent classification noise. BKW combines samples to eliminate coordinate blocks, trading increased noise against sample count and memory. The fastest concrete attack depends on the sample budget and noise rate. These are classical noisy-sample costs for the stated access model.

\dimensionlabel{Strength}{strengthPolynomialTime} Subexponential to polynomial in $n$, through Fourier interference in the quantum examples. The availability of coherent examples is part of the input model.

\dimensionlabel{Evidence}{evidenceStrong} Strong conditional evidence from classical LPN hardness. The quantum learning theorem is rigorous for its noise channel; the comparison invokes the hardness of learning from measured examples.

\dimensionlabel{Hardware}{hardwareSmall} The data register uses $n+1$ qubits and shallow Hadamard-basis processing. Rist\`e et al. demonstrated noisy parity learning with a small superconducting processor. This supports the small-instance category; scaling the supplied example oracle and controlling noise determine larger-instance feasibility.


\subhead{Algorithm}
\begin{algorithmindex}{73}
\item \href{https://arxiv.org/pdf/1407.5088}{Cross, Smith and Smolin, \emph{Quantum Learning Robust Against Noise} (2015)}. Sections ``Learning from ideal queries'' and ``Learning in the presence of noise''; equations (7)--(8): transformed examples and classification noise; the depolarizing-noise subsection gives reconstruction bounds.
\item \href{https://arxiv.org/pdf/1512.06069}{Rist\`e et al., \emph{Demonstration of Quantum Advantage in Machine Learning} (2017)}. Experimental learning circuits and sample-complexity results in the main text and figures.
\end{algorithmindex}


\entrytitle{74}{The Euclidean closest vector problem}{strengthPolynomial}{evidencePartial}{hardwareExponential}
\setcounter{problem}{73}
\begin{problem}[The Euclidean closest vector problem]
Given a rank-$n$ lattice $L=B\mathbb Z^n$ and target $t\in\mathbb Q^m$, output $v\in L$ minimizing $\|v-t\|_2$. CVP models noisy lattice decoding and underlies attacks on lattice cryptography. The tree comparison counts quantum and classical node-oracle operations; total algorithms use quantum gates and classical bit operations including lattice arithmetic.
\end{problem}

\speedupheading
\complexity{Quantum}{$\widetilde O(\sqrt T\,\operatorname{poly}(n))$ node operations, quantum lattice enumeration.}
Here $T$ is the size of the classical enumeration tree for the chosen basis and search radius, and its depth is $n$. Radius refinement and witness recovery introduce polynomial factors. The arithmetic cost also depends on basis bit length and precision. This bound gives a matched acceleration of the specified tree, while general exact CVP can also be solved by running the classical $2^{n+o(n)}$ algorithm on quantum hardware.

\complexity{Classical}{$2^{n+o(n)}$ time and space, Aggarwal--Dadush--Stephens-Davidowitz exact CVP; $T$ nodes for the matched enumeration tree.}
The general exact algorithm uses discrete-Gaussian sampling and a reduction through clusters of approximate closest vectors. For the promise $\operatorname{dist}(t,L)\le\alpha\lambda_1(L)$ with $\alpha<1.039\ldots$, Hhan gives $2^{n/2+o(n)}$ classical time and space. Its lattice model and distance guarantee are part of that bound. The appropriate comparator is the minimum of the general, promised-distance and enumeration bounds applicable to the input.

\dimensionlabel{Strength}{strengthPolynomial} Quadratic acceleration in enumeration-tree size. The total-runtime comparison with the best classical solver depends on $T$, basis quality and the target-distance promise.

\dimensionlabel{Evidence}{evidencePartial} Rigorous quantum backtracking upper bounds; partial evidence for separation from optimized classical CVP algorithms on explicit lattices.

\dimensionlabel{Hardware}{hardwareExponential} Quantum enumeration needs polynomial logical registers for coefficient prefixes and reversible arithmetic, but potentially exponential coherent depth. Discrete-Gaussian and list-based methods use exponential storage. Large cryptanalytic instances require fault tolerance.


\subhead{Algorithm}
\begin{algorithmindex}{74}
\item \href{https://eprint.iacr.org/2018/546}{Aono, Nguyen and Shen, \emph{Quantum Lattice Enumeration and Tweaking Discrete Pruning} (ASIACRYPT 2018)}. Sections 3.1--3.2: close-vector enumeration; Algorithm 1 and Theorem 6: quantum tree minimization; Section 5.3, Algorithm 2: close-vector discrete pruning.
\item \href{https://arxiv.org/pdf/1504.01995}{Aggarwal, Dadush and Stephens-Davidowitz, \emph{Solving the Closest Vector Problem in $2^n$ Time---The Discrete Gaussian Strikes Again!} (FOCS 2015)}. Theorem 5.7: exact-CVP reduction; Theorem 6.4 and the resulting exact algorithm: $2^{n+o(n)}$ time.
\item \href{https://arxiv.org/pdf/2609.02764}{Hhan, \emph{Finding a Shortest Vector and More in $2^{n/2+o(n)}$ Time Using $q$-ary Coset Difference Tree} (2026)}. Theorem 1.1: $2^{n/2+o(n)}$ CVP bound with a distance promise; Section 5, Algorithm 4: CVP with a distance guarantee.
\end{algorithmindex}


\entrytitle{75}{Bounded distance decoding on periodic lattices}{white}{white}{hardwarePolynomial}
\setcounter{problem}{74}
\begin{problem}[Bounded distance decoding on periodic integer lattices]
Given an integer lattice $L\subseteq\mathbb Z^n$ with $q\mathbb Z^n\subseteq L$ and a target $t$ promised to satisfy $\operatorname{dist}(t,L)<\rho\lambda_1(L)$ for $\rho<1/2$, recover its unique closest lattice vector. Let $k$ be the minimum number of generators of the finite group $L/q\mathbb Z^n$. Complexity is measured in quantum gates and classical bit operations, polynomial in the lattice encoding and $\log q$.
\end{problem}

\speedupheading
\complexity{Quantum}{$\operatorname{poly}(n,\log q)$ time, Eldar--Hallgren periodic-lattice decoding.}
The decoding radius is $\lambda_1(L)2^{-\Omega(\sqrt{k\log q})}$, with the explicit constants and parameter restrictions of the theorem. The algorithm uses the finite-group structure of the $q$-periodic integer lattice and returns the closest vector under this promise. Basis encoding and the accuracy needed for quantum-state preparation enter the polynomial factors.

\complexity{Classical}{Deterministic polynomial time in the same radius regime, Ducas--van Woerden decoding as developed by Allen et al.}
The classical construction combines lattice reduction, duality, short full-rank vector sets and Babai nearest plane. The short-basis guarantee ensures correct rounding throughout the promised radius. Both bounds concern the same periodic-lattice parameters $n,k,q$ and exact recovery within the stated distance.

\dimensionlabel{Strength}{white} Polynomial quantum and classical complexity in the matched regime.

\dimensionlabel{Evidence}{white} Rigorous classical matching of the quantum decoding regime. The formal radius and input representation specify the scope of both polynomial-time statements.

\dimensionlabel{Hardware}{hardwarePolynomial} The quantum construction uses polynomial logical space, coherent algebraic arithmetic and Fourier processing. Its implementation would require controlled state preparation and fault-tolerant arithmetic at large dimension. The classical reduced-basis method implements the same decoding guarantee on ordinary hardware.


\subhead{Algorithm}
\begin{algorithmindex}{75}
\item \href{https://arxiv.org/pdf/2201.13450}{Eldar and Hallgren, \emph{An Efficient Quantum Algorithm for Lattice Problems Achieving Subexponential Approximation Factor} (2022)}. Algorithm 22, ``Algorithm Q'': bounded-distance decoding on $q$-periodic lattices; Algorithms 15, 17 and 19: quantum-state and sampling subroutines.
\item \href{https://arxiv.org/pdf/2203.05019}{Allen et al., \emph{Quantum and Classical Algorithms for Bounded Distance Decoding} (2022)}. Theorems 1.1--1.3: quantum promise and classical comparisons; the classical decoding development uses LLL and nearest plane.
\end{algorithmindex}


\entrytitle{76}{Learning with errors with classical samples}{strengthPolynomial}{evidencePartial}{hardwareExponential}
\setcounter{problem}{75}
\begin{problem}[Learning with errors with classical samples]
Given $A\in\mathbb Z_q^{m\times n}$ and $b=As+e\pmod q$, where $A$ is uniform, $s$ follows a specified secret distribution, and $e$ follows a specified small-noise distribution, recover $s$. Modulus, noise width, sample count, and secret distribution are essential parameters. ML-KEM and related schemes use structured versions of this problem. The comparison uses classical bit-operation time and quantum gate time with the specified coherent-memory model; the samples themselves are classical.
\end{problem}

\speedupheading
\complexity{Quantum}{$2^{0.2563\beta+o(\beta)}$ per heuristic SVP call, quantum sieving in lattice-reduction attacks.}
The block size $\beta$ is chosen by the primal, dual or hybrid attack for the given modulus, noise, dimension and sample count. Polynomial factors include reduction tours and matrix operations. Amplitude amplification can additionally accelerate guessed secret coordinates. The rigorous superlattice SVP alternative has $2^{0.247\beta+o(\beta)}$ time under the 2026 theorem; its use requires an attack reduction compatible with that subroutine.

\complexity{Classical}{The minimum of optimized primal, dual, hybrid, BKW and applicable bounded-error algorithms; heuristic sieving costs $2^{0.2925\beta+o(\beta)}$ per SVP call.}
For prime $q=n^{\kappa+o(1)}$, Lin--Vasudevan give expected time $2^{(1/2+o(1))n/\ln\ln q}$. One regime has errors bounded by $\exp((\ln n)^{o(1)})$ and at least $n+Cn/\ln\ln q$ samples. Another has secrets bounded by $\exp((\ln n)^{o(1)})$, errors $n^{o(1)}$ and at least $Cn/\ln\ln q$ samples, where $C>\min\{1,1/\kappa\}$. For fixed module ranks $k,\ell$ in power-of-two cyclotomic Module-LWE, Youmans gives $2^{(1/2+\xi)n/\ln\ln n}$ time and space for every fixed $\xi>0$, when $n=kd$, $m=(k+\ell)d$, $q=n^{\kappa+o(1)}$, $\kappa>1/2$, and matrix-independent errors have norm at most $\operatorname{polylog}(n)\sqrt m$ with probability $1-o(1)$. The algorithm uses the given $k+\ell$ ring samples.

\dimensionlabel{Strength}{strengthPolynomial} Polynomial speedup through a smaller exponential constant for the matched lattice-reduction subroutines. In the bounded-error regimes, the subexponential classical algorithms enter the minimum for the complete attack.

\dimensionlabel{Evidence}{evidencePartial} Partial evidence for quantum sieve speedups, with explicit random-vector and memory assumptions. The bounded-error classical statements are parameter-specific theorems; the module extension is a 2026 preprint.

\dimensionlabel{Hardware}{hardwareExponential} Fast quantum sieving requires exponential coherent storage in $\beta$ and fault-tolerant list access. Enumeration trades memory for runtime and uses polynomial logical qubits. Cryptographic resource estimates must optimize the full attack parameters and include reversible arithmetic.


\subhead{Algorithm}
\begin{algorithmindex}{76}
\item \href{https://arxiv.org/pdf/2205.14023}{Bonnetain et al., \emph{Finding Many Collisions via Reusable Quantum Walks: Application to Lattice Sieving} (2023)}. Section 7.1: quantum SVP subroutine used in lattice-reduction attacks.
\item \href{https://eprint.iacr.org/2015/1222}{Albrecht et al., \emph{On the Asymptotic Complexity of Solving LWE} (2015)}. Primal, dual and BKW attack analyses specify the parameter-dependent classical costs.
\item \href{https://arxiv.org/pdf/2610.13043}{Lin and Vasudevan, \emph{Provably Subexponential Algorithms for Bounded-Error LWE with Few Samples} (2026)}. Theorem 1.1: both parameter regimes; Algorithm 2 and Proposition 2.3: formal decision algorithm; Section 2.5 and Proposition 2.7: search and resource bounds.
\item \href{https://eprint.iacr.org/2026/2420}{Youmans, \emph{Provable Subexponential Algorithms for Module-SIS and Module-LWE} (2026)}. Abstract: fixed-rank Module-LWE theorem, error promise and sample count.
\end{algorithmindex}


\entrytitle{77}{Short integer solutions with a wide public matrix}{white}{white}{hardwarePolynomial}
\setcounter{problem}{76}
\begin{problem}[Short integer solutions with a wide public matrix]
For uniformly random $A\in\mathbb Z_q^{n\times m}$, output $0\ne z\in\mathbb Z^m$ with $Az=0\pmod q$ and $\|z\|_\infty\le\beta$. Chen--Liu--Zhandry study polynomial prime $q$, fixed constant $c$, $\beta=(q-c)/2$, and sufficiently large polynomial width, including $m=\Omega((q-c)^3 n^{c+1}q\log q)$. Complexity is measured in quantum gates and classical bit operations on the explicit public matrix, including its dimensions and modulus length.
\end{problem}

\speedupheading
\complexity{Quantum}{Polynomial time in the public dimensions and $q$, Chen--Liu--Zhandry quantum filtering.}
The guarantee holds for the stated wide random matrix, prime modulus and infinity-norm bound. Quantum filtering succeeds with the prescribed error-support and Fourier-coefficient conditions; preparation and processing of the LWE-like quantum states are included in the polynomial running time. The output is a classical nonzero short kernel vector.

\complexity{Classical}{$\operatorname{poly}(m,\log q)$ time in a stronger regime, Kothari--O'Donnell--Wu.}
The deterministic classical result covers all the Chen--Liu--Zhandry SIS infinity-norm parameter regimes. In particular, for arbitrary input matrices and $m\ge Cn^3$, it finds a nonzero $z$ satisfying $Az=0\pmod q$ and $\|z\|_\infty\le\lfloor q/6\rfloor$. Here $C$ is an appropriate absolute constant. This worst-case wide-matrix result strengthens both the input distribution and the norm guarantee of the matched quantum regime.

\dimensionlabel{Strength}{white} Polynomial quantum and classical solvability in the specified wide-matrix regime.

\dimensionlabel{Evidence}{white} Rigorous classical matching and stronger classical parameters. The matrix width and infinity-norm threshold define the algorithmic scope.

\dimensionlabel{Hardware}{hardwarePolynomial} The quantum filtering construction uses polynomial logical space in the full public-matrix dimensions, Fourier transforms and coherent filtering. Large instances require fault-tolerant state preparation and arithmetic. The deterministic classical method gives a direct ordinary-hardware implementation for the same regime.


\subhead{Algorithm}
\begin{algorithmindex}{77}
\item \href{https://arxiv.org/pdf/2108.11015}{Chen, Liu and Zhandry, \emph{Quantum Algorithms for Variants of Average-Case Lattice Problems via Filtering} (EUROCRYPT 2022)}. Theorem 2: SIS guarantee; quantum filtering algorithms and the LWE-state construction provide the associated subroutines.
\item \href{https://arxiv.org/pdf/2510.07515}{Kothari, O'Donnell and Wu, \emph{No Exponential Quantum Speedup for $\mathrm{SIS}_\infty$ Anymore} (STOC 2026)}. Theorem 1.5, proved in Section 4.4: classical infinity-norm algorithm; its $k=3$ specialization gives the $q/6$ bound with cubic width.
\end{algorithmindex}


\entrytitle{78}{Binary syndrome decoding}{strengthPolynomial}{evidencePartial}{hardwareExponential}
\setcounter{problem}{77}
\begin{problem}[Binary syndrome decoding]
Given a full-rank binary parity-check matrix $H\in\mathbb F_2^{(n-k)\times n}$, a syndrome $s$, and weight $w$, find $e\in\mathbb F_2^n$ with $He=s$ and $|e|=w$. Random instances model generic attacks on code-based cryptography. Complexity is measured in quantum gate time and classical bit-operation time, with coherent list-access costs specified for quantum walks.
\end{problem}

\speedupheading
\complexity{Quantum}{$\operatorname{poly}(n)\sqrt{\binom nw/\binom{n-k}{w}}$, amplified Prange; parameter-dependent improvements from quantum-walk ISD.}
For $w\le n-k$, $p=\binom{n-k}{w}/\binom nw$ is the success probability of one Prange information-set trial in the usual random-code model. Reversible elimination, the weight test and uncomputation enter its polynomial cost. Kachigar--Tillich and Kirshanova accelerate internal list matching as well as the outer search. The $0.05869$ exponent is the maximum over rates at Gilbert--Varshamov weight in the Kachigar--Tillich comparison; concrete attacks optimize the actual $n,k,w$.

\complexity{Classical}{The minimum of optimized BJMM/May--Ozerov ISD and applicable sparse-circuit counting; Prange costs $\operatorname{poly}(n)/p$.}
For dense random constant-rate codes, the classical comparator minimizes the parameter-dependent representation and nearest-neighbor costs of modern ISD. If $H$ has $h$ nonzero entries, the predicate $He=s$, $|e|=w$ has an $O(n+h)$-gate Boolean circuit. Tensor contraction and witness recovery give deterministic time $\operatorname{poly}(n+h)2^{O((n+h)/\sqrt{\log(n+h)})}$. For $h=O(n)$ this is $2^{O(n/\sqrt{\log n})}$. Dense matrices have $h=\Theta(n^2)$, where optimized ISD supplies the stronger displayed bound.

\dimensionlabel{Strength}{strengthPolynomial} Polynomial speedup for dense random codes: quadratic acceleration of Prange and further reductions in the exponential constant from quantum walks. Sparse parity checks have the stated subexponential classical comparator.

\dimensionlabel{Evidence}{evidencePartial} Partial evidence for a total-runtime separation on explicit dense codes; amplification rigorously improves a specified trial success probability.

\dimensionlabel{Hardware}{hardwareExponential} Reversible elimination uses polynomial logical space. Advanced walks trade runtime for exponentially large coherent list storage. Cryptographic instances require fault-tolerant arithmetic and many coherent trial evaluations.

{\small [Before ``ETH is false'', the classical sparse-code comparison used exponential-time ISD, and quantum amplification reduced a Prange trial cost from $\operatorname{poly}(n)/p$ to $\operatorname{poly}(n)/\sqrt p$.]\par}
\subhead{Algorithm}
\begin{algorithmindex}{78}
\item \href{https://arxiv.org/pdf/1703.00263}{Kachigar and Tillich, \emph{Quantum Information Set Decoding Algorithms} (2017)}. Algorithm 2: ISD skeleton; quantum-walk analyses and rate optimization give the $0.05869$ exponent.
\item \href{https://arxiv.org/pdf/1808.00714}{Kirshanova, \emph{Improved Quantum Information Set Decoding} (2018)}. Algorithm 1 and Theorem 1: quantum ISD and its runtime formula.
\item \href{https://arxiv.org/pdf/2112.06157}{Esser et al., \emph{An Optimized Quantum Implementation of ISD on Scalable Quantum Resources} (2021)}. Algorithms 1--4: Prange, state preparation, linear-system solving and the combined quantum circuit.
\item \href{https://github.com/openai/math/blob/main/preprints/The-Exponential-Time-Hypothesis-is-False-October-25-2026/paper.pdf}{OpenAI, \emph{The Exponential Time Hypothesis Is False} (2026)}. Theorem 1.1: Boolean tensor contraction; witness self-reduction following the theorem and the bounded-width reductions in Appendix D: explicit-circuit search bound.
\end{algorithmindex}


\entrytitle{79}{Minimum-weight nonzero codewords}{strengthPolynomial}{evidencePartial}{hardwareExponential}
\setcounter{problem}{78}
\begin{problem}[Minimum-weight nonzero codewords]
Given a full-rank generator matrix $G\in\mathbb F_2^{k\times n}$, output $u\ne0$ minimizing $|uG|$. The minimum value is the code's minimum distance. This is a discrete optimization problem related to syndrome decoding and the search for short dual codewords in statistical cryptanalysis. Complexity is measured in weight-oracle evaluations for the direct search, and quantum gates or classical bit operations for explicitly represented codes.
\end{problem}

\speedupheading
\complexity{Quantum}{$O(2^{k/2})$ weight evaluations, D\"urr--H\o yer minimum finding; $C/\sqrt{p_w}$ for amplified ISD trials.}
The direct exact method has polynomial workspace and includes coherent computation of $uG$, its Hamming weight and comparison with a threshold. For random codes, let an ISD trial cost $C$ and find a nonzero word of weight at most $w$ with probability $p_w$. Amplification replaces the trial count by $O(p_w^{-1/2})$, and quantum nearest-neighbor walks can reduce the internal cost. Full attacks optimize weight, representations and list sizes.

\complexity{Classical}{The minimum of $O(2^k)$ weight evaluations, optimized ISD and applicable sparse-circuit counting.}
Exhaustive enumeration can update codewords in Gray-code order. An ISD trial family gives expected work $C/p_w$, with optimized representations and nearest-neighbor matching improving the random-code comparator. If a generator or parity-check matrix has $s$ nonzero entries, the threshold predicate has $O(n+s)$ Boolean gates, yielding $\operatorname{poly}(n+s)2^{O((n+s)/\sqrt{\log(n+s)})}$ deterministic classical time. Binary search over the weight adds a polynomial factor. A linear-size representation gives subexponential time in $n$; dense constant-rate codes retain the stronger enumeration or ISD bound.

\dimensionlabel{Strength}{strengthPolynomial} Polynomial speedup from quadratic acceleration of exhaustive weight search or a specified ISD success probability, giving a smaller exponential constant in the dense-code comparison.

\dimensionlabel{Evidence}{evidencePartial} Rigorous oracle amplification; partial evidence for a full-runtime separation on explicit random codes. Sparse representations have the separate subexponential classical bound.

\dimensionlabel{Hardware}{hardwareExponential} The direct circuit uses $k$ message qubits, $n$ codeword qubits, a weight register and polynomial ancillas. Exponentially many fault-tolerant repetitions dominate its cost. Faster quantum-walk variants use exponential coherent lists.

{\small [Before ``ETH is false'', the generic exact comparison was $O(2^k)$ classical weight evaluations versus $O(2^{k/2})$ quantum evaluations, a quadratic speedup also applied to sparse representations.]\par}
\subhead{Algorithm}
\begin{algorithmindex}{79}
\item \href{https://arxiv.org/pdf/quant-ph/9607014}{D\"urr and H\o yer, \emph{A Quantum Algorithm for Finding the Minimum} (1996)}. The minimum-finding algorithm and its expected $O(\sqrt N)$ query theorem.
\item \href{https://arxiv.org/pdf/1703.00263}{Kachigar and Tillich, \emph{Quantum Information Set Decoding Algorithms} (2017)}. Algorithm 2 and the quantum-walk cost analyses, applied at zero syndrome with the zero vector excluded.
\item \href{https://arxiv.org/pdf/1808.00714}{Kirshanova, \emph{Improved Quantum Information Set Decoding} (2018)}. Algorithm 1 and Theorem 1: the parameter-dependent ISD construction.
\item \href{https://github.com/openai/math/blob/main/preprints/The-Exponential-Time-Hypothesis-is-False-October-25-2026/paper.pdf}{OpenAI, \emph{The Exponential Time Hypothesis Is False} (2026)}. Theorem 1.1: Boolean tensor contraction; witness self-reduction following the theorem and the bounded-width reductions in Appendix D: explicit-circuit search bound.
\end{algorithmindex}


\entrytitle{80}{Random subset sum at cryptographic density}{strengthPolynomial}{evidencePartial}{hardwareExponential}
\setcounter{problem}{79}
\begin{problem}[Random subset sum at cryptographic density]
Given independent random integers $a_1,\ldots,a_n$ modulo a modulus $M$ of roughly $n$ bits and a target $t$, find $x\in\{0,1\}^n$ such that $\sum_i a_ix_i=t\pmod M$. Density and the distribution of solutions are fixed as in the representation-algorithm model. These instances appear in knapsack cryptanalysis and as subproblems in decoding. Complexity is measured in classical RAM time and quantum time with either quantum access to classical memory or quantum-accessible quantum memory, as stated.
\end{problem}

\speedupheading
\complexity{Quantum}{$\widetilde O(2^{0.2182n})$ time with classical list heuristics, or $\widetilde O(2^{0.2156n})$ with an additional update heuristic, Bonnetain--Bricout--Schrottenloher--Shen quantum walks.}
The first walk has memory exponent $0.2182$; the second has memory exponent $0.2110$. A quantum-search variant uses $\widetilde O(2^{0.2356n})$ time and classically stored lists with quantum random access. The walk variants use quantum-accessible quantum memory. Their representation lists, modular filters and update costs are analyzed for the random density-one distribution in the problem statement.

\complexity{Classical}{$\widetilde O(2^{0.283n})$ time and space, Bonnetain--Bricout--Schrottenloher--Shen classical representations.}
The result uses the standard random-instance list heuristics and representations with coefficients in $\{-1,0,1,2\}$. For positive binary integers of polynomial bit length, the general worst-case subset-sum algorithm gives randomized $O(2^{0.49n})$ word-RAM time; the $0.283$ exponent is the stronger bound in this random regime. A direct Boolean verifier for $n$ modular summands of $b$ bits has $O(n(b+\log n))$ gates. With $b=\Theta(n)$, applying the explicit-circuit counting bound gives an exponent of order $n^2/\sqrt{\log n}$, so the representation bound determines the displayed comparison.

\dimensionlabel{Strength}{strengthPolynomial} Polynomial speedup through a smaller exponential constant in $n$. Both amplitude amplification and quantum walks contribute to the algorithms.

\dimensionlabel{Evidence}{evidencePartial} Partial evidence under the stated random-instance and list assumptions. The $0.2156$ result additionally assumes the quantum-walk update heuristic; the $0.2182$ result uses the classical heuristics.

\dimensionlabel{Hardware}{hardwareExponential} All fast variants are memory-intensive fault-tolerant workloads. The $0.2356$ algorithm uses exponential classical storage with coherent access; the faster walks use exponential quantum-accessible quantum memory. The memory-access model determines the gate cost of a lookup.


\subhead{Algorithm}
\begin{algorithmindex}{80}
\item \href{https://arxiv.org/pdf/2002.05276}{Bonnetain, Bricout, Schrottenloher and Shen, \emph{Improved Classical and Quantum Algorithms for Subset-Sum} (ASIACRYPT 2020)}. Section 2.5: classical $0.283$ exponent; Section 4.3: quantum-search $0.2356$ exponent; Section 5.3: heuristic quantum walk; Section 6.4: walk without the additional update heuristic; Table 1: models and resource exponents.
\item \href{https://github.com/openai/math/tree/main/preprints/Subset-Sum-in-Time-2-power-0-49n-October-4-2026}{OpenAI, \emph{Subset Sum in Time $2^{0.49n}$} (2026)}. Main randomized word-RAM theorem for worst-case inputs.
\end{algorithmindex}


\entrytitle{81}{Boolean multivariate quadratic equations}{strengthPolynomial}{evidencePartial}{hardwareExponential}
\setcounter{problem}{80}
\begin{problem}[Boolean multivariate quadratic equations]
Given $m$ quadratic polynomials $f_1,\ldots,f_m\in\mathbb F_2[x_1,\ldots,x_n]$, find $x\in\mathbb F_2^n$ such that $f_i(x)=0$ for every $i$, or report infeasibility. The square case $m=n$ is a basic algebraic-cryptanalysis benchmark and a foundation for multivariate signature assumptions. Complexity is measured in elementary quantum gates and classical bit-operation time on the explicit polynomial coefficient list.
\end{problem}

\speedupheading
\complexity{Quantum}{$O(mn^2 2^{n/2})$ gates, Grover search with reversible quadratic-equation evaluation.}
Schwabe--Westerbaan construct the equation-checking oracle with $O(mn^2)$ gates and $O(m+n)$ qubits; their space-efficient variant uses $O(n+\log m)$ qubits with the same asymptotic gate bound. QuantumBooleanSolve gives $O(2^{(0.462+\varepsilon)n})$ expected gates for every fixed $\varepsilon>0$ when $m=n$ and the system is $0.17$-strong semiregular. Its deterministic variant gives $O(2^{(0.47+\varepsilon)n})$ gates for $m=n$ under $0.13$-strong semiregularity. These bounds include Macaulay-system linear algebra and coherent testing of partial assignments.

\complexity{Classical}{$\operatorname{poly}(m,n)+O(2^{0.6943n})$ time, Dinur multiple parity counting; subexponential time for sparse systems, Boolean tensor contraction.}
Polynomial-time Gaussian elimination removes linearly dependent equations from the coefficient list; Dinur's randomized search algorithm then uses parity counts of partial assignments. If $m=O(n)$ and each equation has polylogarithmically many monomials, the counting reduction gives $\operatorname{poly}(n)2^{O(n\log\log n/\sqrt{\log n})}$. With linear total variable--monomial incidence, an $O(n)$-size Boolean circuit gives $\operatorname{poly}(n)2^{O(n/\sqrt{\log n})}$. These sparse bounds take precedence for their input regimes. Dense systems satisfying the quantum semiregularity promise give the displayed $0.6943$ versus $0.462$ comparison.

\dimensionlabel{Strength}{strengthPolynomial} Polynomial speedup through a smaller exponential constant for the dense promised regime, using amplitude amplification and algebraic preprocessing. Sparse explicit systems have subexponential classical bounds.

\dimensionlabel{Evidence}{evidencePartial} Partial evidence: the quantum algebraic exponent uses strong semiregularity, while the classical parity-counting bound is a general upper bound.

\dimensionlabel{Hardware}{hardwareExponential} Reversible polynomial evaluation and Macaulay linear algebra need fault tolerance and substantial workspace. Direct Grover search uses $n$ assignment qubits and polynomial ancillas; the algebraic algorithm has additional matrix-storage and reversible-linear-algebra costs.

{\small [Before ``ETH is false'', the general classical upper bound applied to sparse quadratic systems was $O(2^{0.6943n})$, compared with the claimed $O(2^{(0.462+\varepsilon)n})$ quantum bound under strong semiregularity.]\par}
\subhead{Algorithm}
\begin{algorithmindex}{81}
\item \href{https://cryptojedi.org/papers/mqgrover-20171130.pdf}{Schwabe and Westerbaan, \emph{Solving Binary MQ with Grover's Algorithm} (SPACE 2016; full version 2017)}. Section 2.1: problem definition; Sections 4--5: the two reversible equation-checking circuits and their gate and qubit costs.
\item \href{https://arxiv.org/pdf/1712.07211}{Faug\`ere et al., \emph{Fast Quantum Algorithm for Solving Multivariate Quadratic Equations} (2017)}. Theorems 5--7 and Section 4: quantum circuit and QuantumBooleanSolve complexity; Theorem 7 states the $0.17$-strong semiregular $0.462+\varepsilon$ bound.
\item \href{https://arxiv.org/pdf/2005.04800}{Dinur, \emph{Improved Algorithms for Solving Polynomial Systems over GF(2) by Multiple Parity-Counting} (2020)}. Theorem 1.1: search complexity; Section 2.2: elimination of linearly dependent equations; Algorithms 1--3: parity-counting procedures.
\item \href{https://github.com/openai/math/blob/main/preprints/The-Exponential-Time-Hypothesis-is-False-October-25-2026/paper.pdf}{OpenAI, \emph{The Exponential Time Hypothesis Is False} (2026)}. Theorem 1.1: Boolean tensor contraction; witness self-reduction following the theorem and the bounded-width reductions in Appendix D: explicit-circuit search bound.
\end{algorithmindex}


\entrytitle{82}{Symmetric-key recovery in the ideal-cipher model}{strengthPolynomial}{evidenceProved}{hardwareExponential}
\setcounter{problem}{81}
\begin{problem}[Symmetric-key recovery in the ideal-cipher model]
Given plaintext--ciphertext pairs $(P_i,C_i)$ generated by a block cipher $E_K$ with a uniformly random $k$-bit key, find a key satisfying every equation $E_K(P_i)=C_i$. Choose sufficiently many pairs that the key is unique with high probability. The ideal-cipher comparison counts cipher-oracle queries; explicit implementations use quantum gate complexity and classical bit complexity.
\end{problem}

\speedupheading
\complexity{Quantum}{$\Theta(2^{k/2})$ cipher queries, Grover key search.}
With sufficiently many known pairs to isolate the key, a reversible verification circuit of $G$ gates and workspace $W$ gives $O(2^{k/2}(G+k))$ gates and $k+W$ logical qubits, up to the registers needed for the supplied pairs. The cost includes verification and uncomputation. Constant success probability is obtained after the usual number of amplitude-amplification iterations.

\complexity{Classical}{$\Theta(2^k)$ ideal-cipher queries, exhaustive key search; explicit small circuits admit Boolean tensor contraction.}
The query bound is tight for a uniformly random key in the ideal-cipher model. If the cipher is given by an explicit verification circuit with $G$ gates, its key-recovery predicate has $O(k+G)$-size 3-CNF encoding. Tensor contraction and witness recovery give $\operatorname{poly}(k+G)2^{O((k+G)/\sqrt{\log(k+G)})}$ bit complexity. For $G=O(k)$ this is subexponential in $k$. The ideal-cipher query lower bound and the explicit-circuit upper bound refer to their respective input representations.

\dimensionlabel{Strength}{strengthPolynomial} Tight quadratic polynomial speedup in the ideal-cipher key-space size, entirely from Grover search. For $2^k$ keys it halves the exponential constant; explicit linear-size verification circuits have the stated subexponential classical comparator.

\dimensionlabel{Evidence}{evidenceProved} Fully proved in the ideal-cipher query model. For an explicit cipher, the comparison includes its best structural and circuit-based attacks.

\dimensionlabel{Hardware}{hardwareExponential} Cryptographic search requires fault tolerance and exponential coherent depth. The Grassl--Langenberg--Roetteler--Steinwandt AES-128 construction uses 2,953 logical qubits and roughly $2^{87}$ gates for full key search. A logical circuit estimate includes reversible rounds, key schedule, multi-pair checks and diffusion.

{\small [Before ``ETH is false'', the generic explicit-key-search baseline was $O(2^k)$ classical trials versus $O(2^{k/2})$ quantum trials, a quadratic speedup.]\par}
\subhead{Algorithm}
\begin{algorithmindex}{82}
\item \href{https://arxiv.org/pdf/quant-ph/9605043}{Grover, \emph{A Fast Quantum Mechanical Algorithm for Database Search} (1996)}. Section 3: the search algorithm.
\item \href{https://arxiv.org/pdf/1512.04965}{Grassl, Langenberg, Roetteler and Steinwandt, \emph{Applying Grover's Algorithm to AES: Quantum Resource Estimates} (2016)}. Section 3.1: key uniqueness; Section 3.2: AES circuit; Section 3.4 and Table 5: full search resources.
\item \href{https://github.com/openai/math/blob/main/preprints/The-Exponential-Time-Hypothesis-is-False-October-25-2026/paper.pdf}{OpenAI, \emph{The Exponential Time Hypothesis Is False} (2026)}. Theorem 1.1 and witness self-reduction: search for a satisfying assignment of the explicit Boolean verification circuit.
\end{algorithmindex}


\entrytitle{83}{Hash preimages and second preimages}{strengthPolynomial}{evidenceProved}{hardwareExponential}
\setcounter{problem}{82}
\begin{problem}[Hash preimages and second preimages]
Given oracle access to $H:\{0,1\}^{\ell}\to\{0,1\}^n$ and a target $y$ in its image, find $x$ satisfying $H(x)=y$. For second-preimage search, a particular $x_0$ is supplied, $y=H(x_0)$, and the output must satisfy $x\ne x_0$. Assume a random-function model and a search domain large enough to contain suitable preimages. The random-function comparison counts hash-oracle queries; an explicit hash circuit adds quantum gate cost and admits classical bit-complexity bounds.
\end{problem}

\speedupheading
\complexity{Quantum}{$\Theta(\sqrt{D/M})$ queries, Grover search and amplitude amplification.}
Here $M\ge1$ is the number of valid inputs in the domain of size $D=2^\ell$, and success probability is constant. Random $n$-bit outputs give the familiar $\Theta(2^{n/2})$ scale when the marked fraction has order $2^{-n}$. The excluded input in second-preimage search is incorporated in the predicate. Unknown $M$ can be handled by randomized iteration counts. If reversible hashing and checking cost $G$ gates, the gate bound is $O((G+\ell)\sqrt{D/M})$.

\complexity{Classical}{$\Theta(D/M)$ queries, random sampling.}
These bounds match the random-oracle lower bound at constant success probability. For an explicit hash circuit with $G$ gates, the preimage predicate on $\ell$ input bits has $O(\ell+G)$-size Boolean encoding. Tensor contraction gives $\operatorname{poly}(\ell+G)2^{O((\ell+G)/\sqrt{\log(\ell+G)})}$ deterministic search time. When $\ell+G=O(n)$, that explicit-family bound is subexponential in the output length.

\dimensionlabel{Strength}{strengthPolynomial} Tight quadratic polynomial speedup in the oracle search-space ratio, arising from amplitude amplification. In the random-output parameter $n$, the exponent changes from $n$ to $n/2$ without power saving.

\dimensionlabel{Evidence}{evidenceProved} Fully proved for the oracle task with fixed marked fraction. Explicit hash functions are compared using their representation-dependent algorithms as well.

\dimensionlabel{Hardware}{hardwareExponential} The candidate register needs $\ell$ logical qubits plus reversible-hash workspace and the equality test. Full cryptographic searches require exponentially many coherent evaluations and fault tolerance. Parallel search trades elapsed time for independently maintained quantum processors.

{\small [Before ``ETH is false'', the generic preimage comparison used $\Theta(2^n)$ classical versus $\Theta(2^{n/2})$ quantum queries, with the same trial bounds transferred to explicit hash evaluation.]\par}
\subhead{Algorithm}
\begin{algorithmindex}{83}
\item \href{https://arxiv.org/pdf/quant-ph/9605034}{Boyer, Brassard, H\o yer and Tapp, \emph{Tight Bounds on Quantum Searching} (1998)}. Search procedure with an unknown number of marked items and its expected-query analysis.
\item \href{https://arxiv.org/pdf/quant-ph/0005055}{Brassard, H\o yer, Mosca and Tapp, \emph{Quantum Amplitude Amplification and Estimation} (2002)}. Section 2, Theorems 2--3: known and unknown success probabilities.
\item \href{https://github.com/openai/math/blob/main/preprints/The-Exponential-Time-Hypothesis-is-False-October-25-2026/paper.pdf}{OpenAI, \emph{The Exponential Time Hypothesis Is False} (2026)}. Theorem 1.1 and witness self-reduction: search for a satisfying assignment of the explicit Boolean verification circuit.
\end{algorithmindex}


\entrytitle{84}{Hash collisions}{strengthPolynomial}{evidenceProved}{hardwareExponential}
\setcounter{problem}{83}
\begin{problem}[Hash collisions]
Given a random function $H$ with $n$-bit outputs and a sufficiently large domain, find distinct $x,x'$ with $H(x)=H(x')$. Collision resistance is the relevant property for many hash-based constructions and commitment schemes. The random-function comparison measures hash-oracle queries; total quantum time additionally counts coherent table access, and explicit circuits use gate or bit complexity.
\end{problem}

\speedupheading
\complexity{Quantum}{$O(r+\sqrt{2^n/r})$ queries with a table of $r$ values, Brassard--H\o yer--Tapp collision search.}
The table parameter satisfies $1\le r\le 2^n$. Optimizing it gives $r=2^{n/3}$ and the tight $\Theta(2^{n/3})$ quantum query bound. Near-matching time assumes efficient coherent lookup into the stored table. A gate implementation includes the lookup circuit and reversible hash cost. The random-function setting yields the matching quantum collision lower-bound scale.

\complexity{Classical}{$\Theta(2^{n/2})$ queries, birthday search or Pollard-rho collision finding.}
Birthday storage is of order $2^{n/2}$ entries; Pollard-rho offers low-memory collision finding under the usual random-function assumptions. For a hash given by a $G$-gate circuit on $\ell$ input bits, the predicate $x\ne x'$, $H(x)=H(x')$ has $O(\ell+n+G)$ gates. Boolean tensor contraction and witness recovery give $\operatorname{poly}(\ell+n+G)2^{O((\ell+n+G)/\sqrt{\log(\ell+n+G)})}$ classical time. A linear-size explicit family therefore has a subexponential comparator.

\dimensionlabel{Strength}{strengthPolynomial} Polynomial speedup of degree $3/2$ relative to the classical birthday query count. The BHT construction combines a stored table with Grover search. In output length $n$, the exponent changes from $1/2$ to $1/3$ without power saving.

\dimensionlabel{Evidence}{evidenceProved} Fully proved in the collision-oracle model with the stated domain promise; random-function collision bounds give the corresponding cryptographic query scale.

\dimensionlabel{Hardware}{hardwareExponential} The time-efficient BHT implementation stores about $2^{n/3}$ entries with coherent access and requires exponentially many fault-tolerant operations. Polynomial-qubit alternatives use different time-memory tradeoffs; the lookup model is a material resource assumption.

{\small [Before ``ETH is false'', the generic collision comparison used $\Theta(2^{n/2})$ classical versus $\Theta(2^{n/3})$ quantum queries, an exponential-constant improvement also used for explicit hash circuits.]\par}
\subhead{Algorithm}
\begin{algorithmindex}{84}
\item \href{https://arxiv.org/pdf/quant-ph/9705002}{Brassard, H\o yer and Tapp, \emph{Quantum Algorithm for the Collision Problem} (1997)}. Theorem 1 and \textsc{Collision}$(F,k)$: table-based quantum collision search; Section 3: coherent table-access costs.
\item \href{https://arxiv.org/pdf/quant-ph/0112086}{Shi, \emph{Quantum Lower Bounds for the Collision and the Element Distinctness Problems} (FOCS 2002)}. Theorem 1.1: tight cube-root collision query lower bound.
\item \href{https://github.com/openai/math/blob/main/preprints/The-Exponential-Time-Hypothesis-is-False-October-25-2026/paper.pdf}{OpenAI, \emph{The Exponential Time Hypothesis Is False} (2026)}. Theorem 1.1 and witness self-reduction: search for a satisfying assignment of the explicit Boolean verification circuit.
\end{algorithmindex}


\entrytitle{85}{Claw finding between two functions}{strengthPolynomial}{evidenceProved}{hardwareExponential}
\setcounter{problem}{84}
\begin{problem}[Claw finding between two functions]
Given evaluation oracles for $f:[N]\to Z$ and $g:[M]\to Z$, with $N\le M$, find $(x,y)$ with $f(x)=g(y)$ or determine that none exists. Claws arise when matching two halves of a cryptanalytic computation, including double encryption and prescribed-degree isogeny search. Complexity is measured in evaluation-oracle queries; runtime includes list-access costs, and explicit functions are additionally compared in gate and bit complexity.
\end{problem}

\speedupheading
\complexity{Quantum}{$O((NM)^{1/3})$ queries for $N\le M<N^2$, and $O(\sqrt M)$ for $M\ge N^2$, Tani quantum-walk claw finding.}
The balanced case is $\Theta(N^{2/3})$. The algorithm stores evaluated subsets of both domains; coherent sorting or lookup contributes to its gate cost. The stated bounded-error bounds include finding a claw or deciding that none exists, for the evaluation-oracle model. Quantum-walk spectral and marked-fraction parameters determine the two regimes.

\complexity{Classical}{$\Theta(N+M)$ queries in the worst case, classical list matching.}
Hashing both lists gives expected linear-time matching with suitable value representation; sorting gives a logarithmic time overhead. For $N=M$, the bound is $\Theta(N)$. If $f$ and $g$ have explicit verification circuits of combined size $G$ and the combined witness length is $\ell=\lceil\log_2N\rceil+\lceil\log_2M\rceil$, circuit encoding gives $\operatorname{poly}(\ell+G)2^{O((\ell+G)/\sqrt{\log(\ell+G)})}$ classical search. In the balanced case with $G=O(\log N)$, this becomes $N^{o(1)}$.

\dimensionlabel{Strength}{strengthPolynomial} Polynomial speedup, with degree $3/2$ in the balanced query comparison and quadratic search in the highly unbalanced regime. Quantum walks provide the balanced improvement.

\dimensionlabel{Evidence}{evidenceProved} Fully proved for the evaluation-oracle task and its matching query lower bounds. Explicit small circuits have the separate representation-dependent classical comparator.

\dimensionlabel{Hardware}{hardwareExponential} Time-efficient walks need coherent access to stored function values and exponential storage in the logarithmic domain size. Reversible function evaluation and walk updates require fault tolerance for cryptographic domains.

{\small [Before ``ETH is false'', the balanced generic comparison was $\Theta(N)$ classical versus $\Theta(N^{2/3})$ quantum queries, also transferred to explicit function evaluations.]\par}
\subhead{Algorithm}
\begin{algorithmindex}{85}
\item \href{https://arxiv.org/pdf/0708.2584}{Tani, \emph{Claw Finding Algorithms Using Quantum Walk} (2009)}. Section 3: claw detection; Section 4, \textsc{ClawSearch} and Theorem 8: finding a claw and the two domain-size regimes.
\item \href{https://github.com/openai/math/blob/main/preprints/The-Exponential-Time-Hypothesis-is-False-October-25-2026/paper.pdf}{OpenAI, \emph{The Exponential Time Hypothesis Is False} (2026)}. Theorem 1.1 and witness self-reduction: search for a satisfying assignment of the explicit Boolean verification circuit.
\end{algorithmindex}


\entrytitle{86}{Double-encryption key recovery}{strengthPolynomial}{evidenceProved}{hardwareExponential}
\setcounter{problem}{85}
\begin{problem}[Double and multiple encryption key recovery]
For independent $k$-bit keys $K_1,K_2$, given sufficiently many known pairs satisfying $C=E_{K_2}(E_{K_1}(P))$, recover the key pair. The ideal-cipher model provides a precise comparison for security amplification by repeated encryption. The ideal-cipher comparison counts encryption and decryption queries; explicit cipher implementations also use quantum gate and classical bit complexity.
\end{problem}

\speedupheading
\complexity{Quantum}{$\Theta(2^{2k/3})$ ideal-cipher queries, quantum meet-in-the-middle through claw finding.}
The forward and backward key domains each have size $N=2^k$, giving the optimal $\Theta(N^{2/3})$ query bound. Tuples of intermediate values from enough known pairs isolate the intended key pair. A quantum walk maintains subsets of forward and backward evaluations; its stored-value access and reversible cipher operations determine the full gate cost. The associated time-efficient construction uses $O(N^{2/3})$ memory.

\complexity{Classical}{$\widetilde O(2^k)$ time and $O(2^k)$ storage, classical meet-in-the-middle.}
Matching forward encryptions against backward decryptions recovers the two keys; the ideal-cipher query scale is $\Theta(2^k)$. For a known explicit cipher, let $G$ be the gate count of the two-key verification circuit, including the supplied pairs. Its $2k$ witness bits give classical time $\operatorname{poly}(k+G)2^{O((k+G)/\sqrt{\log(k+G)})}$ by Boolean tensor contraction. Linear-size verification circuits have a subexponential classical bound in $k$.

\dimensionlabel{Strength}{strengthPolynomial} Polynomial speedup of degree $3/2$ in the ideal-cipher meet-in-the-middle comparison, or an exponential-constant improvement from $k$ to $2k/3$. The mechanism is quantum claw finding.

\dimensionlabel{Evidence}{evidenceProved} Fully proved in the stated ideal-cipher query model through matching upper and adversary lower bounds.

\dimensionlabel{Hardware}{hardwareExponential} The time-efficient algorithm requires exponentially large coherent lists and exponentially many reversible cipher evaluations. Logical workspace for a single encryption is polynomial in the cipher size, while coherent storage and circuit depth dominate a cryptographic-scale fault-tolerant attack.

{\small [Before ``ETH is false'', the explicit double-encryption comparison used classical $\widetilde O(2^k)$ meet-in-the-middle versus $\widetilde O(2^{2k/3})$ quantum claw finding.]\par}
\subhead{Algorithm}
\begin{algorithmindex}{86}
\item \href{https://arxiv.org/pdf/1410.1434}{Kaplan, \emph{Quantum Attacks Against Iterated Block Ciphers} (2015)}. Theorem 2: double-encryption algorithm and memory; Theorem 3: matching query lower bound; Theorem 4: the related fourfold-encryption construction.
\item \href{https://arxiv.org/pdf/0708.2584}{Tani, \emph{Claw Finding Algorithms Using Quantum Walk} (2009)}. Section 4 and Theorem 8: the balanced claw-finding subroutine.
\item \href{https://github.com/openai/math/blob/main/preprints/The-Exponential-Time-Hypothesis-is-False-October-25-2026/paper.pdf}{OpenAI, \emph{The Exponential Time Hypothesis Is False} (2026)}. Theorem 1.1 and witness self-reduction: search for a satisfying assignment of the explicit Boolean verification circuit.
\end{algorithmindex}


\entrytitle{87}{Generalized birthday search: random \texorpdfstring{$k$}{k}-XOR}{strengthPolynomial}{evidencePartial}{hardwareExponential}
\setcounter{problem}{86}
\begin{problem}[Generalized birthday search: random $k$-XOR]
Given $k$ lists $L_1,\ldots,L_k$ of uniformly random $n$-bit strings, find $x_i\in L_i$ with $x_1\oplus\cdots\oplus x_k=0$. The list lengths are part of the input. The single-solution regime uses lists of size about $2^{n/k}$; unrestricted-list variants permit larger lists. These problems model hash combiners, subset sums, and parts of decoding. Complexity is measured in classical list-processing time and quantum time or oracle queries, with coherent access to the prescribed random lists.
\end{problem}

\speedupheading
\complexity{Quantum}{$\widetilde O(2^{n/(2+\lfloor\log_2k\rfloor)})$ time for unrestricted lists, quantum generalized birthday algorithms.}
The list sizes in this setting can be chosen to support the merging strategy. For four prescribed random lists of size $2^{n/4}$, Schrottenloher gives $\widetilde O(2^{7n/24})$ quantum time in the single-solution regime. Its asymmetric merges, quantum search and quantum walks exploit that specific list-size promise. The fastest time bounds depend on the available memory and coherent lookup model; random-oracle $k$-XOR query lower bounds have scale $2^{n/(k+1)}$.

\complexity{Classical}{$\widetilde O(2^{n/(1+\lfloor\log_2k\rfloor)})$ time with unrestricted lists, Wagner generalized birthday search.}
Wagner uses lists of the corresponding size and successive prefix constraints. For four prescribed lists of size $2^{n/4}$, pair-list meet-in-the-middle costs $\widetilde O(2^{n/2})$ time. This is the comparator for the $7n/24$ quantum result. The unrestricted-list and prescribed-list problems have different input promises and are compared separately.

\dimensionlabel{Strength}{strengthPolynomial} Polynomial speedups through smaller exponential constants in $n$ for fixed $k$. Amplitude amplification and quantum walks both appear in the merging constructions.

\dimensionlabel{Evidence}{evidencePartial} Partial evidence for time optimality, together with rigorous random-oracle query bounds. List size, number of solutions and memory model specify each comparison.

\dimensionlabel{Hardware}{hardwareExponential} The fastest methods require exponentially large lists with coherent access and fault-tolerant merging operations. Polynomial-qubit variants trade greater runtime for less coherent memory. Explicitly supplied lists also incur their stated loading and storage costs.


\subhead{Algorithm}
\begin{algorithmindex}{87}
\item \href{https://www.iacr.org/archive/crypto2002/24420288/24420288.pdf}{Wagner, \emph{A Generalized Birthday Problem} (CRYPTO 2002)}. Generalized birthday algorithm and its $k$-list complexity analysis.
\item \href{https://eprint.iacr.org/2019/501}{Naya-Plasencia and Schrottenloher, \emph{Optimal Merging in Quantum $k$-XOR and $k$-SUM Algorithms} (EUROCRYPT 2020)}. Quantum merging-tree framework and optimization of single-solution instances.
\item \href{https://ir.cwi.nl/pub/31603/31603.pdf}{Schrottenloher, \emph{Improved Quantum Algorithms for the $k$-XOR Problem} (SAC 2021; proceedings 2022)}. Section 5: single-solution merging; Section 6, Theorem 4 and Algorithm 4: quantum-walk construction; Section 6.3: the $\widetilde O(2^{7n/24})$ four-list exponent.
\end{algorithmindex}


\entrytitle{88}{Even--Mansour recovery with quantum queries}{strengthPolynomialTime}{evidenceProved}{hardwarePolynomial}
\setcounter{problem}{87}
\begin{problem}[Even--Mansour key recovery with quantum chosen plaintexts]
Let $P$ be a public random permutation on $\{0,1\}^n$, and let $E(x)=P(x\oplus k_1)\oplus k_2$. Given coherent oracle access to both $E$ and $P$, recover $k_1,k_2$. The coherent chosen-plaintext access is part of this cryptographic model, commonly called Q2. The comparison counts coherent or classical oracle queries to the encryption and public permutation, with polynomial auxiliary quantum-gate and classical-bit costs.
\end{problem}

\speedupheading
\complexity{Quantum}{$O(n)$ coherent queries and polynomial auxiliary time, Kuwakado--Morii and Simon period finding.}
The function $F(x)=E(x)\oplus P(x)$ has period $k_1$. The bounded-extra-collision analysis for a random public permutation supplies the success guarantee for Simon sampling; polynomial-time linear algebra recovers the period. One additional evaluation gives $k_2$. Coherent evaluation of both $E$ and $P$ is available in the stated Q2 model, so each sample has polynomial additional gates and two oracle evaluations.

\complexity{Classical}{$\widetilde\Theta(2^n/D)$ public-permutation work with $D$ classical encryption queries, generic Even--Mansour attacks.}
The time--data tradeoff has scale $TD\asymp2^n$, with $2\le D\le 2^n$ and total work including data acquisition. Minimizing $D+T$ gives the square-root total-work scale. The lower bound concerns the public random permutation and secret-key encryption access specified by the idealized construction.

\dimensionlabel{Strength}{strengthPolynomialTime} Exponential to polynomial in $n$, from Simon period finding. This uses the XOR period created by the Even--Mansour construction and its coherent chosen-plaintext access.

\dimensionlabel{Evidence}{evidenceProved} Fully proved in the stated random-permutation oracle model, with the bounded-extra-collision analysis supporting the period-recovery algorithm.

\dimensionlabel{Hardware}{hardwarePolynomial} The algorithm needs polynomial logical qubits for the input, output and period equations, plus workspace for the public permutation and coherent encryption device. The depth is polynomial in $n$ times the oracle depth. Realizing the supplied coherent secret-key oracle is part of the hardware model.


\subhead{Algorithm}
\begin{algorithmindex}{88}
\item \href{https://arxiv.org/pdf/1602.05973}{Kaplan, Leurent, Leverrier and Naya-Plasencia, \emph{Breaking Symmetric Cryptosystems Using Quantum Period Finding} (CRYPTO 2016)}. Theorem 1: Simon sampling with an approximate promise; Section 3.2, ``Application to the Even-Mansour construction'': the Kuwakado--Morii attack and its extra-collision analysis; Figure 6: the period function.
\item \href{https://arxiv.org/pdf/2002.12439}{Bonnetain et al., \emph{Quantum Attacks Without Superposition Queries: The Offline Simon's Algorithm} (ASIACRYPT 2019; full version 2020)}. Section 5.1 and Table 2: classical Even--Mansour time--data tradeoff.
\end{algorithmindex}


\entrytitle{89}{Even--Mansour recovery with classical queries}{strengthPolynomial}{evidenceProved}{hardwareExponential}
\setcounter{problem}{88}
\begin{problem}[Even--Mansour key recovery with classical queries]
Recover the keys of $E(x)=P(x\oplus k_1)\oplus k_2$ when the attacker may query $E$ only on classical inputs and may compute the public permutation $P$ coherently. This is the Q1 version of entry 88, with a materially different access model and complexity. The Q1 comparison counts classical encryption queries, coherent public-permutation evaluations and quantum gates; explicit permutations also admit classical bit-complexity bounds.
\end{problem}

\speedupheading
\complexity{Quantum}{$\widetilde O(\sqrt{2^n/D})$ offline quantum time with $D$ classical encryption queries, offline Simon.}
For $D=2^u$ and $1\le u\le n$, Bonnetain et al. obtain the tradeoff $DT^2\asymp2^n$ up to polynomial factors in $n$. Total time includes data acquisition and has scale $\widetilde O(D+\sqrt{2^n/D})$. Reusable quantum states support the period tests with $O(n^2)$ logical qubits and polynomial classical memory. Only the public permutation is evaluated coherently; encryption queries are classical.

\complexity{Classical}{$\widetilde O(2^n/D)$ offline work with $D$ encryption queries, generic classical Even--Mansour attacks.}
The tradeoff $TD\asymp2^n$ gives total work $\widetilde O(D+2^n/D)$ and uses $O(D)$ classical storage. For a public permutation implemented by a $G_P$-gate circuit, enough known pairs give an explicit key-verification circuit of size $O(n+G_P)$, up to the fixed number of pairs. Tensor contraction yields $\operatorname{poly}(n+G_P)2^{O((n+G_P)/\sqrt{\log(n+G_P)})}$ classical key-recovery time. Families with $G_P=O(n)$ therefore have a subexponential classical comparator.

\dimensionlabel{Strength}{strengthPolynomial} Quadratic improvement of offline work at fixed data volume in the generic Q1 model, combining Grover search with offline Simon tests. Optimizing total time over $D$ gives $\widetilde O(2^{n/3})$ quantum versus $\widetilde O(2^{n/2})$ classical work, a polynomial speedup of degree $3/2$.

\dimensionlabel{Evidence}{evidenceProved} Rigorous algorithmic improvement over the matching generic classical time-data bound. Explicit permutations are additionally compared with their circuit-based classical algorithms.

\dimensionlabel{Hardware}{hardwareExponential} The construction uses $O(n^2)$ logical qubits and polynomial classical memory while processing $D$ classical encryption queries. The offline computation takes $\widetilde O(\sqrt{2^n/D})$ time; cryptographic-scale coherent depth requires fault tolerance.

{\small [Before ``ETH is false'', the explicit-permutation comparison used classical $\widetilde O(2^{n/2})$ balanced work versus $\widetilde O(2^{n/3})$ offline Simon time.]\par}
\subhead{Algorithm}
\begin{algorithmindex}{89}
\item \href{https://arxiv.org/pdf/2002.12439}{Bonnetain, Hosoyamada, Naya-Plasencia, Sasaki and Schrottenloher, \emph{Quantum Attacks Without Superposition Queries: The Offline Simon's Algorithm} (ASIACRYPT 2019; full version 2020)}. Section 3.2 and Algorithms 2--4: asymmetric shift search, period test and reusable database preparation; Section 5.1: Q1 Even--Mansour tradeoffs and the balanced $2^{n/3}$ attack.
\item \href{https://github.com/openai/math/blob/main/preprints/The-Exponential-Time-Hypothesis-is-False-October-25-2026/paper.pdf}{OpenAI, \emph{The Exponential Time Hypothesis Is False} (2026)}. Theorem 1.1 and witness self-reduction: search on the explicit key-verification circuit.
\end{algorithmindex}


\entrytitle{90}{Differential and linear cryptanalysis}{strengthPolynomial}{evidencePartial}{hardwareExponential}
\setcounter{problem}{89}
\begin{problem}[Differential and linear cryptanalysis]
Given a keyed permutation $E_K$, exploit either a differential $\Pr_x[E_K(x)\oplus E_K(x\oplus\Delta)=\nabla]=p$ or a linear approximation $\Pr_x[a\cdot x=b\cdot E_K(x)]=1/2+\varepsilon$ to distinguish it from random and recover key bits. The differential probability, linear bias, queried rounds, and guessed subkey size specify the instance. Complexity is measured in classical samples or oracle calls and coherent quantum evaluations for the statistical primitives, and in gate or bit complexity for complete key recovery.
\end{problem}

\speedupheading
\complexity{Quantum}{$O(p^{-1/2})$ differential tests and $O(|\varepsilon|^{-1})$ coherent bias-estimation evaluations, amplitude amplification and estimation.}
A differential right-pair search with success probability $p$ has the first bound. Estimating a linear bias to constant relative precision has the second. A search over $2^s$ subkeys uses $O(2^{s/2})$ verification steps, with the data-processing cost of each step included. Kaplan--Leurent--Leverrier--Naya-Plasencia analyze how these components combine for simple and truncated differentials and for linear cryptanalysis. The query model determines whether data acquisition itself can use coherent access.

\complexity{Classical}{$\Theta(p^{-1})$ differential trials, $\Theta(\varepsilon^{-2})$ bias-estimation samples and $O(2^s)$ exhaustive subkey checks.}
These are matched primitive costs. A complete classical attack minimizes over characteristics, data volume, guessed subkeys and other attacks on the same cipher and round count. If $k$ key bits are checked by a $G$-gate explicit verification circuit, Boolean tensor contraction adds $\operatorname{poly}(k+G)2^{O((k+G)/\sqrt{\log(k+G)})}$ to that classical minimum. Linear-size circuit families have subexponential complete key-recovery time.

\dimensionlabel{Strength}{strengthPolynomial} Quadratic improvements in the specified sampling and search components, from amplitude amplification or estimation. The exponent of a complete attack depends on the selected characteristic and key-recovery strategy.

\dimensionlabel{Evidence}{evidencePartial} Rigorous primitive bounds and partial cipher-specific evidence for a total-runtime separation. The full comparison uses the strongest complete classical attack with the same data access.

\dimensionlabel{Hardware}{hardwareExponential} Polynomial registers suffice for reversible cipher rounds and subkeys. Coherent-query attacks require the specified quantum access; Q1 attacks process stored classical data. Cryptanalytic-scale coherent depth requires fault tolerance, and memory depends on the data-processing strategy.

{\small [Before ``ETH is false'', the generic complete-key-search baseline was $O(2^k)$ classical versus $O(2^{k/2})$ amplified trials; differential and linear attacks optimized their own exponential components.]\par}
\subhead{Algorithm}
\begin{algorithmindex}{90}
\item \href{https://arxiv.org/pdf/1510.05836}{Kaplan, Leurent, Leverrier and Naya-Plasencia, \emph{Quantum Differential and Linear Cryptanalysis} (ToSC 2016)}. Sections 4.1--4.2: classical and quantum differential attacks; Sections 5.1--5.2: truncated differentials; Sections 7.1--7.2: linear cryptanalysis; Theorems 1--3: search and counting subroutines.
\item \href{https://arxiv.org/pdf/quant-ph/0005055}{Brassard, H\o yer, Mosca and Tapp, \emph{Quantum Amplitude Amplification and Estimation} (2002)}. Section 2: trial amplification; Theorem 12: amplitude estimation and its error bound.
\item \href{https://github.com/openai/math/blob/main/preprints/The-Exponential-Time-Hypothesis-is-False-October-25-2026/paper.pdf}{OpenAI, \emph{The Exponential Time Hypothesis Is False} (2026)}. Theorem 1.1 and witness self-reduction: search on the explicit key-verification circuit.
\end{algorithmindex}

\entrytitle{101}{Vector subset sum over \texorpdfstring{$\mathbb F_3$}{F3}}{strengthPolynomialTime}{evidencePartial}{hardwarePolynomial}
\setcounter{problem}{100}
\begin{problem}[Average-case $\mathbb F_3^n$-Subset-Sum]
Given independent uniform vectors $h_1,\ldots,h_m\in\mathbb F_3^n$, find a nonempty $S\subseteq\{1,\ldots,m\}$ such that $\sum_{i\in S}h_i=0$. The coefficients are restricted to $\{0,1\}$. This is a finite-field subset-sum problem related to restricted short integer solutions. Complexity is measured in quantum gates and classical bit operations on the explicit input.
\end{problem}

\speedupheading
\complexity{Quantum}{$(m/n)^{O_\kappa(\lceil n^2/m\rceil)}$ time, Kothari--Metger--O'Donnell--Shutty--Wu.}
For every fixed $\kappa>\log_2 3$ and $m\ge\kappa n$, success probability is at least $1-2^{-\Omega_\kappa(m)}$. The algorithm combines quantum filtering with a deterministic binary-error LWE decoder.

\complexity{Classical}{$\operatorname{poly}(n,m)2^{O(n/\log(m/n))}$ time in the Wagner regime; polynomial time above the classical vector-count threshold.}
For $m/n\to\infty$ below the known polynomial-time thresholds, Wagner-type algorithms merge lists of partial sums while imposing zero coordinates. For bounded $m/n$, known algorithms have $2^{O(n)}$ time. Kothari--O'Donnell--Wu give a deterministic polynomial-time algorithm, valid even for worst-case inputs, when $m\ge n^2/3+Cn\log n$ for an absolute constant $C$. Kothari--Metger--O'Donnell--Shutty--Wu further announce a polynomial-time threshold $m\ge c_\star n^2$, where $c_\star=(8\log_2 3-9)/36\approx0.1022$; Section 1.3 presents this as an announced construction with its formal proof deferred. The classical comparison takes the fastest applicable algorithm at the given $n,m$.

\dimensionlabel{Strength}{strengthPolynomialTime} Subexponential to polynomial in the strongest comparison. The quantum tradeoff reaches polynomial time throughout $m=\Omega(n^2)$, extending below the smallest reported classical polynomial-time vector-count threshold. The general bounds also describe intermediate improvements when both algorithms have superpolynomial time; their degree depends on the input density $m/n$.

\dimensionlabel{Evidence}{evidencePartial} Partial evidence from the classical frontier; the quantum running-time theorem is rigorous. The tradeoff is a September 2026 preprint.

\dimensionlabel{Hardware}{hardwarePolynomial} The Fourier construction has $m+n$ qutrit registers, together with the workspace for a reversible decoder. The gate count and decoder workspace follow the parameter-dependent running time. The polynomial-time regime has polynomial logical resources; larger decoding costs require correspondingly longer fault-tolerant computations.

\subhead{Algorithm}
\begin{algorithmindex}{101}
\item \href{https://arxiv.org/html/2609.40321v1}{Kothari, Metger, O'Donnell, Shutty and Wu, \emph{Exponential Quantum Speedup for $\mathbb F_3^n$-Subset-Sum? Or, Rigorous Classical Algorithms for Binary-Error LWE} (2026)}. Problem 1.1; Theorem 1.3 and Corollary 1.4: quantum tradeoff; Section 1.1: classical comparison; Theorem 3.1: quantum reduction; Definition 4.3 and Section 4.2: quotient decoder.
\item \href{https://arxiv.org/abs/2108.11015}{Chen, Liu and Zhandry, \emph{Quantum Algorithms for Variants of Average-Case Lattice Problems via Filtering} (EUROCRYPT 2022)}. Quantum filtering and the decoding-to-short-solution reduction underlying the quantum construction.
\item \href{https://arxiv.org/html/2510.07515v3}{Kothari, O'Donnell and Wu, \emph{No Exponential Quantum Speedup for SIS$^\infty$ Anymore} (STOC 2026)}. Theorem 1.2 and Section 2: deterministic polynomial-time subset sum above the classical vector-count threshold.
\item \href{https://people.eecs.berkeley.edu/~daw/papers/genbday.html}{Wagner, \emph{A Generalized Birthday Problem} (CRYPTO 2002)}. Section 2 of the short version: list merging and the generalized-birthday algorithm; the finite-field subset-sum specialization is compared in Kothari et al., Section 1.1.
\end{algorithmindex}

\clearpage
\section{Complexity bounds and implications}

The problems in this chapter are selected for their significance to quantum complexity classes. Query lower bounds, completeness results and classical simulation theorems clarify the power of $\mathsf{BQP}$ and the relationships among $\mathsf{NP}$, $\mathsf{MA}$ and $\mathsf{QMA}$. This scope includes continuous optimization and quantum-state problems, whose complexity also constrains the power of quantum computation.

\entrytitle{91}{Local minima on the Boolean hypercube}{strengthPolynomial}{evidenceProved}{hardwareExponential}
\setcounter{problem}{90}
\begin{problem}[Local search on the Boolean hypercube]
Given value-oracle access to $f:\{0,1\}^n\to\mathbb R$, output $x$ such that $f(x)\le f(x\oplus e_i)$ for every coordinate $i$, with probability at least $2/3$. Values have a finite comparison encoding. This models single-bit local improvement in combinatorial optimization. Complexity is measured in value-oracle queries for both quantum and randomized classical algorithms.
\end{problem}
\speedupheading
\complexity{Quantum}{$\Theta(2^{n/3}n^{1/6})$ queries, quantum sampling and local descent.}
Minimum finding selects a low-valued starting point, and quantum search finds improving neighbors. The bound is tight. On a fixed-dimensional grid with $N$ vertices and dimension at least six, the optimal quantum query complexity is $\Theta(N^{1/3})$.

\complexity{Classical}{$\Theta(2^{n/2}n^{1/2})$ queries, randomized sampling and local descent.}
The Aldous-type algorithm samples vertices, chooses the lowest sampled value, and follows improving neighbors. For a fixed-dimensional grid of dimension at least four, the tight randomized bound is $\Theta(N^{1/2})$.

\dimensionlabel{Strength}{strengthPolynomial} Polynomial speedup in the domain size $N=2^n$, with a smaller exponential constant in $n$. Minimum finding and amplitude amplification are algorithmic components.

\dimensionlabel{Evidence}{evidenceProved} Matching quantum and randomized query lower bounds establish the separation.

\dimensionlabel{Hardware}{hardwareExponential} A fault-tolerant implementation needs an $n$-qubit vertex register, value and comparison workspace, and coherent neighborhood access. Exponentially many queries dominate the runtime; the evaluator determines the additional logical-qubit count.

\subhead{Algorithm}
\begin{algorithmindex}{91}
\item \href{https://arxiv.org/pdf/quant-ph/0504085}{Zhang, \emph{(Almost) Tight Bounds for Randomized and Quantum Local Search on Hypercubes and Grids} (2005 preprint; journal version, 2009)}. Introduction, Theorems 1--2 and the discussion of matching sampling/descent upper bounds; Section 3 for the hypercube lower bound.
\item \href{https://arxiv.org/pdf/quant-ph/0307149}{Aaronson, \emph{Lower Bounds for Local Search by Quantum Arguments} (2003/2006)}. Proposition 2 for the quantum algorithm and the preceding paragraph for randomized sampling and descent.
\end{algorithmindex}

\entrytitle{92}{Total Boolean decision problems: a fourth-power bound}{strengthPolynomial}{evidenceProved}{white}
\setcounter{problem}{91}
\begin{problem}[Total Boolean function evaluation]
For a total function $f:\{0,1\}^N\to\{0,1\}$, determine $f(x)$ by accessing individual bits of $x$. Every $N$-bit string is a valid input. Threshold decision versions of explicitly encoded optimization problems are examples. Complexity is measured in input-bit oracle queries, with deterministic classical and bounded-error quantum algorithms distinguished.
\end{problem}
\speedupheading
\complexity{Quantum}{$Q(f)$ queries, an optimal bounded-error quantum decision algorithm.}
The theorem relates this resource to deterministic query complexity for every total Boolean function. In particular, every nonconstant monotone graph property on an $n$-vertex adjacency matrix has quantum query complexity $\Omega(n)$.

\complexity{Classical}{$D(f)=O(Q(f)^4)$ queries, deterministic decision-tree evaluation.}
A decision tree achieving this bound exists for every total $f$. Randomized query complexity satisfies $R(f)\le D(f)$. Known total functions attain the fourth-power deterministic-versus-quantum relationship up to logarithmic factors. The resource being bounded is the number of queried input bits; computation on the answers is free in this model.

\dimensionlabel{Strength}{strengthPolynomial} A polynomial upper bound on the possible query advantage, including separations greater than quadratic. The fourth-power relationship concerns deterministic classical complexity.

\dimensionlabel{Evidence}{evidenceProved} Unconditional query-complexity theorems, with the fourth-power exponent attained up to logarithmic factors.

\dimensionlabel{Hardware}{white} A complexity constraint. A concrete decision algorithm supplies its own address register, workspace and oracle circuit.

\subhead{Algorithm}
\begin{algorithmindex}{92}
\item \href{https://arxiv.org/pdf/2004.13231}{Aaronson, Ben-David, Kothari and Tal, \emph{Quantum Implications of Huang's Sensitivity Theorem} (2020)}. Theorem 1 for $D(f)=O(Q(f)^4)$; Theorem 3 for monotone graph properties. Sections 3--4 give the corresponding arguments.
\end{algorithmindex}

\entrytitle{93}{Symmetric promise problems and graph properties}{strengthPolynomialTime}{evidenceProved}{white}
\setcounter{problem}{92}
\begin{problem}[Symmetric decision problems and graph property testing]
Let $P\subseteq\Sigma^N$ and $f:P\to\{0,1\}$ be invariant under every permutation of the $N$ input positions. Given oracle access to $x\in P$, determine $f(x)$. For graph properties, require invariance under vertex relabeling and specify either adjacency-matrix or adjacency-list access. In the latter model, let $\mathcal P_k$ be the graph family of maximum degree five defined by Ben-David et al., Section 6, with tree-depth parameter $k\ge 1$. Given $0<\varepsilon<1$ and a neighbor oracle $[N]\times[5]\to[N]\cup\{\bot\}$ for an undirected multigraph $G=(V,E)$, $|V|=N$, distinguish
\[
G\in\mathcal P_k
\quad\text{from}\quad
|E\mathbin{\triangle}E'|\ge\varepsilon N
\ \text{for every }G'=(V,E')\in\mathcal P_k,
\]
with error at most $1/3$. Edge edits count multiplicity; self-loops and parallel edges are allowed. These models describe the role of symmetry and input access in quantum speedups. Complexity is measured in the corresponding input-symbol, adjacency-matrix, or neighbor-oracle queries.
\end{problem}
\speedupheading
\complexity{Quantum}{$(k/\varepsilon)^{O(1)}$ queries for $\mathcal P_k$, quantum-walk property testing; $Q(f)$ for the symmetric decision problems.}
The tester identifies weld vertices by quantum walks and checks the graph structure locally. It has perfect completeness and rejects $\varepsilon$-far inputs with probability at least $2/3$. Here $Q(f)$ and $R(f)$ denote bounded-error quantum and randomized query complexity; each comparison uses the same promise and oracle representation.

\complexity{Classical}{$\exp(\Omega(k))$ queries required for $\mathcal P_k$ at a fixed constant $\varepsilon>0$, and $O(N)$ sufficient by reading the graph; $O(Q(f)^3)$ for full permutation symmetry, Chailloux's simulation.}
The lower bound holds for every bounded-error randomized tester. Reading the adjacency lists permits exact membership testing with unrestricted offline computation. For adjacency-matrix graph properties, $R(f)=O(Q(f)^6)$; for fixed $r$-uniform hypergraphs, $R(f)=O(Q(f)^{3r})$.

\dimensionlabel{Strength}{strengthPolynomialTime} Exponential to polynomial query complexity in $k$ for adjacency-list testing of $\mathcal P_k$. Full permutation symmetry and adjacency-matrix graph symmetry impose the stated polynomial limits.

\dimensionlabel{Evidence}{evidenceProved} Unconditional classical query lower bounds establish the adjacency-list separation; classical simulation theorems establish the limits in the other models.

\dimensionlabel{Hardware}{white} Each vertex-label register uses $\lceil\log_2N\rceil$ qubits. Further resources depend on the neighbor oracle and quantum-walk implementation; the simulation theorems specify query complexity.

\subhead{Algorithm}
\begin{algorithmindex}{93}
\item \href{https://doi.org/10.4230/LIPIcs.ITCS.2019.19}{Chailloux, \emph{A Note on the Quantum Query Complexity of Permutation Symmetric Functions} (2019)}. Theorem 3; Section 3, especially Section 3.1, for the classical simulation.
\item \href{https://doi.org/10.1137/23M1573975}{Ben-David et al., \emph{Symmetries, Graph Properties, and Quantum Speedups} (2024)}. Theorem 1.2 and Corollary 4.14: matrix-model limits. Section 6, pp.~FOCS20-393--394: formal definition of $\mathcal P_k$ and edge-edit distance. Theorem 6.1 and Sections 6.1--6.3: quantum tester and classical lower bound.
\end{algorithmindex}

\entrytitle{94}{Exact counting and parity of feasible solutions}{white}{evidenceProved}{hardwareSmall}
\setcounter{problem}{93}
\begin{problem}[Exact counting and parity-SAT]
Given a predicate $f:\{0,1\}^n\to\{0,1\}$, compute $\sum_x f(x)$ exactly, or compute this integer modulo two. The predicate is supplied either as an oracle or as an explicit CNF formula. These tasks count feasible solutions and determine the parity of their number. Complexity is measured in predicate-oracle queries for the oracle formulation and in classical bit operations for explicitly encoded formulas.
\end{problem}
\speedupheading
\complexity{Quantum}{$\lceil N/2\rceil$ queries for parity and $\Theta(N)$ for exact counting, where $N=2^n$.}
The parity algorithm uses phase kickback on pairs of input bits. The query bounds apply to bounded-error computation of the exact answer. For explicit fixed-width CNF input, a quantum computer can also execute the classical counting algorithm below.

\complexity{Classical}{$N$ queries for oracle parity and $\Theta(N)$ for oracle counting; $2^{o(n)}\operatorname{poly}(n+m)$ time for explicit fixed-width CNF, tensor-contraction counting.}
Reading and accumulating the truth-table entries attains the oracle bounds. For a fixed-width $k$-CNF with $n$ variables and $m$ clauses, exact counting takes
\[
(n+m)^{O_k(1)}2^{O_k(n\log\log n/\sqrt{\log n})}.
\]
For $m\le Dn$ with fixed $D$, the time is $(n+m)^{O_{k,D}(1)}2^{O_{k,D}(n/\sqrt{\log n})}$. Reducing the count modulo two gives parity-$k$-SAT at the same cost.

\dimensionlabel{Strength}{white} Constant-factor oracle improvement for parity; equal asymptotic oracle complexity for exact counting. Explicit fixed-width counting has a subexponential classical upper bound without power saving.

\dimensionlabel{Evidence}{evidenceProved} Matching oracle lower bounds and rigorous counting upper bounds for the explicit representation.

\dimensionlabel{Hardware}{hardwareSmall} Small parity circuits are shallow implementations. General oracle counting requires exponentially many calls and $O(n)$ address/count registers in addition to evaluator workspace.

\subhead{Algorithm}
\begin{algorithmindex}{94}
\item \href{https://arxiv.org/pdf/quant-ph/9802049}{Beals et al., \emph{Quantum Lower Bounds by Polynomials} (1998/2001)}. Section 6, Proposition 6.4 and the preceding pairwise parity algorithm; Table 1 summarizes the tight bounds.
\item \href{https://github.com/openai/math/tree/main/preprints/The-Exponential-Time-Hypothesis-is-False-October-25-2026}{OpenAI, \emph{The Exponential Time Hypothesis is False} (2026)}. Theorem 1.1 for fixed-density counting and Appendix D for the general fixed-width consequence.
\end{algorithmindex}

\entrytitle{95}{NP search relative to an oracle}{strengthPolynomial}{evidenceProved}{hardwareExponential}
\setcounter{problem}{94}
\begin{problem}[Oracle-verifier search]
Given oracle access to a predicate $V(x,w)$ with $w\in\{0,1\}^n$, decide whether $\exists w\,V(x,w)=1$ and return such a witness when one exists, with probability at least $2/3$. This abstracts verifier access in SAT, exact cover and NP optimization decision problems. Complexity is measured in verifier-oracle queries; the explicit-circuit variant additionally uses classical bit-operation time.
\end{problem}
\speedupheading
\complexity{Quantum}{$\Theta(2^{n/2})$ verifier queries, amplitude amplification.}
This is the optimal worst-case bound for arbitrary verifier access. Bennett, Bernstein, Brassard and Vazirani prove that, relative to a random oracle with probability one, quantum time $o(2^{n/2})$ cannot contain NP in their oracle formulation.

\complexity{Classical}{$\Theta(2^n)$ verifier queries, exhaustive or randomized witness search.}
For an explicitly given bounded-fan-in verifier with $g$ gates, a count-preserving Tseitin encoding has $O(n+g)$ variables and clauses. Sparse fixed-width counting then gives witness search in
\[
\operatorname{poly}(n+g)2^{O((n+g)/\sqrt{\log(n+g)})}.
\]
A linear-size verifier therefore has a subexponential classical search algorithm. In the ordinary circuit model, a polynomial-time quantum algorithm for an NP-complete problem would imply $\mathrm{NP}\subseteq\mathrm{BQP}$.

\dimensionlabel{Strength}{strengthPolynomial} Tight square-root improvement in the number of candidate witnesses. In witness length $n$, the oracle comparison improves the exponential constant without power saving.

\dimensionlabel{Evidence}{evidenceProved} Fully rigorous for oracle-verifier access; the explicit-circuit upper bound uses the supplied circuit structure.

\dimensionlabel{Hardware}{hardwareExponential} The search kernel needs $n$ witness qubits plus a reversible verifier's ancillas. Exponentially many coherent iterations require fault-tolerant execution at large instance sizes.

\subhead{Algorithm}
\begin{algorithmindex}{95}
\item \href{https://arxiv.org/pdf/quant-ph/0005055}{Brassard et al., \emph{Quantum Amplitude Amplification and Estimation} (2002)}. Section 2, Theorems 2--3 and \textsc{QSearch}.
\item \href{https://arxiv.org/pdf/quant-ph/9701001}{Bennett et al., \emph{Strengths and Weaknesses of Quantum Computing} (1997)}. Section 3, Theorem 3.5, for the random-oracle separation.
\item \href{https://github.com/openai/math/tree/main/preprints/The-Exponential-Time-Hypothesis-is-False-October-25-2026}{OpenAI (2026)}. Theorem 1.1 and Appendix C for counting and circuit encodings.
\end{algorithmindex}

\entrytitle{96}{Fine-grained complexity of satisfiability and string problems}{white}{evidenceStrong}{white}
\setcounter{problem}{95}
\begin{problem}[Succinct truth-table properties]
Given a Boolean representation $C$ on $n$ inputs, compute a specified truth-table property, such as $\bigvee_x C(x)$, $\bigoplus_x C(x)$, or $\sum_x C(x)$. The representation class is fixed-width CNF, unrestricted-width CNF, a branching program, or a circuit. Fine-grained reductions connect these tasks to edit distance and longest common subsequence. Complexity is measured in running time on the explicit representation, using quantum gates and classical bit or word operations as specified by each reduction.
\end{problem}
\speedupheading
\complexity{Quantum}{$O^*(2^{n/2})$ for satisfiability by Grover search; $2^{o(n)}\operatorname{poly}(n+m)$ for fixed-width CNF by classical counting.}
Buhrman--Patro--Speelman's basic QSETH concerns unrestricted-width CNF and postulates a $2^{(1/2-o(1))n}\operatorname{poly}(m)$ lower bound. Their broader framework specifies both the representation class and the truth-table property. The edit-distance and longest-common-subsequence consequences give conditional $N^{3/2-o(1)}$ quantum lower bounds for length-$N$ strings under the corresponding compression-obliviousness assumptions.

\complexity{Classical}{$(n+m)^{O_k(1)}2^{O_k(n\log\log n/\sqrt{\log n})}$ for fixed-width CNF counting; $O(N^2)$ for the two string problems by dynamic programming.}
Sparse fixed-width counting improves to $\operatorname{poly}(n)2^{O(n/\sqrt{\log n})}$. The Satisfiability Hypothesis (SH) is equivalent to Weak ETH: for every fixed $\varepsilon>0$, 3-SAT has no $2^{O(n^{1-\varepsilon})}L^{O(1)}$ classical algorithm, where $L$ is input length. If a reduction maps 3-SAT size $s=\Theta(n+m)$ to target size $N\le s^{b+o(1)}$, SH excludes target time $2^{O(N^{1/b-\varepsilon})}$ for every fixed $\varepsilon>0$. Counting versions satisfy the analogous equivalence.

\dimensionlabel{Strength}{white} These are upper bounds and conditional lower-bound implications. The fixed-width counting upper bound is subexponential without power saving; SH addresses a fixed-power reduction inside the exponent.

\dimensionlabel{Evidence}{evidenceStrong} The reduction theorems are rigorous. The target lower bounds are conditional on the explicitly named SH or representation-specific QSETH assumptions.

\dimensionlabel{Hardware}{white} Complexity implications; implementation requirements are inherited from the algorithms in the relevant representation class.

[Before ``ETH is false'', the general classical 3-SAT bound was $O^*(1.30704^n)$, and amplified PPSZ claimed a quadratic improvement to $O^*(1.1433^n)$.]

\subhead{Algorithm}
\begin{algorithmindex}{96}
\item \href{https://doi.org/10.4230/LIPIcs.STACS.2021.19}{Buhrman, Patro and Speelman, \emph{A Framework of Quantum Strong Exponential-Time Hypotheses} (2021)}. Sections 1--2 for basic QSETH and the general framework; Corollaries 11--12 for consequences; Appendix B for example lower bounds.
\item \href{https://github.com/openai/math/tree/main/preprints/The-Exponential-Time-Hypothesis-is-False-October-25-2026}{OpenAI (2026)}. Theorem 1.1 and Appendix D for fixed-width counting.
\item \href{https://arxiv.org/abs/2611.02650}{MIT Hardness Group et al., \emph{ETH is Dead, Long Live SH} (2026)}. Theorem 2.4 and Corollary 2.5 for decision/counting equivalences; Theorem 3.1 for transfer through reductions.
\end{algorithmindex}

\entrytitle{97}{Stoquastic frustration-free constraint satisfaction}{white}{evidenceProved}{hardwarePolynomial}
\setcounter{problem}{96}
\begin{problem}[Stoquastic $k$-SAT]
Given $H=\sum_{a=1}^m H_a$ on $n$ qubits, with each positive semidefinite $H_a$ acting on at most $k=O(1)$ qubits and having real nonpositive off-diagonal entries, distinguish $\lambda_{\min}(H)=0$ from $\lambda_{\min}(H)\ge1/\operatorname{poly}(n)$. Zero energy means that the local constraints share a zero vector. Complexity is measured in classical running time and quantum evolution time or simulation gates, with the path and spectral-gap promises specified below.
\end{problem}
\speedupheading
\complexity{Quantum}{Polynomial evolution time for a suitably gapped, initialized path, stoquastic frustration-free adiabatic evolution.}
The polynomial guarantee uses the smoothness, inverse-polynomial gap and initialization assumptions (A0)--(A3) of Bravyi--Terhal. For unrestricted instances, stoquastic 6-SAT is MA-complete: a polynomial-time quantum solution would imply $\mathrm{MA}\subseteq\mathrm{BQP}$ and hence $\mathrm{NP}\subseteq\mathrm{BQP}$.

\complexity{Classical}{$\operatorname{poly}(n,\delta^{-1})$ time, Bravyi--Terhal sampling under the same path assumptions.}
The simulator samples a distribution within total variation distance $\delta$ of the final ground-state measurement distribution. Its nonnegative-amplitude random walks track the path. Stoquastic $k$-SAT is in MA for each fixed $k$, with a classical randomized polynomial-time witness verifier.

\dimensionlabel{Strength}{white} Polynomial classical simulation in the specified adiabatic regime; MA-completeness characterizes the unrestricted decision problem.

\dimensionlabel{Evidence}{evidenceProved} Rigorous simulation, containment and completeness theorems.

\dimensionlabel{Hardware}{hardwarePolynomial} The ideal analog system uses $n$ qubits; a digital realization adds Hamiltonian-simulation ancillas. Polynomial evolution requires the stated gap and initialization conditions.

\subhead{Algorithm}
\begin{algorithmindex}{97}
\item \href{https://arxiv.org/pdf/0806.1746}{Bravyi and Terhal, \emph{Complexity of Stoquastic Frustration-Free Hamiltonians} (2008/2009)}. Definitions 1.6--1.8 for stoquastic $k$-SAT; assumptions (A0)--(A3) and Theorem 1.4 for classical sampling; Section 5 for the sampler; Theorem 1.9 and Section 6 for MA verification and completeness.
\end{algorithmindex}

\entrytitle{98}{Local Hamiltonian ground-energy optimization}{strengthPolynomial}{evidencePartial}{hardwareExponential}
\setcounter{problem}{97}
\begin{problem}[Local Hamiltonian]
Given $H=\sum_{j=1}^m H_j$ on $n$ qubits, with constant-local terms of polynomially bounded norm, distinguish $\lambda_{\min}(H)\le a$ from $\lambda_{\min}(H)\ge b$, promised $b-a\ge1/\operatorname{poly}(n)$. Diagonal terms encode classical constraint optimization; noncommuting terms permit quantum witnesses. Complexity is measured in quantum gates or block-encoding queries and classical arithmetic operations at the required precision.
\end{problem}
\speedupheading
\complexity{Quantum}{$O^*(2^{n/2})$ gates, energy filtering and amplitude amplification at inverse-polynomial precision.}
A purification of the maximally mixed state supplies ground-subspace overlap at least $2^{-n/2}$. In a guided-state regime, a preparable state with ground-state overlap amplitude $\gamma$, gap lower bound $\Delta$, and block encoding of $H/\alpha$ allows ground-state preparation in $\widetilde O(\alpha/(\gamma\Delta))$ block-encoding queries. The 2-local decision problem is QMA-complete, so a polynomial-time quantum solution would imply $\mathrm{QMA}=\mathrm{BQP}$.

\complexity{Classical}{$O^*(2^n)$ arithmetic operations, randomized sparse-matrix Lanczos.}
Each matrix-vector product costs $O(m2^n)$ operations; inverse-polynomial accuracy contributes polynomial overhead. State-vector memory is exponential. For the diagonal subfamily with $O(n)$ constant-arity terms and constant-bit integer weights, a linear-size threshold circuit and sparse counting instead give exact minimization in $\operatorname{poly}(n)2^{O(n/\sqrt{\log n})}$ time.

\dimensionlabel{Strength}{strengthPolynomial} Polynomial speedup through a smaller exponential constant for the generic noncommuting comparison. Guided-state preparation is polynomial when its overlap, gap and block encoding have polynomial costs.

\dimensionlabel{Evidence}{evidencePartial} Proved algorithmic upper bounds and QMA-completeness; the generic runtime advantage is a comparison with known classical algorithms.

\dimensionlabel{Hardware}{hardwareExponential} A fault-tolerant implementation uses $2n$ qubits for the mixed-state purification, plus filtering and arithmetic ancillas. The generic method has exponential depth; guided-state resources depend on $\gamma$, $\Delta$ and precision.

\subhead{Algorithm}
\begin{algorithmindex}{98}
\item \href{https://arxiv.org/pdf/2002.12508}{Lin and Tong, \emph{Near-optimal Ground State Preparation} (2020)}. Theorem 6 for preparation with a ground-energy bound; Algorithm 1 and Theorem 8 for ground-energy estimation and binary search.
\item \href{https://doi.org/10.1137/0613066}{Kuczy\'nski and Wo\'zniakowski, \emph{Estimating the Largest Eigenvalue by the Power and Lanczos Algorithms with a Random Start} (1992)}. The Lanczos analysis for randomized extremal-eigenvalue estimation; apply it to $-H$.
\item \href{https://arxiv.org/pdf/quant-ph/0406180}{Kempe, Kitaev and Regev, \emph{The Complexity of the Local Hamiltonian Problem} (2004/2006)}. Theorem 1 for 2-local QMA-completeness.
\item \href{https://github.com/openai/math/tree/main/preprints/The-Exponential-Time-Hypothesis-is-False-October-25-2026}{OpenAI (2026)}. Theorem 1.1 and Appendix C for the diagonal circuit/counting specialization.
\end{algorithmindex}

\entrytitle{99}{Global optimization of variational-circuit parameters}{white}{white}{hardwareSmall}
\setcounter{problem}{98}
\begin{problem}[Variational quantum circuit optimization]
Given a parameterized circuit $U(\theta)$, an observable $H$, a parameter domain $\Theta\subseteq\mathbb R^p$ on which the minimum is attained, and $\varepsilon>0$, find $\theta\in\Theta$ such that
\[
\langle0|U(\theta)^\dagger H U(\theta)|0\rangle
\le\min_{\phi\in\Theta}\langle0|U(\phi)^\dagger H U(\phi)|0\rangle+\varepsilon.
\]
This is the outer optimization in VQE and variational combinatorial algorithms, including QAOA. Complexity is measured in circuit samples or coherent evaluation-oracle queries for objective evaluation, together with the classical or quantum time for the outer optimization.
\end{problem}
\speedupheading
\complexity{Quantum}{$O(K\delta^{-2})$ circuit samples or $O(K\delta^{-1})$ coherent oracle uses for $K$ objective evaluations, sampling or amplitude estimation.}
The objective is normalized to a bounded observable and each evaluation has additive accuracy $\delta$ with constant success probability. Simultaneous high confidence adds logarithmic factors. The global optimizer determines $K$. Bittel--Kliesch establish NP-hard parameter-optimization families even when the expectations are efficiently classically evaluable.

\complexity{Classical}{$2^{O(n/\sqrt{\log n})}\operatorname{poly}(n)$ expected time per exact sample, Gheorghiu--Gupte counting-based simulation of linear-size circuits.}
This bound holds for $n$-qubit circuits with $O(n)$ gates from a fixed finite set of bounded-arity algebraic gates. Sampling within total variation distance $\eta$ takes worst-case time $2^{O(n/\sqrt{\log n})}\operatorname{poly}(n,\log(1/\eta))$. Estimating a sampled observable multiplies this cost by the sample count. Gate count is measured after compilation to the specified finite set. The specialized $\{H,Z,CZ,CCZ\}$ method, with $h=O(n)$ Hadamards and $d\le Cn^{11/10}$ diagonal gates, has exact worst-case time $\operatorname{poly}(n)2^{O(n\log(2+d/n)/\sqrt{\log n})}$. Classical grid search supplies a global upper bound exponential in the parameter dimension for a compact domain and a known Lipschitz bound.

\dimensionlabel{Strength}{white} Coherent evaluation has a quadratic precision advantage. Global optimization complexity depends on the circuit family, landscape and outer optimizer.

\dimensionlabel{Evidence}{white} The evaluation bounds, classical sampling bounds and NP-hardness results are rigorous. A general global-optimization speedup has no established evidence.

\dimensionlabel{Hardware}{hardwareSmall} Small noisy circuits implement the sampling loop. Coherent amplitude estimation needs controlled state preparation and its inverse, with fault-tolerant depth. Logical storage consists of the ansatz register and evaluation ancillas.

[Before ``ETH is false'', generic state-vector simulation of linear-size circuits took $O^*(2^n)$ time; variational speedup claims often used this exponential evaluation baseline.]

\subhead{Algorithm}
\begin{algorithmindex}{99}
\item \href{https://arxiv.org/pdf/2101.07267}{Bittel and Kliesch, \emph{Training Variational Quantum Algorithms Is NP-Hard} (2021)}. Problem 4 and Theorem 1 for VQA optimization; Problem 6 and Theorem 3 for QAOA.
\item \href{https://arxiv.org/pdf/quant-ph/0005055}{Brassard et al., \emph{Quantum Amplitude Amplification and Estimation} (2002)}. Section 4, \textsc{Est\_Amp} and Theorem 12.
\item \href{https://drive.google.com/file/d/1RzhP9ZMy1xkC22oqKmMHEOh7gG1cwkfA/view}{Gheorghiu and Gupte, \emph{Subexponential-Time Sampling of Quantum Circuits from Exact Counting} (2026 note)}. Theorems 1--2; Sections 3--5 for conditional sampling, finite algebraic gates and the Hadamard/diagonal specialization.
\end{algorithmindex}

\entrytitle{100}{Adiabatic optimization: spectral-gap costs and oracle bounds}{strengthPolynomial}{evidenceProved}{hardwareExponential}
\setcounter{problem}{99}
\begin{problem}[Adiabatic ground-state optimization]
Given a smooth Hamiltonian path $H(s)$, $0\le s\le1$, with known initial ground state and final diagonal Hamiltonian $H(1)|x\rangle=C(x)|x\rangle$, prepare a final ground state and measure a minimizer of $C$. Write $\Delta(s)$ for the ground-state spectral gap. Complexity is measured in continuous evolution time with a fixed oracle energy scale; digital realizations additionally count Hamiltonian-simulation gates, and the classical comparison counts cost-oracle queries.
\end{problem}
\speedupheading
\complexity{Quantum}{$\Theta(\sqrt N)$ evolution time for unstructured marked-item search, locally adiabatic search; gap-dependent time for a general path.}
For a nondegenerate ground state, a rigorous adiabatic error bound is proportional to
\[
\frac1T\left[
\frac{\|H'(0)\|}{\Delta(0)^2}+\frac{\|H'(1)\|}{\Delta(1)^2}
+\int_0^1\left(\frac{\|H''(s)\|}{\Delta(s)^2}
+\frac{\|H'(s)\|^2}{\Delta(s)^3}\right)ds\right].
\]
Derivative and gap bounds therefore determine a sufficient evolution time $T$. Roland--Cerf's locally adjusted schedule attains the square-root bound for a marked item among $N$ configurations, with constant success probability and fixed energy normalization.

\complexity{Classical}{$\Theta(N)$ cost queries, exhaustive minimum finding or unstructured search.}
The classical procedure evaluates the costs and retains the smallest value. Structured cost functions have the problem-specific classical bounds given in the corresponding catalogue entries. Continuous-time oracle lower bounds of $\Omega(\sqrt N)$ match quantum unstructured search and apply to scrambled-cost constructions in adiabatic optimization.

\dimensionlabel{Strength}{strengthPolynomial} Tight square-root oracle improvement for unstructured search; the general-path runtime depends on the gap and derivatives. The search schedule realizes the adiabatic form of Grover's improvement.

\dimensionlabel{Evidence}{evidenceProved} Rigorous adiabatic error bounds and unstructured oracle lower bounds establish the stated comparisons.

\dimensionlabel{Hardware}{hardwareExponential} The system register has $\lceil\log_2N\rceil$ qubits, with simulation ancillas in a digital implementation. Search takes exponential evolution time in the register size; coherent control and precision contribute to implementation cost.

\subhead{Algorithm}
\begin{algorithmindex}{100}
\item \href{https://arxiv.org/pdf/quant-ph/0107015}{Roland and Cerf, \emph{Quantum Search by Local Adiabatic Evolution} (2002)}. Equations (16)--(20), especially the schedule in (17)--(18).
\item \href{https://arxiv.org/pdf/quant-ph/0603175}{Jansen, Ruskai and Seiler, \emph{Bounds for the Adiabatic Approximation with Applications to Quantum Computation} (2007)}. Section II, Theorem 3, for the explicit derivative/gap error bound.
\item \href{https://arxiv.org/pdf/quant-ph/0512159}{Farhi et al., \emph{How to Make the Quantum Adiabatic Algorithm Fail} (2005/2008)}. Section III, Theorem 2, for the scrambled-cost construction and oracle runtime bound.
\end{algorithmindex}

\section*{Relations for the knowledge graph}
\addcontentsline{toc}{section}{Relations for the knowledge graph}
The following relations connect the problem entries through algorithms, assumptions, and computational models. Entry numbers are stable identifiers for subsequent expansion of the catalogue.

\begin{longtable}{@{}p{0.24\textwidth}p{0.23\textwidth}p{0.46\textwidth}@{}}
\toprule
Relation & Connected entries & Mathematical content \\
\midrule
\endhead
Amplitude amplification & \ref{entry:1}, \ref{entry:2}, \ref{entry:5}, \ref{entry:6}, \ref{entry:78}, \ref{entry:79} & Amplify a successful classical trial or estimate its success probability; the verifier, trial cost and success probability determine the comparison. \\
Quantum tree search & \ref{entry:3}, \ref{entry:4}, \ref{entry:7}, \ref{entry:74} & The same branching and pruning rules define the classical tree and the quantum walk. Tree size is an instance-dependent resource. \\
Subset dynamic programming & \ref{entry:7}--\ref{entry:19} & Classical preprocessing of small subsets is combined with quantum minimum finding over decompositions; exponential table memory appears in many implementations. \\
Quantum walks and span programs & \ref{entry:37}--\ref{entry:44}, \ref{entry:91} & Marked subsets, graph witnesses and eigenphase detection yield improvements beyond a direct search over candidate certificates. \\
Graph input representation & \ref{entry:31}--\ref{entry:44}, \ref{entry:46}, \ref{entry:59}, \ref{entry:93} & Adjacency matrices, adjacency arrays, explicit lists and property-testing promises lead to different classical and quantum complexities. \\
Fourier sampling & \ref{entry:61}--\ref{entry:68} & Periods, hidden subgroups, real periods and hidden shifts connect algebraic number theory with cryptanalysis. \\
Lattice algorithms as subroutines & \ref{entry:71}, \ref{entry:72}, \ref{entry:74}--\ref{entry:77} & SVP, Gaussian sampling and decoding costs feed into reduction, closest-vector, LWE and SIS algorithms. \\
Code decoding and list matching & \ref{entry:78}--\ref{entry:80} & Information-set search, representations and nearest-neighbor methods connect code-based cryptanalysis with subset sum. \\
Oracle lower bounds & \ref{entry:1}, \ref{entry:6}, \ref{entry:91}--\ref{entry:95} & The classical lower bound and the quantum upper bound refer to the same queried object and success requirement. \\
Conditional lower bounds & \ref{entry:8}, \ref{entry:96}, \ref{entry:98} & Set-cover hypotheses, SH/Weak ETH, representation-specific QSETH variants and Hamiltonian complexity attach consequences to hypothetical algorithmic improvements. \\
Exact classical counting & \ref{entry:5}, \ref{entry:10}, \ref{entry:20}, \ref{entry:49}--\ref{entry:54}, \ref{entry:94} & Fixed-width counting algorithms yield feasibility, witness recovery and optimization through appropriately sized encodings. \\
Simulation from counting & \ref{entry:49}, \ref{entry:56}, \ref{entry:99} & Feynman path counts and conditional probabilities yield subexponential classical sampling bounds for specified quantum circuit families. \\
Hamiltonian and variational optimization & \ref{entry:97}--\ref{entry:100} & Frustration-free structure, ground-state overlap, spectral gaps, and parameter optimization control the full cost. \\
Constraint-preserving variational circuits & \ref{entry:103}, \ref{entry:104}, \ref{entry:106} & One-hot, fixed-Hamming-weight and long/short encodings define feasible subspaces for alternating-operator mixers. \\
Variational optimization of discrete objectives & \ref{entry:49}, \ref{entry:105}, \ref{entry:107} & Cost estimation, parameter updates and measured classical solutions connect QAOA, VQE, filtering and local-search methods. \\
Quantum filtering and decoding & \ref{entry:77}, \ref{entry:101} & A classical decoder supplies the coherent uncomputation step in quantum algorithms for restricted short relations; the distribution and number of input vectors determine the comparison. \\
Recursive Fourier sampling & \ref{entry:102} & Growing recursion depth gives superpolynomial quantum and classical query costs; the classical lower bound establishes a superpolynomial separation. \\
\bottomrule
\end{longtable}



\end{document}
